% Options for packages loaded elsewhere
\PassOptionsToPackage{unicode}{hyperref}
\PassOptionsToPackage{hyphens}{url}
\PassOptionsToPackage{dvipsnames,svgnames,x11names}{xcolor}
%
\documentclass[
  11pt,
]{article}
\usepackage{amsmath,amssymb}
\usepackage{iftex}
\ifPDFTeX
  \usepackage[T1]{fontenc}
  \usepackage[utf8]{inputenc}
  \usepackage{textcomp} % provide euro and other symbols
\else % if luatex or xetex
  \usepackage{unicode-math} % this also loads fontspec
  \defaultfontfeatures{Scale=MatchLowercase}
  \defaultfontfeatures[\rmfamily]{Ligatures=TeX,Scale=1}
\fi
\usepackage{lmodern}
\ifPDFTeX\else
  % xetex/luatex font selection
  \setmainfont[]{Liberation Serif}
  \setsansfont[]{Liberation Sans}
  \setmonofont[]{DejaVu Sans Mono}
\fi
% Use upquote if available, for straight quotes in verbatim environments
\IfFileExists{upquote.sty}{\usepackage{upquote}}{}
\IfFileExists{microtype.sty}{% use microtype if available
  \usepackage[]{microtype}
  \UseMicrotypeSet[protrusion]{basicmath} % disable protrusion for tt fonts
}{}
\makeatletter
\@ifundefined{KOMAClassName}{% if non-KOMA class
  \IfFileExists{parskip.sty}{%
    \usepackage{parskip}
  }{% else
    \setlength{\parindent}{0pt}
    \setlength{\parskip}{6pt plus 2pt minus 1pt}}
}{% if KOMA class
  \KOMAoptions{parskip=half}}
\makeatother
\usepackage{xcolor}
\usepackage[margin=0.8in]{geometry}
\usepackage{longtable,booktabs,array}
\usepackage{calc} % for calculating minipage widths
% Correct order of tables after \paragraph or \subparagraph
\usepackage{etoolbox}
\makeatletter
\patchcmd\longtable{\par}{\if@noskipsec\mbox{}\fi\par}{}{}
\makeatother
% Allow footnotes in longtable head/foot
\IfFileExists{footnotehyper.sty}{\usepackage{footnotehyper}}{\usepackage{footnote}}
\makesavenoteenv{longtable}
\setlength{\emergencystretch}{3em} % prevent overfull lines
\providecommand{\tightlist}{%
  \setlength{\itemsep}{0pt}\setlength{\parskip}{0pt}}
\setcounter{secnumdepth}{-\maxdimen} % remove section numbering
\ifLuaTeX
\usepackage[bidi=basic]{babel}
\else
\usepackage[bidi=default]{babel}
\fi
\babelprovide[main,import]{american}
\ifPDFTeX
\else
\babelfont{rm}[]{Liberation Serif}
\fi
% get rid of language-specific shorthands (see #6817):
\let\LanguageShortHands\languageshorthands
\def\languageshorthands#1{}
\usepackage{adjustbox}
\emergencystretch=3em
\AtBeginEnvironment{longtable}{\small\sloppy}
\usepackage{stmaryrd}
% Printing shims for run-on conditional-bar tokens in the supplied Biology source.
\providecommand{\midL}{\mid L}
\providecommand{\midN}{\mid N}
\sloppy
\ifLuaTeX
  \usepackage{selnolig}  % disable illegal ligatures
\fi
\usepackage{bookmark}
\IfFileExists{xurl.sty}{\usepackage{xurl}}{} % add URL line breaks if available
\urlstyle{same}
\hypersetup{
  pdflang={en-US},
  colorlinks=true,
  linkcolor={blue},
  filecolor={Maroon},
  citecolor={Blue},
  urlcolor={blue},
  pdfcreator={LaTeX via pandoc}}

\author{}
\date{}

\begin{document}

\section{COHERENCE UNIVERSALISM: METAPHYSICS AND
EPISTEMOLOGY}\label{coherence-universalism-metaphysics-and-epistemology}

\subsection{Coherence Under Constraint, Explanatory Order, and Warrant
Across
Registers}\label{coherence-under-constraint-explanatory-order-and-warrant-across-registers}

\textbf{Gaura Kiśora Dās Rader}\\
\textbf{Heaven≡Earth Foundation}\\
\textbf{White Paper Series}

\textbf{Version 0.10}\\
\textbf{October 2026}

\textbf{Release status.} This manuscript is part of the complete Version
0.10 proposal. Numbered principles identify the framework's current
formulations; their claim status, external validation, and authorization
for consequential application remain separate questions.

\textbf{Freeze status.} This manuscript is synchronized to the frozen
Version 0.10 release registry, evidence package, and review gates.
Future changes require a new versioned release.

\begin{center}\rule{0.5\linewidth}{0.5pt}\end{center}

\section{Abstract}\label{abstract}

Coherence Universalism (CU) is an integrative framework organized around
a structural question: what must be related, constrained, preserved, and
made correctable for a physical pattern, living organization, candidate
perspective, person, institution, or artificial system to remain
intelligible and viable through change? Version 0.10 answers this
question without assuming one universal scalar, one universal dynamics,
or one derivation from physics to value. Its mature unity is
methodological and architectural: one analytic grammar, many domain
realizations. Domains are related by shared declaration requirements and
explicit Translation Licenses rather than by an assumed identity of
quantity, geometry, mechanism, or value.

Coherence is defined as integrated relational organization under
specified constraints. A substantive assessment has the dependency form

\[\adjustbox{max width=\linewidth}{$\displaystyle 
C_{\mathcal D}=C_{\mathcal D}(x;K,G,\pi,T),
$}\]

where the domain, state or trajectory, constraints, relational
organization, scale or access map, and horizon are declared. The
operational realization may be a scalar, profile, partial ordering,
viable region, threshold family, or trajectory property. Scalarization
is a decision procedure, not a discovery of natural commensurability; it
must disclose its weights, floors, affected scope, horizon, and purpose.

Version 0.10 refines the two-register architecture. Description and
commitment are distinguished principally by speech-act force, use, and
assessment standard rather than by assuming that they always carry
different semantic kinds. Descriptive contents answer to truth,
evidence, proof, and correction. Commitment acts answer to warrant,
integrity, consequence, agency, and repair. Established facts constrain
commitments; commitment warrant does not establish facts. Mixed
utterances must be decomposed, and disputed register attribution must be
reported rather than forced.

The formal layer begins only after a relevant fragment has stable,
addressable realization conditions. It distinguishes descriptive
truth-satisfaction \(\vDash_D\) from commitment realization
\(\Vdash_R\). KLM preferential consequence governs inference over those
already interpreted contents; it does not itself solve the Frege--Geach
problem or explain how meaning is constituted. Within a fixed evaluation
model, the viability base, affected-agent scope, horizon, protected
constraints, and ordering remain fixed while premises filter by
intersection. Under Boolean closure, definability, and smoothness,
skeptical consequence satisfies System P. Ranked models support Rational
Monotony when relevant minima are attained.

The practical taxonomy contains five verdicts: strong warrant, mere
admissibility, exclusion, infeasibility, and underdetermination.
Classical negation is an adopted modeling choice within the Boolean
fragment and may flatten practical distinctions between absence of
support and positive support for cessation. Transformative commitments
are represented through trajectories; changed floors, affected scope,
horizon, or ordering constitute declared model revision. Formal
determinacy is always relative to a model and supplies no protection
against captured options, excluded affected parties, or institutionally
controlled correction.

Version 0.10 also restores the Heaven≡Earth core in disciplined form.
The interpretive register adopts dual-aspect coherence monism as its
canonical minimum, meaning realism, a minimum teleology or grain of
reality, love as the relational interior of coherence, and the Immanent
Omega. Omega is understood primarily as a locally realizable form of
universal love---complete in kind but partial in extent---rather than as
a required final singularity or end of history. Its phenomenological
ingredients are treated as humanly accessible through ordinary love,
awe, contemplative practice, collective experience, and psychedelic
states; accessibility is not metaphysical proof, stabilization, or
ethical embodiment.

The commitment register is expanded accordingly. CU commits to truth,
disciplined registers, revision, sacred regard, reciprocal dependence,
the grain of coherence as vocation, love as the ground and aim of
conscious life, meaning as discovery and participation, freedom of
meaning-revision, redemptive possibility without retroactive
justification of evil, Heaven on Earth as a revisable social attractor,
and enactment of Omega in present relations and institutions.

The Coherence Ladder remains a revisable dependency graph and the
Mandala an interpretive systems map. Neither is a proof by exhaustion, a
universal chronology, or a derivation of ethics from cosmology. Version
0.10 preserves the framework's scientific discipline while recovering
the worldview and practice that motivated the project.

\subsection{Relation to the series}\label{relation-to-the-series}

This paper provides the methodological, logical, metaphysical, and
epistemological architecture shared by the Coherence Universalism v0.10
series. The domain papers own their scientific, clinical, social,
ethical, and engineering claims. \emph{Foundations} maintains the
synchronized principle registry. \emph{Denying Knowledge to Make Room
for Commitment} is the active philosophical companion and controls the
extended exposition, worked cases, and regressions of the
commitment-logic clarification when consistent with M\&E and
\emph{Foundations}. \emph{The Architecture of Non-Egoic Artificial Mind}
is the active AI/consciousness architecture companion, subordinate to
the Consciousness, Ethics, and Artificial Intelligence authorities and
adding no independent principle authority.

\subsection{Note on Version 0.10}\label{note-on-version-0.10}

The manuscript presents the mature framework directly. It distinguishes
descriptive truth from commitment warrant; strong warrant from mere
admissibility; fixed-model inference from model revision; technical
dependency from interpretive traversal; and scientific evidence from
metaphysical and spiritual commitment. It restores the original
sixty-four-stage Coherence Mandala under explicit claim discipline and
includes the complete v0.10 proposals on structured evolutionary
possibility, Open Myth, and mythic--constitutional alignment with
bounded statuses.

The Yi Jing version of the Mandala remains a distinctive contemplative
contribution of \emph{You Are the Universe Becoming Conscious of
Itself}. It is not part of the canonical identity or evidential support
of the sixty-four philosophical stages.

\section{Contents}\label{contents}

\begin{enumerate}
\def\labelenumi{\arabic{enumi}.}
\tightlist
\item
  The Integrative Project\\
\item
  Method: Explanatory Order and Claim Architecture\\
\item
  Coherence Under Constraint\\
\item
  Constraint, Accessibility, Scale, and Identity\\
\item
  Truth, Evidence, and Epistemic Integrity\\
\item
  Content, Force, and Two Registers\\
\item
  Fixed-Base Nonmonotonic Consequence\\
\item
  Commitment Verdicts and Practical Warrant\\
\item
  Dynamic and Probabilistic Extensions\\
\item
  The Two-Stage Decision Architecture\\
\item
  The Epistemology Ladder\\
\item
  Governance of Inquiry and Action\\
\item
  The Coherence Ladder and Mandala\\
\item
  Gödel, Tarski, Löb, and the Limits of Formalization\\
\item
  Spiritual and Metaphysical Interpretation\\
\item
  The Commitment Register: Truth, Love, Meaning, Redemption, and
  Vocation\\
\item
  Research Program, Failure Conditions, and Conclusion
\end{enumerate}

Appendix A. Formal Definitions and Results\\
Appendix B. Countermodels and Obstruction Results\\
Appendix C. Claim-Status and Translation Protocol\\
Appendix D. Domain Dependency Map\\
Appendix E. Glossary\\
Appendix F. Cross-Domain Governance Matrix\\
References

\begin{center}\rule{0.5\linewidth}{0.5pt}\end{center}

\section{Domain Analytic Frame and CAP
Record}\label{domain-analytic-frame-and-cap-record}

\textbf{Target:} Coherence Universalism Version 0.10 release\\
\textbf{Status:} Release declaration; scope-limited and revisable under
the declared review gates\\
\textbf{Use rule:} Unresolved fields are left explicit rather than
completed by analogy.

\subsection{Paper and model identity}\label{paper-and-model-identity}

\begin{longtable}[]{@{}
  >{\raggedright\arraybackslash}p{(\columnwidth - 4\tabcolsep) * \real{0.3333}}
  >{\raggedright\arraybackslash}p{(\columnwidth - 4\tabcolsep) * \real{0.3333}}
  >{\raggedright\arraybackslash}p{(\columnwidth - 4\tabcolsep) * \real{0.3333}}@{}}
\toprule\noalign{}
\begin{minipage}[b]{\linewidth}\raggedright
Field
\end{minipage} & \begin{minipage}[b]{\linewidth}\raggedright
Declaration
\end{minipage} & \begin{minipage}[b]{\linewidth}\raggedright
Status
\end{minipage} \\
\midrule\noalign{}
\endhead
\bottomrule\noalign{}
\endlastfoot
Paper & Coherence Universalism: Metaphysics and Epistemology &
Specified \\
Domain owner & Framework / Metaphysics and Epistemology & Specified by
the Version 0.10 release corpus \\
Primary model or analysis & The shared claim architecture, translation
discipline, and fixed-base warrant models & Paper-relative \\
Primary claim status & Definitions, methodological principles, imported
formal results, bounded interpretations, and commitments & Mixed
statuses remain explicit \\
Review gate & Philosophy, logic, formal methods, owning-domain review,
and affected-community review for consequential applications &
Proportionate to claim and consequential reach \\
Methodological authority & CU-FP2, CU-M2, CU-M11, CU-M12, CU-M13,
CU-M14, and CU-M15 & Controlling Version 0.10 methodological
framework \\
\end{longtable}

\subsection{Shared Analytic Frame}\label{shared-analytic-frame}

\begin{longtable}[]{@{}
  >{\raggedright\arraybackslash}p{(\columnwidth - 4\tabcolsep) * \real{0.3333}}
  >{\raggedright\arraybackslash}p{(\columnwidth - 4\tabcolsep) * \real{0.3333}}
  >{\raggedright\arraybackslash}p{(\columnwidth - 4\tabcolsep) * \real{0.3333}}@{}}
\toprule\noalign{}
\begin{minipage}[b]{\linewidth}\raggedright
\#
\end{minipage} & \begin{minipage}[b]{\linewidth}\raggedright
Slot and domain declaration
\end{minipage} & \begin{minipage}[b]{\linewidth}\raggedright
Status and evidence
\end{minipage} \\
\midrule\noalign{}
\endhead
\bottomrule\noalign{}
\endlastfoot
1 & \textbf{Unit and boundary.} The focal unit is a declared claim,
model, translation, evaluation context, or commitment practice; the
paper is not itself a physical system model. & Specified in Sections
1--3 and 6--12 \\
2 & \textbf{State or trajectory space.} Depending on the analysis, the
object is a conceptual dependency graph, a set of worlds, a set of
trajectories, a profile, or a versioned framework state. & Plural by
design \\
3 & \textbf{Environment.} Domain sciences, lived practices,
institutions, evidence systems, and affected communities supply
constraints not generated by the formalism. & Explicitly external to the
formal layer \\
4 & \textbf{Constraints and viable region.} Truth, claim status,
protected floors, affected-agent scope, fixed-base semantics, register
discipline, and revision conditions constrain use. & Specified;
substantive contents remain domain-owned \\
5 & \textbf{Actions, options, and transitions.} Inferential moves,
cross-domain translations, commitment adoption, option selection, and
declared model revision are distinguished. & Specified conceptually and
formally where applicable \\
6 & \textbf{Dynamics or update rule.} No universal dynamics is asserted.
Fixed-base preferential filtering, trajectory lifting, and versioned
revision are separate model families. & Model-specific \\
7 & \textbf{Observation and access.} Evidence, proof, testimony,
records, first-person report, access maps, and institutional observation
channels must be declared. & Specified as a governance requirement \\
8 & \textbf{Perturbation and intervention class.} Counterexamples, new
evidence, premise additions, changed floors, option-set changes, and
institutional challenges are distinct perturbations. & Specified at the
methodological level \\
9 & \textbf{Equilibria, basins, or regimes.} Preferential minima may
exist inside declared evaluation models; no framework-wide dynamical
attractor or basin is claimed. & Not applicable globally \\
10 & \textbf{Local--global status.} Primarily qualitative and
coordinate-wise; exact or inequality relations belong only to separately
licensed models. & Governed by proposed CU-M13 \\
11 & \textbf{Scalar status.} No universal scalar. Measurement scalars,
decision scalarizations, and risk functionals may be used only as typed,
local instruments. & Governed by revised CU-FP2 \\
12 & \textbf{Statistical-image and allocation status.} The framework
profile is not a statistical distribution or allocation. Fixed-base
world sets, policy-induced outcome distributions, and declared
probability models may qualify only as question-relative subobjects. &
No paper-wide simplex license \\
13 & \textbf{Translation status.} Cross-domain relations are classified
from analogy through empirically validated instantiation. & Governed by
revised CU-M2 \\
14 & \textbf{Open gaps.} Operational measures, validated target-domain
mappings, and evidence of explanatory gain over separate theories remain
incomplete. & Open research program \\
\end{longtable}

\subsection{Compact Coherence Analysis Protocol
record}\label{compact-coherence-analysis-protocol-record}

\begin{longtable}[]{@{}
  >{\raggedright\arraybackslash}p{(\columnwidth - 4\tabcolsep) * \real{0.3333}}
  >{\raggedright\arraybackslash}p{(\columnwidth - 4\tabcolsep) * \real{0.3333}}
  >{\raggedright\arraybackslash}p{(\columnwidth - 4\tabcolsep) * \real{0.3333}}@{}}
\toprule\noalign{}
\begin{minipage}[b]{\linewidth}\raggedright
CAP module
\end{minipage} & \begin{minipage}[b]{\linewidth}\raggedright
Declaration
\end{minipage} & \begin{minipage}[b]{\linewidth}\raggedright
Boundary
\end{minipage} \\
\midrule\noalign{}
\endhead
\bottomrule\noalign{}
\endlastfoot
System declaration & Name the claim, model, domain, boundary, scale,
horizon, identity criterion, and represented relations. & Required
before comparison \\
Observation declaration & Separate observed, inferred, imported, latent,
excluded, and institutionally controlled evidence. & Formal determinacy
does not repair observation failure \\
Evaluation declaration & Declare profiles, floors, option sets,
orderings, scalar type, uncertainty, and preserved incomparability. &
Screening and scoring remain distinct \\
Power and authorship declaration & Name who framed the question,
supplied categories, controlled options and records, and may challenge
or obtain remedy. & Publicity is not contestability \\
Revision declaration & Distinguish premise addition from changed facts,
floors, affected scope, horizon, semantics, or ordering. & Changed model
does not inherit warrant \\
\end{longtable}

\subsection{License summary}\label{license-summary}

\begin{longtable}[]{@{}
  >{\raggedright\arraybackslash}p{(\columnwidth - 4\tabcolsep) * \real{0.3333}}
  >{\raggedright\arraybackslash}p{(\columnwidth - 4\tabcolsep) * \real{0.3333}}
  >{\raggedright\arraybackslash}p{(\columnwidth - 4\tabcolsep) * \real{0.3333}}@{}}
\toprule\noalign{}
\begin{minipage}[b]{\linewidth}\raggedright
License
\end{minipage} & \begin{minipage}[b]{\linewidth}\raggedright
Declaration
\end{minipage} & \begin{minipage}[b]{\linewidth}\raggedright
Current status
\end{minipage} \\
\midrule\noalign{}
\endhead
\bottomrule\noalign{}
\endlastfoot
Scalar license & Only local measurement, decision, or risk instruments;
no natural total order is presumed. & Conditional \\
Local--global license & No paper-wide quantitative theorem; relation
status must be declared model by model. & Qualitative / coordinate-wise
by default \\
Statistical Image License & Possible for declared worlds, outcome
distributions, belief models, or policy distributions; never for the
framework as a whole. & Object-specific \\
Allocation License & Possible for real resource, burden, or probability
allocations; not for the framework profile. & Object-specific \\
Refinement and Coarse-Graining License & Permitted only where category
splitting, merging, or stochastic mapping has declared semantics and
support treatment. & Conditional \\
Fisher-Rao License & No paper-wide license. Classical statistical
subobjects may inherit Fisher-Rao geometry under CU-M14 assumptions. &
Object-specific only \\
Translation license & Levels 0--4 distinguish analogy, shared schema,
property preservation, exact structure, and empirical validation. &
Framework-canonical proposal \\
\end{longtable}

\subsection{Open gaps and review
disposition}\label{open-gaps-and-review-disposition}

\begin{longtable}[]{@{}
  >{\raggedright\arraybackslash}p{(\columnwidth - 4\tabcolsep) * \real{0.3333}}
  >{\raggedright\arraybackslash}p{(\columnwidth - 4\tabcolsep) * \real{0.3333}}
  >{\raggedright\arraybackslash}p{(\columnwidth - 4\tabcolsep) * \real{0.3333}}@{}}
\toprule\noalign{}
\begin{minipage}[b]{\linewidth}\raggedright
Gap
\end{minipage} & \begin{minipage}[b]{\linewidth}\raggedright
Consequence
\end{minipage} & \begin{minipage}[b]{\linewidth}\raggedright
Required work
\end{minipage} \\
\midrule\noalign{}
\endhead
\bottomrule\noalign{}
\endlastfoot
Target-domain validation & Prevents claims of completed universal
theory. & Independent operationalization and empirical comparison \\
Authorship and model capture & Transparent formal models can remain
institutionally captured. & Governance, participation, and independent
correction \\
Semantic boundary & Preferential consequence presupposes stable
realization conditions. & Further semantics and practice-sensitive
analysis \\
\end{longtable}

\begin{quote}
\textbf{Propagation rule.} This card makes the paper's existing objects
and limits visible. It does not strengthen any empirical, mathematical,
normative, metaphysical, or spiritual claim beyond the wording supported
in the manuscript.
\end{quote}

\begin{center}\rule{0.5\linewidth}{0.5pt}\end{center}

\section{1. The Integrative Project}\label{the-integrative-project}

\subsection{1.1 The problem specialization does not
solve}\label{the-problem-specialization-does-not-solve}

For much of human intellectual history, questions about nature, mind,
ethics, politics, and spiritual practice belonged to one inquiry. Vedic
discussions of ṛta and dharma, Daoist accounts of the Way, Buddhist
analyses of suffering and dependent origination, Greek investigations of
nature and the good, and medieval syntheses of theology and natural
philosophy all assumed that how reality is, how human beings know, and
how they should live cannot be permanently isolated from one another.

Modern specialization broke these syntheses for good reasons.
Mathematical physics, evolutionary biology, experimental psychology,
historical criticism, and formal logic advanced by refusing inherited
answers that could not survive evidence. Their autonomy created forms of
precision the older unities did not possess.

But specialization solved one problem by exposing another. Technical
knowledge expanded faster than our ability to relate facts, purposes,
constraints, and responsibilities. A society may understand how to build
a powerful system without understanding who should authorize it, whose
interests it serves, which harms it externalizes, or how its mistakes
can be corrected. A person may possess accurate information without an
adequate framework for grief, commitment, responsibility, or meaning. An
institution may optimize what it can measure while degrading the
conditions that make its objective worth pursuing.

The need for integration therefore remains. The question is not whether
to return to a premodern unity. It is whether a modern synthesis can
preserve the standards that specialization earned.

\subsection{1.2 Integration without
reduction}\label{integration-without-reduction}

Coherence Universalism begins with four refusals.

It refuses to reduce biology, consciousness, morality, or meaning to one
physical quantity.

It refuses to treat every recurrence of a similar pattern as proof that
the same mechanism operates everywhere.

It refuses to make spiritual significance depend on a scientific
conclusion that science has not established.

And it refuses the opposite fragmentation in which every domain becomes
self-contained and no one remains responsible for the relations among
them.

The unity sought here has four parts.

First, CU asks shared questions across domains:

\begin{itemize}
\tightlist
\item
  What is being preserved?
\item
  Under which constraints?
\item
  Across which scales and horizons?
\item
  Through which relations?
\item
  At whose cost?
\item
  By what means can error be detected and repaired?
\end{itemize}

Second, it maps dependencies. A living organization presupposes physical
dynamics, but metabolism is not derivable from mechanics alone. A
perspective may presuppose globally available information, but global
availability does not by itself prove experience. Public justification
presupposes agents capable of giving and receiving reasons, but those
conditions do not decide every moral dispute.

Third, it states translation rules. A concept may change meaning as it
crosses domains. Physical persistence becomes biologically significant
only when a system's own processes contribute to maintaining the
conditions of its continuation. Information becomes psychologically
meaningful when it enters perception, memory, affect, identity, and
action. Coordination becomes politically legitimate only when power,
standing, contestation, and remedy are included.

Fourth, it traces recurring failure forms without assuming one universal
cause. Local optimization, insulation from correction, rigidity,
uncontrolled expansion, hidden externalization, and responsibility
without authority appear in different domains through different
mechanisms. Their recurrence supports comparison. It does not erase
difference.

\subsubsection{1.2.1 The mature form of
unification}\label{the-mature-form-of-unification}

The earlier framework made unity immediately visible by repeatedly using
the same objects: a scalar coherence functional, gradients, flows,
landscapes, attractors, and local--global conflict. The reconstruction
showed that these objects cannot be assumed to possess the same
interpretation or mathematical status in physics, biology,
consciousness, psychology, society, ethics, and artificial intelligence.

The mature framework does not recover unity by giving every domain the
same answer. It requires every domain to answer the same classes of
question, disclose the same kinds of modeling choice, and earn every
transfer of formal structure through an explicit license.

Its central methodological claim is therefore:

\begin{quote}
\textbf{One analytic grammar, many domain realizations.}
\end{quote}

The grammar asks what the unit is, what can vary, which constraints and
relations matter, which observations are available, what can change,
which scales and horizons are represented, what is protected, what
remains incomparable, and what would force revision. A common question
does not entail a common mechanism. Repetition of the frame makes both
unity and difference visible.

This also clarifies CU's public architecture:

\begin{itemize}
\tightlist
\item
  \textbf{worldview:} the Mandala and the bounded interpretive account
  of reality, meaning, love, and Omega;
\item
  \textbf{method:} the Coherence Analysis Protocol, Shared Analytic
  Frame, Translation Licenses, typed scalars, local--global declaration,
  and model revision;
\item
  \textbf{vocation:} the commitment register through which truth, sacred
  regard, love, repair, Heaven on Earth, and present enactment orient
  action.
\end{itemize}

None of the three substitutes for the others. The worldview does not
prove the domain sciences; the method does not by itself supply a moral
end; and the vocation does not manufacture descriptive truth.

\subsection{1.3 Coherence under
constraint}\label{coherence-under-constraint}

CU defines coherence as:

\begin{quote}
\textbf{integrated relational organization under specified constraints.}
\end{quote}

The definition is deliberately incomplete until a domain supplies its
variables. A substantive assessment has the dependency form

\[\adjustbox{max width=\linewidth}{$\displaystyle 
C_{\mathcal D}=C_{\mathcal D}(x;K,G,\pi,T),
$}\]

where:

\begin{itemize}
\tightlist
\item
  \(\mathcal D\) identifies the domain;
\item
  \(x\) identifies a state, process, or trajectory;
\item
  \(K\) identifies the relevant constraints;
\item
  \(G\) identifies relational organization;
\item
  \(\pi\) identifies the scale, boundary, or access map;
\item
  \(T\) identifies the temporal horizon.
\end{itemize}

This is a schema, not a universal equation. In one domain the
appropriate object may be a scalar. In another it may be a profile whose
dimensions conflict. In another it may be a viability region, a partial
ordering, a family of thresholds, or a property of a trajectory.

The framework therefore distinguishes coherence from several nearby
ideas:

\begin{itemize}
\tightlist
\item
  coherence is not simply low entropy;
\item
  coherence is not truth;
\item
  coherence is not fitness;
\item
  coherence is not mental health;
\item
  coherence is not social conformity;
\item
  coherence is not moral goodness;
\item
  coherence is not alignment;
\item
  coherence is not consensus.
\end{itemize}

Each of these may relate to a domain-specific form of organization. None
is interchangeable with the framework concept.

\subsection{1.4 Explanatory order}\label{explanatory-order}

Integration requires attention to explanatory order. A concept is
explanatorily prior to another when the proposed use of the second
presupposes the first.

A rate presupposes an ordered parameter. A reduced state presupposes a
distinction between retained and discarded variables. A claim about
local interaction presupposes a neighborhood or factorization. A
judgment of error presupposes a standard against which mismatch can be
registered. Moral responsibility presupposes some relation among
knowledge, control, agency, and consequence.

Explanatory priority is weaker than metaphysical fundamentality. A model
may require a time parameter without proving that time is fundamental. A
subsystem description may require a factorization even if the
factorization is emergent. A practice of justification may require
responsiveness to reasons without proving that reasons are physical
forces.

CU distinguishes four kinds of necessity:

\begin{enumerate}
\def\labelenumi{\arabic{enumi}.}
\tightlist
\item
  \textbf{Conceptual necessity:} a condition is required to use a
  concept intelligibly.
\item
  \textbf{Formal necessity:} a conclusion follows from declared axioms
  and rules.
\item
  \textbf{Empirical or nomological necessity:} a supported theory
  requires a condition within a stated regime.
\item
  \textbf{Practical or normative necessity:} a condition is required for
  an aim, practice, or form of agency.
\end{enumerate}

Confusing these kinds of necessity creates many of the overreaches an
integrative framework is meant to prevent.

\subsection{1.5 Bridges must be visible}\label{bridges-must-be-visible}

A bridge claim crosses from one level, domain, or claim type to another.
Examples include:

\begin{itemize}
\tightlist
\item
  physical persistence to biological self-maintenance;
\item
  causal access to candidate perspective;
\item
  functional stake to phenomenal valence;
\item
  social coordination to legitimacy;
\item
  metaphysical interdependence to spiritual commitment.
\end{itemize}

A bridge may be plausible and important. It is not thereby a deduction.
A responsible bridge states:

\begin{itemize}
\tightlist
\item
  the source construct;
\item
  the target construct;
\item
  the relation proposed;
\item
  the additional assumptions;
\item
  the evidence available;
\item
  the conditions under which the bridge would fail.
\end{itemize}

This requirement does not weaken synthesis. It makes synthesis
criticizable.

\subsection{1.6 Description and
commitment}\label{description-and-commitment}

CU distinguishes two broad registers primarily by \textbf{use,
speech-act force, and standard of assessment}, not by assigning every
stable content to one exhaustive semantic species.

An utterance functions in the \textbf{description register} when it
purports to represent how things are. Descriptive uses answer to
evidence, proof, valid inference, calibration, and correction.
Scientific, mathematical, historical, and ordinary factual assertions
typically function this way.

An utterance functions in the \textbf{commitment register} when its
principal role is to undertake, orient, consecrate, trust, refuse,
promise, organize attention, or structure action and relationship. A
vow, prayer, existential affirmation, or metaphysical commitment can use
content that is also descriptively expressible, but its adoption is
assessed through warrant, integrity, viability, consequence, agency,
revision, and repair.

The permeability between the registers is asymmetric:

\begin{quote}
\textbf{Established facts constrain commitments; commitments do not
establish facts.}
\end{quote}

A commitment may be profound without becoming evidence. A factual
correction may require a commitment to change. A metaphysical
interpretation may integrate a life without acquiring the status of an
empirical result.

This distinction allows the framework to preserve spiritual seriousness
without insulating spirituality from criticism.

\subsection{1.7 Governance belongs inside
inquiry}\label{governance-belongs-inside-inquiry}

Inquiry is not governed only by valid inference. It is also shaped by
who defines the problem, which alternatives are considered, whose
evidence is recorded, which harms are measured, and who can challenge
the result.

CU therefore treats seven governance requirements as part of responsible
inquiry and action:

\begin{enumerate}
\def\labelenumi{\arabic{enumi}.}
\tightlist
\item
  identify affected parties, distributions, and protected conditions;
\item
  examine how the problem and option set were framed;
\item
  preserve practical correction and sufficiently independent challenge;
\item
  align responsibility with authority, knowledge, control, and causal
  reach;
\item
  include affected communities in proportion to governance reach;
\item
  distinguish patiency, agency, responsibility, and sovereignty;
\item
  preserve reversibility, containment, or retained option value under
  severe risk, unless a necessity-and-proportionality exception is
  justified.
\end{enumerate}

These are normative and methodological commitments. They are not
physical laws. Their presence in the framework makes explicit what
consequential reasoning often leaves hidden.

\subsection{1.8 Scope}\label{scope}

This paper develops metaphysics in a restrained sense: it examines the
categories and dependency relations through which CU interprets what
kinds of structure may exist. It develops epistemology in a practical
and formal sense: it asks how claims are classified, corrected,
warranted, and related across registers.

It does not derive the universe from pure logic. It does not prove that
coherence is a substance, field, cosmic purpose, or moral agent. It does
not infer an ought from a physical gradient. It does not make religious
or spiritual commitment immune to evidence.

Its governing methodological claim is:

\begin{quote}
\textbf{An integrative framework becomes credible only when it preserves
the differences among the kinds of claim it integrates.}
\end{quote}

\begin{center}\rule{0.5\linewidth}{0.5pt}\end{center}

\section{2. Method: Explanatory Order and Claim
Architecture}\label{method-explanatory-order-and-claim-architecture}

\subsection{2.1 Explanatory priority}\label{explanatory-priority}

\textbf{Definition.} A concept or structure \(A\) is explanatorily prior
to \(B\) when a proposed definition, model, or use of \(B\) presupposes
\(A\).

Examples are straightforward. A temporal derivative presupposes an
ordered parameter. Entropy production presupposes a trajectory along
which entropy changes. A reduced state presupposes a distinction between
retained and discarded variables. A claim about local interaction
presupposes a neighborhood, factorization, or accessibility structure.
An inference about error presupposes a standard against which a mismatch
can be registered.

Explanatory priority is weaker than metaphysical fundamentality. A clock
parameter may be needed to state a model even if time is not
ontologically fundamental. A factorization may be required for a
subsystem description even if the factorization is emergent. The point
is not to legislate what reality ``really is'' but to expose hidden
assumptions.

\subsection{2.2 Four kinds of necessity}\label{four-kinds-of-necessity}

CU distinguishes four kinds of necessity that are often conflated in
integrative argument.

\subsubsection{2.2.1 Conceptual or semantic
necessity}\label{conceptual-or-semantic-necessity}

A condition is conceptually necessary when denying it undermines the use
of the concept at issue. Distinguishability is necessary for comparison;
some temporal or ordered structure is necessary for rates; some
stability of reference is necessary to reidentify an object.

Such arguments concern intelligibility. They do not determine a unique
physical realization.

\subsubsection{2.2.2 Formal necessity}\label{formal-necessity}

A conclusion is formally necessary relative to axioms when it follows by
a valid proof. Formal necessity is conditional on the language, axioms,
definitions, and inference rules. It is exact within that structure and
silent about whether the structure describes nature.

\subsubsection{2.2.3 Nomological or empirical
necessity}\label{nomological-or-empirical-necessity}

A condition is nomologically necessary when the laws of an empirically
supported theory require it within a stated regime. This form is
answerable to evidence and may be revised when the theory's domain
changes.

\subsubsection{2.2.4 Practical or normative
necessity}\label{practical-or-normative-necessity}

A condition is practically necessary relative to an aim, practice, or
form of agency. Trust-dependent cooperation requires some reliability;
moral justification requires responsiveness to reasons; a
self-correcting institution requires channels through which error can
change policy.

Practical necessity can reveal objective structural constraints without
automatically producing an unconditional moral command. ``If a community
is to preserve correction, it must not punish every dissenter'' is a
structural conditional. The further commitment that correction should be
preserved is not supplied by physics.

\subsection{2.3 The bounded transcendental
method}\label{the-bounded-transcendental-method}

\textbf{Definition.} A transcendental argument asks what conditions must
be in place for a specified activity, object, or form of intelligibility
to be possible (compare Kant, 1781/1787).

CU retains this method but narrows its reach. A transcendental argument
can establish that a practice presupposes some condition; it cannot by
itself establish that the condition is a fundamental constituent of the
cosmos.

For example:

\begin{itemize}
\tightlist
\item
  description presupposes differentiable content and sufficiently stable
  reference;
\item
  inquiry presupposes some sensitivity to correction;
\item
  agency presupposes alternatives represented under constraints;
\item
  justification presupposes standards not wholly determined by the
  conclusion being defended.
\end{itemize}

These are meaningful results. They become overclaims when converted into
statements such as ``therefore reality is a coherence field'' without an
additional bridge.

\subsection{2.4 Minimal commitment}\label{minimal-commitment}

The framework should introduce no stronger claim than the next
explanatory step requires. If temporal ordering is enough, a
thermodynamic arrow should not be assumed. If a partial ordering of
coherence profiles is enough, a universal scalar should not be imposed.
If a metaphysical interpretation is compatible with the mathematics,
compatibility should not be written as entailment.

Minimal commitment is not timidity. It concentrates the framework's risk
where it can be evaluated.

\subsection{2.5 Explicit bridges and Translation
Licenses}\label{explicit-bridges-and-translation-licenses}

\textbf{Methodological principle.} Whenever an argument crosses from one
claim type, domain, level, model class, population, scale, horizon, or
discourse register to another, the bridge premise must be stated.

Typical bridges include:

\begin{itemize}
\tightlist
\item
  conceptual dependence to physical hypothesis;
\item
  mathematical state space to ontological interpretation;
\item
  statistical association to mechanism;
\item
  physical persistence to biological viability;
\item
  information access to candidate perspective;
\item
  functional stake to phenomenal valence;
\item
  social coordination to legitimacy;
\item
  and warranted commitment to descriptive truth.
\end{itemize}

A bridge may be plausible, fruitful, and indispensable. It is not
thereby a deduction. The proposed CU-M2 architecture classifies the
strongest transfer actually earned:

\begin{enumerate}
\def\labelenumi{\arabic{enumi}.}
\tightlist
\item
  \textbf{Level 0 --- Analogy.} Suggestive language or pattern is
  shared; no theorem or evidential status transfers.
\item
  \textbf{Level 1 --- Shared schema.} The same questions or tuple
  structure organize both domains; quantitative conclusions do not
  transfer.
\item
  \textbf{Level 2 --- Property-preserving transfer.} A bounded result
  applies because the target satisfies declared properties required by
  that result.
\item
  \textbf{Level 3 --- Exact structural transfer.} The target belongs to
  the theorem's actual mathematical class and satisfies its hypotheses.
\item
  \textbf{Level 4 --- Empirically validated instantiation.} The
  transferred structure demonstrates explanatory, predictive,
  discriminative, measurement, or intervention value in target-domain
  evidence.
\end{enumerate}

A lower license is not a failed higher license. Analogy and shared
schema can be valuable when their status remains visible. Exact formal
fit does not establish empirical realization, ontological truth, or
normative legitimacy.

\subsection{2.6 Translation without
reduction}\label{translation-without-reduction}

Cross-domain reasoning requires mappings that preserve enough structure
for the target claim to be meaningful. A translation from thermodynamic
constraints to biological self-maintenance must say what changes: energy
becomes metabolically available work; persistence becomes actively
regulated viability; boundaries become maintained rather than merely
given.

A responsible translation identifies:

\begin{enumerate}
\def\labelenumi{\arabic{enumi}.}
\tightlist
\item
  the source construct;
\item
  the target construct;
\item
  the declared Translation License;
\item
  the relation or invariant preserved;
\item
  the new constraints and assumptions introduced;
\item
  the information or distinctions lost;
\item
  the evidence available in the target domain;
\item
  the conditions under which the mapping fails or must be narrowed;
\item
  any change of claim status; and
\item
  the review competence required.
\end{enumerate}

The license attaches to a particular transfer, not to a word. The fact
that \emph{constraint}, \emph{information}, \emph{integration},
\emph{memory}, \emph{agency}, or \emph{repair} appears in several papers
does not give every use the same formal or empirical content.

Without these items, ``the same pattern appears everywhere'' remains
analogy. With them, integration becomes a criticizable research program
rather than a vocabulary of resemblance.

\subsection{2.7 Reflexivity}\label{reflexivity}

CU is itself a coherence-maintaining intellectual structure. It can
become locally stable by insulating itself from criticism, multiplying
principles, or redescribing every counterexample as confirmation.
Reflexivity therefore requires more than saying the framework is
revisable.

A revision must be capable of producing genuine loss:

\begin{itemize}
\tightlist
\item
  a principle may be narrowed;
\item
  a theorem may be withdrawn;
\item
  a metaphor may lose canonical status;
\item
  a companion paper may be superseded;
\item
  a spiritual interpretation may remain meaningful while ceasing to
  function as a scientific claim.
\end{itemize}

A framework that can only absorb criticism has ceased to be corrigible.

\subsection{2.8 Burden of proof}\label{burden-of-proof}

The burden increases with the strength and reach of the claim. The
following sequence is not absolute but is a useful discipline:

\[\adjustbox{max width=\linewidth}{$\displaystyle 
\text{definition}
<
\text{formal construction}
<
\text{conditional recovery}
<
\text{empirical generalization}
<
\text{cross-domain law}
<
\text{fundamental ontology}.
$}\]

A formal construction can be exact without supporting the rightward
claims. A cross-domain law requires independent operationalization and
evidence in each domain together with a justified translation among
them.

\subsection{2.9 The method in one rule}\label{the-method-in-one-rule}

\begin{quote}
State what has been assumed, what has been shown, what has been
imported, what remains interpretive, and what observation or argument
would force revision.
\end{quote}

This is the methodological center of the framework.

\subsection{2.10 The Methodological
Spine}\label{the-methodological-spine}

The method in one rule is operationalized through a visible middle layer
between the Mandala and the domain sciences.

\begin{enumerate}
\def\labelenumi{\arabic{enumi}.}
\tightlist
\item
  \textbf{The Shared Analytic Frame} asks the same declared questions of
  every domain while allowing different state objects, measures,
  dynamics, and evidential strengths.
\item
  \textbf{The Coherence Analysis Protocol} records system, observation,
  evaluation, power and authorship, and revision declarations before a
  result is treated as portable or authoritative.
\item
  \textbf{Typed scalar licenses} distinguish measurement scalars,
  decision scalarizations, and risk functionals; protected floors are
  screened before any permitted scoring.
\item
  \textbf{Local--global declaration} identifies the focal scale,
  containing system, partition, affected units, horizon, externalized
  costs, and strongest mathematical or qualitative relation licensed.
\item
  \textbf{The Allocation License Test} determines whether a numerical
  object is genuinely simplex- or allocation-valued before probability,
  decomposition, coarse-graining, or information geometry is imported.
\item
  \textbf{Translation Licenses} govern how mathematics, mechanisms,
  measures, and interpretations travel between domains.
\item
  \textbf{Model revision and correctability} distinguish learning within
  a fixed model from changes to categories, floors, affected scope,
  options, horizon, semantics, or authority.
\end{enumerate}

\subsubsection{2.10.1 Typed mathematical
universality}\label{typed-mathematical-universality}

Retiring one universal coherence scalar does not imply that no
mathematical structure can recur canonically across domains. A system
may admit one or more \textbf{question-relative statistical images}:
declared distributions of outcomes, frequencies, policies, beliefs,
occupancies, or allocations. When such a subobject satisfies CU-M14's
Statistical Image, Refinement, and Fisher-Rao licenses, classical
information geometry supplies a canonical metric up to scale under the
relevant invariance assumptions (Čencov, 1982; Campbell, 1986; Ay et
al., 2015; Bauer et al., 2016).

The system is not identical with its statistical image. The geometry
does not choose the partition, factorization, reference distribution,
divergence, objective, dynamics, value direction, or interpretation. Nor
does use of the same theorem create cross-domain numerical
commensurability. The mature universal claim is therefore typed:
\textbf{one canonical statistical geometry where licensed, local
coherence content everywhere.}

The methodological maxim is:

\begin{quote}
\textbf{Declare before you compare. Screen before you score. Name who
authored the frame before accepting the verdict.}
\end{quote}

Declaration makes assumptions visible. It does not by itself make a
model true, complete, legitimate, contestable, or safe. Those further
properties require evidence, standing, independent challenge,
correction, and repair.

\begin{center}\rule{0.5\linewidth}{0.5pt}\end{center}

\section{3. Coherence Under
Constraint}\label{coherence-under-constraint-1}

\subsection{3.1 Canonical definition}\label{canonical-definition}

\textbf{Definition --- Coherence is integrated relational organization
under specified constraints.}

Each term matters.

\begin{itemize}
\tightlist
\item
  \textbf{Relational organization:} coherence concerns how components,
  processes, representations, or commitments bear on one another.
\item
  \textbf{Integration:} the relations form a mutually constraining
  organization rather than an accidental collection.
\item
  \textbf{Specified constraints:} what counts as successful organization
  depends on the conditions, functions, or invariants relevant to the
  system.
\end{itemize}

The definition is structural. It does not by itself assign a number,
moral value, purpose, or preferred direction.

\subsection{3.2 The dependency schema}\label{the-dependency-schema}

A substantive coherence assessment must state at least:

\[\adjustbox{max width=\linewidth}{$\displaystyle 
C_{\mathcal D}=C_{\mathcal D}(x;K,G,\pi,T).
$}\]

Here:

\begin{itemize}
\tightlist
\item
  \(\mathcal D\) is the domain;
\item
  \(x\) is the state, trajectory, model, or organization evaluated;
\item
  \(K\) is the relevant constraint set;
\item
  \(G\) is the relational or coupling structure;
\item
  \(\pi\) is a scale, partition, observation map, or access structure;
  and
\item
  \(T\) is the temporal horizon.
\end{itemize}

Additional arguments may be needed: environment \(E\), protected agents
\(A\), function set \(F\), uncertainty model \(P\), or intervention
class \(\mathcal I\).

A notation such as \(C(x)\) is acceptable only when these dependencies
have been fixed by context.

\subsection{3.3 Why one scalar is not
assumed}\label{why-one-scalar-is-not-assumed}

A scalar functional is useful when different considerations can be
validly aggregated into one ordering. But aggregation may conceal
tradeoffs or embed values that should remain explicit. In many domains,
a coherence profile is more honest:

\[\adjustbox{max width=\linewidth}{$\displaystyle 
\mathbf C(x)
=
\bigl(
C_{\mathrm{int}},
C_{\mathrm{rob}},
C_{\mathrm{adeq}},
C_{\mathrm{adapt}},
C_{\mathrm{cross}},
C_{\mathrm{repair}}
\bigr).
$}\]

Possible components include:

\begin{itemize}
\tightlist
\item
  internal integration;
\item
  robustness to perturbation;
\item
  adequacy to external constraints;
\item
  adaptive range;
\item
  cross-scale compatibility;
\item
  recoverability after damage.
\end{itemize}

The profile may support a partial order. Two systems can be
incomparable: one more robust, another more adaptive. A lexicographic
rule may place viability floors before optimization. A Pareto frontier
may be more appropriate than a single maximum.

\textbf{Open question.} CU does not yet possess a generally justified
aggregation rule across these dimensions.

\subsection{3.4 States, trajectories, and
capacities}\label{states-trajectories-and-capacities}

Coherence can be evaluated at three related levels.

\subsubsection{State coherence}\label{state-coherence}

A state may exhibit integrated organization relative to current
constraints. Examples include synchronized activity, compatible
subsystem states, or a well-integrated belief network.

\subsubsection{Trajectory coherence}\label{trajectory-coherence}

A trajectory may preserve relevant organization through time, including
periods of local disruption. This is necessary for learning,
development, repair, and identity.

\subsubsection{Coherence capacity}\label{coherence-capacity}

A system may possess a range of responses that keeps future trajectories
viable under perturbation. Capacity concerns the geometry of available
continuations, not only the path already taken.

Trajectory viability is an important realization of coherence for
organisms, agents, and institutions with meaningful conditions of
continuation. It is not a complete definition for every physical use, so
CU treats it as one domain-specific realization of the broader
structural concept.

\subsection{3.5 Constraint families}\label{constraint-families}

Constraints may be:

\begin{itemize}
\tightlist
\item
  \textbf{kinematic:} defining the admissible state space;
\item
  \textbf{dynamical:} restricting transitions;
\item
  \textbf{boundary or environmental:} imposed by surroundings;
\item
  \textbf{symmetry and conservation constraints:} preserving invariants;
\item
  \textbf{scale-relative:} arising under a coarse description;
\item
  \textbf{functional:} specifying performance conditions;
\item
  \textbf{viability-bearing:} marking states incompatible with continued
  self-maintenance;
\item
  \textbf{normative:} protecting agency, rights, or justificatory
  conditions.
\end{itemize}

The same mathematical form can play different roles. A boundary imposed
externally on a physical subsystem is not yet a norm the subsystem
recognizes. A loss function used by an engineer is not necessarily an
intrinsic stake of the model.

\subsection{3.6 Integration and
adequacy}\label{integration-and-adequacy}

Internal integration and external adequacy are distinct.

A system can be internally consistent but badly coupled to reality. A
database can contain accurate facts without integrating them into a
usable model. A community can coordinate efficiently around a false
narrative. A spiritual commitment can provide meaning while suppressing
grief.

This motivates at least a two-dimensional epistemic profile:

\[\adjustbox{max width=\linewidth}{$\displaystyle 
\mathbf C_{\mathrm{ep}}=(A,I),
$}\]

where \(A\) is adequacy to what is being tracked and \(I\) is
integration among the relevant elements.

The product

\[\adjustbox{max width=\linewidth}{$\displaystyle 
C=A\,I
$}\]

may be used as an illustrative scalarization when \(A\) and \(I\) are
normalized, commensurable, and jointly sufficient for the application.
It is \textbf{not} a framework-level identity. Other aggregations,
thresholds, or partial orders may be more appropriate.

\subsection{3.7 Coherence and entropy}\label{coherence-and-entropy}

Coherence is not generally low entropy. \emph{Coherence Universalism:
Physics}, Version 0.10 proves a restricted paired result for block
averaging: within-block statistical differentiation decreases and
Shannon entropy increases under the same operation, with a shared
equality condition. The two changes are not equal, and neither quantity
alone is substantive coherence.

In living systems, organized self-maintenance may require entropy
production and export. In learning, an increase in explored
possibilities can support later integration. In institutions, excessive
uniformity can be brittle. Entropy, uncertainty, variation, and freedom
of movement can enable or undermine coherence depending on constraints
and horizon.

\subsection{3.8 Coherence and stability}\label{coherence-and-stability}

Stability is neither necessary nor sufficient in a simple form. A rigid
delusion can be stable. A developing organism can remain coherent
through continuous change. A metastable state can preserve function over
the relevant horizon. A stable physical ground state may be coherent
without being adaptive.

The appropriate question is: stable with respect to which perturbations,
preserving which relations, over what time?

\subsection{3.9 Coherence and moral
goodness}\label{coherence-and-moral-goodness}

Coherence does not automatically confer moral value. A criminal
organization may be internally coordinated. A tumor can preserve local
organization. An authoritarian institution can maintain procedural
regularity. Moral evaluation introduces protected agents, asymmetries of
power, externalized costs, justificatory integrity, and local--global
compatibility.

The relation is therefore:

\[\adjustbox{max width=\linewidth}{$\displaystyle 
\text{moral adequacy}
\Rightarrow
\text{some forms of coherence are required},
$}\]

but not

\[\adjustbox{max width=\linewidth}{$\displaystyle 
\text{coherence}
\Rightarrow
\text{moral goodness}.
$}\]

\subsection{3.10 The Shared Analytic
Frame}\label{the-shared-analytic-frame}

Plural measures do not prevent a shared analytic grammar. Every
substantive domain model should answer the following questions where
they are meaningful:

\begin{enumerate}
\def\labelenumi{\arabic{enumi}.}
\tightlist
\item
  \textbf{Unit and boundary:} What is the candidate system, process, or
  governed object?
\item
  \textbf{State or trajectory space:} What can vary?
\item
  \textbf{Environment:} What coupled conditions lie outside the focal
  boundary?
\item
  \textbf{Constraints and viable region:} What must remain possible,
  protected, or identity-preserving?
\item
  \textbf{Actions, options, and transitions:} What changes can occur,
  which alternatives are omitted, and through whose agency?
\item
  \textbf{Dynamics or update rule:} What drives change, if a dynamics is
  specified?
\item
  \textbf{Observation and access:} What can be measured, known,
  reported, inferred, or recovered?
\item
  \textbf{Perturbation and intervention class:} Which disturbances or
  interventions are in scope?
\item
  \textbf{Equilibria, basins, or regimes:} Which are actually
  demonstrated?
\item
  \textbf{Local--global status:} Is the relation an exact decomposition,
  a monotonicity or inequality result, a coordinate-wise profile
  relation, or a qualitative analysis?
\item
  \textbf{Scalar status:} Is there no scalar, a measurement scalar, a
  decision scalarization, a risk functional, or another declared type?
\item
  \textbf{Allocation status:} Which quantities, if any, possess a
  genuine allocation or probability-simplex structure?
\item
  \textbf{Translation status:} Which license supports each imported
  theorem, metric, mechanism, or interpretation?
\item
  \textbf{Open gaps:} Which slots remain unspecified, unmeasured,
  inapplicable, disputed, or unvalidated?
\end{enumerate}

These questions provide the invariant architecture of CU. Domain papers
earn their membership in the framework by answering them precisely, not
by attaching the word \emph{coherence} to familiar phenomena.

Honest incompleteness is part of the result. Ethics need not possess a
natural flow equation. Psychology need not possess one general phase
portrait. Social Dynamics may identify regimes without demonstrating
basins. A blank or contested slot can show exactly where the framework
has and has not earned structure.

\subsection{3.11 The central translation
problem}\label{the-central-translation-problem}

\textbf{Open question.} The primary theoretical task of CU is to
determine when domain-specific models are related by analogy, shared
schema, property-preserving transfer, exact structural transfer, or
empirically validated instantiation.

Success would require:

\begin{itemize}
\tightlist
\item
  operational measures in each target domain;
\item
  a declared mapping and license level;
\item
  demonstrated invariance or correspondence under specified maps;
\item
  nontrivial predictions, constraints, or explanatory compression;
\item
  known failure and narrowing conditions;
\item
  and evidence that the shared framework explains or enables more than
  the separate domain theories already do.
\end{itemize}

Until then, CU is a disciplined integrative research program rather than
a completed universal theory. The methodological achievement is not a
proof that all domains instantiate one object. It is a governed account
of which relations have been established, which are proposed, and what
would be required for them to strengthen.

\section{4. Constraint, Accessibility, Scale, and
Identity}\label{constraint-accessibility-scale-and-identity}

\subsection{4.1 Constraint is constitutive, not merely
restrictive}\label{constraint-is-constitutive-not-merely-restrictive}

In ordinary language, a constraint is often understood as something that
limits an otherwise free system. In CU, the term has a wider and more
precise use.

\textbf{Definition.} A constraint is any stable condition that restricts
the admissible states, transitions, interpretations, or continuations of
a system relative to a question being asked.

Some constraints are externally imposed. Others are generated and
maintained by the system. Some close possibilities; others create them.
The rules of a language exclude many sequences of marks, but thereby
make indefinitely many meaningful statements possible. A cell membrane
limits molecular exchange, but the regulated boundary makes metabolism
possible. A promise constrains action, but can create trust and
coordinated futures that would otherwise be unavailable.

This yields a central CU distinction:

\begin{itemize}
\tightlist
\item
  \textbf{restrictive constraint:} removes possibilities without adding
  an organization capable of using the restriction;
\item
  \textbf{generative constraint:} removes possibilities in a way that
  stabilizes new capacities, relations, or meanings.
\end{itemize}

The distinction is functional and scale-relative. The same boundary may
be generative for one system and destructive for another. A tightly
controlled institution may coordinate an emergency response while
suppressing the distributed correction needed over a longer horizon.

\subsection{4.2 Constraint classes}\label{constraint-classes}

CU distinguishes at least six classes.

\begin{enumerate}
\def\labelenumi{\arabic{enumi}.}
\tightlist
\item
  \textbf{Admissibility constraints} determine which states or
  transitions belong to the modeled space.
\item
  \textbf{Viability constraints} identify conditions whose violation
  ends or seriously degrades the continued organization of a
  self-maintaining system.
\item
  \textbf{Constitutive constraints} help make the system the kind of
  system it is; removing them changes its identity rather than merely
  its performance.
\item
  \textbf{Instrumental constraints} are imposed for a task or
  optimization process and can be changed without necessarily changing
  the system's identity.
\item
  \textbf{Protected floors} specify outcomes or rights that may not be
  traded away in an aggregate optimization.
\item
  \textbf{Epistemic constraints} determine what evidence, methods, or
  inferential standards a claim must answer to.
\end{enumerate}

These categories may overlap. A norm against fabricating evidence is
simultaneously constitutive of scientific inquiry, a protected floor for
epistemic integrity, and an institutional rule.

The category of a constraint must be argued rather than inferred from
its mathematical form. A loss term in an artificial system is an
instrumental objective unless the architecture gives the system an
endogenous, persistence-bearing stake in satisfying it. A biological
need is not merely a number supplied by an observer; it is coupled to
the organism's continued organization.

\subsection{4.3 Accessibility is
relational}\label{accessibility-is-relational}

\textbf{Definition.} Accessibility is the availability of a distinction,
state variable, correlation, reason, or action possibility to a
specified subsystem, observer, scale, or operation.

Accessibility is therefore not an intrinsic property of information
considered in isolation. It has the general form

\[\adjustbox{max width=\linewidth}{$\displaystyle 
\mathsf A(z\mid s,\pi,\mathcal O,T),
$}\]

where \(z\) is the structure at issue, \(s\) is the accessing system,
\(\pi\) is the observational or coarse-graining map, \(\mathcal O\) is
the permitted operation class, and \(T\) is the horizon.

A correlation may be present in a global state while inaccessible to a
local subsystem. A reason may be available in memory but inaccessible
under threat. An institutional record may exist while being practically
inaccessible to those affected by it. A language model may contain a
representation that is available to some downstream computations but not
to a report channel.

The framework must therefore distinguish:

\begin{itemize}
\tightlist
\item
  existence of structure;
\item
  recoverability in principle;
\item
  access under a specified operation;
\item
  access in time to guide action; and
\item
  social or political permission to use the information.
\end{itemize}

Conflating these forms produces errors across domains. The statement
that information is ``not lost'' globally does not imply that any
situated agent can recover it. The statement that a person ``knows''
something in one context does not imply that the knowledge is
action-guiding under another. The statement that a public process is
formally transparent does not imply that it is epistemically legible.

\subsection{4.4 Scale and partitions}\label{scale-and-partitions}

Every claim about parts and wholes presupposes a partition,
factorization, or scale of description. Let

\[\adjustbox{max width=\linewidth}{$\displaystyle 
\pi:X\rightarrow Y
$}\]

map a fine-grained state space into a reduced description. The map may
identify microstates, aggregate variables, select observables, or define
an institutional reporting level.

Three questions must then be kept separate:

\begin{enumerate}
\def\labelenumi{\arabic{enumi}.}
\tightlist
\item
  What distinctions exist in \(X\)?
\item
  Which distinctions survive under \(\pi\)?
\item
  Does the physical or social process itself implement \(\pi\), or is
  \(\pi\) only an observer's description?
\end{enumerate}

\emph{Coherence Universalism: Physics}, Version 0.10 demonstrates why
this matters. Projection to a coarse state space, lifting back by a
block-uniform maximum-entropy representation, and actual physical
evolution are different operations. The descriptive loss of a
distinction is not automatically its physical destruction.

The same discipline applies elsewhere. A psychological diagnostic
category compresses a person; a national statistic compresses
communities; a benchmark compresses model behavior. The reduction may be
useful while hiding the very structure required for explanation.

\subsection{4.5 Local and global
coherence}\label{local-and-global-coherence}

The phrase \emph{local coherence} refers to organization evaluated
within a restricted subsystem, population, scale, partition, horizon,
observation map, or objective. \emph{Global coherence} refers to
organization evaluated across the larger coupled system relevant to the
claim.

Neither is absolute. The global system for one analysis may be local
within another. The crucial question is whether apparent improvement at
one level is purchased by damage, risk, labor, entropy production, loss
of agency, or repair cost exported beyond the selected boundary.

Let \(L\) denote a subsystem and \(G\) a containing system. A local
improvement

\[\adjustbox{max width=\linewidth}{$\displaystyle 
\Delta C_L>0
$}\]

is not a global improvement unless effects on coupled systems and later
horizons are included. The evaluation may require a vector:

\[\adjustbox{max width=\linewidth}{$\displaystyle 
\Delta\mathbf C
=
(\Delta C_L,\Delta C_{G\setminus L},\Delta C_{T_1},\Delta C_{T_2},\ldots).
$}\]

A corporation can improve quarterly performance by transferring
ecological and psychological costs outward. A person can reduce
immediate anxiety by narrowing the evidence they permit themselves to
encounter. An artificial agent can optimize a stated metric while
degrading the institution or environment that gave the metric meaning.

The strongest relation must be declared rather than implied:

\begin{enumerate}
\def\labelenumi{\arabic{enumi}.}
\tightlist
\item
  \textbf{Exact decomposition} applies only when the target quantity
  belongs to a proved decomposable class under the declared partition
  and weighting.
\item
  \textbf{Monotonicity or inequality} applies only when the functional
  and map satisfy the theorem's conditions; interpretation as
  improvement, loss, masking, or harm requires a further bridge.
\item
  \textbf{Coordinate-wise profile relation} reports separately which
  dimensions possess a result and which tradeoffs remain unresolved.
\item
  \textbf{Qualitative or structural analysis} identifies aggregation
  masking, cost transfer, scale reversal, delayed burden, or repair
  asymmetry without claiming one numerical identity.
\end{enumerate}

\textbf{{[}Bridge principle{]} Local--global compatibility.} A claim of
improvement is incomplete until the relevant scales, partitions,
affected units, externalized costs, observation limits, and horizons
have been identified.

Mathematical aggregation and normative noncompensation are distinct. A
theorem may show how a selected quantity behaves under averaging. It
does not decide which beings possess standing or which losses may not be
traded away. Ethics and governance supply protected floors; they are not
derived from convexity, decomposition, or information geometry.

\subsection{4.6 Identity through
transformation}\label{identity-through-transformation}

CU approaches identity through continuity of identity-relevant
constraints, relations, capacities, and histories.

\textbf{Definition.} An identity criterion specifies which relations,
capacities, histories, or constraints must remain invariant---or remain
connected by an admissible continuity relation---for a system to count
as the same system through change.

There is no assumption that every kind of identity has the same
criterion. Physical object identity may rely on worldline continuity or
conserved structure. Organism identity may involve autopoietic and
developmental continuity. Personal identity may involve memory,
embodiment, self-modeling, social recognition, and narrative continuity
in different proportions. Institutional identity may survive complete
turnover of personnel while failing after a change in charter or
practice.

A general schema is

\[\adjustbox{max width=\linewidth}{$\displaystyle 
x_t \equiv_{K_I,H} x_{t'}
\quad\Longleftrightarrow\quad
\text{there exists an admissible continuity path preserving }K_I
\text{ over history }H.
$}\]

Here \(K_I\) is the identity-relevant constraint set and \(H\) records
the path. This is a schema, not a universal metric.

The synchronized identity architecture is organized by five
subquestions.

\subsubsection{CU-FP8a --- Identity-relevant constraint
set}\label{cu-fp8a-identity-relevant-constraint-set}

A persistence claim must name the relations, capacities, boundaries,
commitments, or continuities whose loss would defeat the identity
judgment. A label, database key, product name, or uninterrupted external
record is not sufficient by itself.

\subsubsection{CU-FP8b --- Continuity across component
change}\label{cu-fp8b-continuity-across-component-change}

Material or computational components may change while an identity
continues, but the continuity process must be specified. Migration,
repair, copying, branching, reconstruction, and replacement are not
interchangeable merely because they preserve information.

\subsubsection{CU-FP8c --- Historical
dependence}\label{cu-fp8c-historical-dependence}

Two systems with indistinguishable present descriptions can differ in
identity when their histories differ. An analysis must state whether
history is constitutive, causally relevant, compressed into the present
state, or irrelevant to the practical question.

\subsubsection{CU-FP8d --- Endogenous
maintenance}\label{cu-fp8d-endogenous-maintenance}

Some robust identities are partly constituted by the system's own
participation in maintaining the conditions of its continuation. This is
important for living and agentic systems, but it is not assumed
necessary for every object, institution, software process, or thin
perspective.

\subsubsection{CU-FP8e --- Transformation and termination
criteria}\label{cu-fp8e-transformation-and-termination-criteria}

Every identity account should state what counts as continuity,
transformation, branching, copying, suspension, replacement, and
termination. The criteria may be purpose-relative, but the purpose,
authority, and people affected by the classification must be visible.

\subsection{4.7 Identity is not persistence
alone}\label{identity-is-not-persistence-alone}

Mere duration does not establish identity. A pattern may be recreated
after interruption without being the same token. A legal institution may
persist in name while violating the purposes that constituted it. A
person may undergo radical transformation while retaining enough
embodied, historical, and relational continuity to remain the same
person.

Conversely, identity need not require material sameness. Living systems
replace matter continuously. Cultural practices persist through
re-enactment. Software can migrate substrates.

The governing question is therefore not ``what substance remains?'' but:

\begin{itemize}
\tightlist
\item
  which relations are constitutive;
\item
  which transformations are identity-preserving;
\item
  which history matters;
\item
  who or what has standing to set the criterion; and
\item
  what practical decision the criterion is serving.
\end{itemize}

\subsection{4.8 Path dependence and
memory}\label{path-dependence-and-memory}

A state variable may summarize a history without containing it.
Introducing velocity produces a richer Markovian state, but does not by
itself capture all path dependence. Hysteresis, irreversible learning,
trauma, institutional precedent, and biological development may require
latent variables, changing constraint sets, or explicit history
functionals.

A general path-dependent coherence assessment may be written

\[\adjustbox{max width=\linewidth}{$\displaystyle 
C_t=C(x_t;K_t,G_t,\pi_t,T,H_{0:t}),
$}\]

where \(H_{0:t}\) is the relevant history. The notation is a reminder
that past events can alter not only the current state but the space of
available futures.

\textbf{Open question.} CU requires domain-specific criteria for
distinguishing compressed state memory from irreducible historical
dependence.

\subsection{4.9 Self-modeling, self-maintenance, and
ego}\label{self-modeling-self-maintenance-and-ego}

The consciousness and AI program requires three distinctions.

\begin{itemize}
\tightlist
\item
  \textbf{Self-modeling:} the system represents or tracks aspects of its
  own current activity.
\item
  \textbf{Self-maintenance:} the system acts so that a persisting
  organization of itself continues.
\item
  \textbf{Egoic organization:} the system treats a bounded,
  diachronically persistent self-model as a privileged object whose
  preservation governs interpretation and action.
\end{itemize}

These can dissociate. A system may self-model episodically without
maintaining a persistent identity. An organism may self-maintain with
little reflective self-modeling. A person may possess continuity without
rigidly defending a reified self-conception.

The distinction blocks an important inference: showing that an
architecture constrains persistent self-maintenance does not show that
it lacks all perspective. It may instead constrain egoic persistence
while leaving thin, momentary, or distributed forms of self-modeling
possible.

\subsection{4.10 Identity and
responsibility}\label{identity-and-responsibility}

Responsibility requires more than causal contribution. In ordinary moral
practice it usually presupposes enough continuity to own prior
commitments, learn from consequences, answer demands for justification,
and participate in repair.

A system that has no persisting self across episodes may affect the
world without being a full bearer of responsibility. It may still be a
moral patient if experience is present, and its designers and operators
may bear responsibility for its effects.

This creates the four-way distinction later governed by CU-G6:

\[\adjustbox{max width=\linewidth}{$\displaystyle 
\begin{aligned}
\text{patiency} &\not\Rightarrow \text{agency}, &
\text{patiency} &\not\Rightarrow \text{responsibility}, &
\text{patiency} &\not\Rightarrow \text{sovereignty},\\
\text{agency} &\not\Rightarrow \text{patiency}, &
\text{agency} &\not\Rightarrow \text{responsibility}, &
\text{agency} &\not\Rightarrow \text{sovereignty},\\
\text{responsibility} &\not\Rightarrow \text{patiency}, &
\text{responsibility} &\not\Rightarrow \text{agency}, &
\text{responsibility} &\not\Rightarrow \text{sovereignty},\\
\text{sovereignty} &\not\Rightarrow \text{patiency}, &
\text{sovereignty} &\not\Rightarrow \text{agency}, &
\text{sovereignty} &\not\Rightarrow \text{responsibility}.
\end{aligned}
$}\]

The four statuses require separate evidence and authorization. No one
entails any of the others without further conditions.

\begin{center}\rule{0.5\linewidth}{0.5pt}\end{center}

\section{5. Truth, Evidence, and Epistemic
Integrity}\label{truth-evidence-and-epistemic-integrity}

\subsection{5.1 Truth is not a special case of
coherence}\label{truth-is-not-a-special-case-of-coherence}

Truth and coherence answer to different questions.\footnote{Earlier
  formulations sometimes treated truth as the maximally coherent limit
  of inquiry. Version 0.10 rejects that identity while retaining
  coherence as an important epistemic property.}

\textbf{Definition.} A description-register claim is true when the world
is as the claim represents it to be, under the claim's semantics and
scope conditions.

This definition leaves open philosophical disputes about correspondence,
deflationism, realism, and model-mediated representation. It fixes the
practical asymmetry required by inquiry: the integration of a belief
system does not make a false claim true.

Coherence remains epistemically important. It can expose contradiction,
connect evidence, guide model selection, and make a body of claims
usable. But coherence is a property of an organization of claims,
models, or practices; truth concerns whether a particular truth-apt
claim obtains.

Thus:

\[\adjustbox{max width=\linewidth}{$\displaystyle 
\text{truth}\not\equiv\text{coherence},
\qquad
\text{coherence}\not\Rightarrow\text{truth}.
$}\]

A sufficiently rich epistemology needs both.

\subsection{5.2 Evidence and constraint}\label{evidence-and-constraint}

\textbf{Definition.} Evidence is information that bears on the
comparative support for truth-apt claims under a specified evidential
model.

Evidence is not whatever increases subjective confidence. It must be
connected to the claim by a defensible likelihood, causal, measurement,
testimonial, or inferential relation. The relevant standard varies by
domain, but the relation must be open to scrutiny.

An evidential assessment should specify:

\begin{itemize}
\tightlist
\item
  the claim and alternatives;
\item
  the observation or source;
\item
  the assumed connection between source and claim;
\item
  the background model;
\item
  possible defeaters;
\item
  measurement uncertainty; and
\item
  the scope over which the inference is licensed.
\end{itemize}

This prevents ``coherence with the framework'' from becoming a
substitute for evidence.

\subsection{5.3 Accuracy, consistency, integration, and
calibration}\label{accuracy-consistency-integration-and-calibration}

Epistemic evaluation is multidimensional. At minimum, it can
distinguish:

\begin{itemize}
\tightlist
\item
  \textbf{accuracy:} agreement between claims and what is observed or
  otherwise established;
\item
  \textbf{logical consistency:} absence of contradiction within a
  specified formalization;
\item
  \textbf{integration:} mutual explanatory and inferential connection
  among claims;
\item
  \textbf{calibration:} alignment between confidence and reliability;
\item
  \textbf{scope discipline:} use of claims only where their assumptions
  hold;
\item
  \textbf{revisability:} capacity to change appropriately under
  defeating information;
\item
  \textbf{action adequacy:} capacity to guide successful engagement
  where action is relevant.
\end{itemize}

These dimensions can conflict. A simple model may be well calibrated but
less explanatory. A theory may integrate many observations while
containing speculative bridges. A person may act skillfully without
possessing an explicit account of how.

The framework should resist converting these dimensions into one score
before the tradeoffs are justified.

\subsection{5.4 Epistemic integrity}\label{epistemic-integrity}

\textbf{Definition.} Epistemic integrity is the maintained capacity of
an agent or institution to let relevant evidence, valid criticism, and
disclosed uncertainty alter its claims, confidence, models, and actions
without selectively disabling the standards it invokes.

Epistemic integrity does not require immediate certainty or emotional
detachment. It requires symmetry of standards and genuine exposure to
correction.

Its practical indicators include:

\begin{itemize}
\tightlist
\item
  stating the conditions under which a claim would be revised;
\item
  applying comparable standards to favored and disfavored conclusions;
\item
  preserving records of failed predictions and withdrawn claims;
\item
  distinguishing confidence from identity;
\item
  rewarding the discovery of error rather than only the defense of
  position;
\item
  making inferential bridges visible; and
\item
  protecting dissent channels from retaliation.
\end{itemize}

Epistemic and justificatory integrity are necessary conditions of
defensible inquiry, but they are not outcome guarantees. A fully candid
agent can be factually wrong. A transparent institution can pursue an
unjust objective. A procedurally consistent system can exclude affected
parties or offer only harmful options. Integrity keeps reasons open to
correction; it does not supply truth, moral standing, protected floors,
or legitimacy by itself.

\subsection{5.5 Error and revision}\label{error-and-revision}

An error is not simply a decrease in systemic coherence. Correcting a
deeply embedded false belief can initially destabilize a person or
institution. The disruption may be necessary for later integration.

CU therefore rejects a simple update rule of the form ``accept whatever
immediately increases coherence.'' Revision must be evaluated across
dimensions and horizons. A responsible update can involve:

\begin{enumerate}
\def\labelenumi{\arabic{enumi}.}
\tightlist
\item
  registering the mismatch;
\item
  preserving enough local stability to avoid collapse;
\item
  revising the claim or model;
\item
  reorganizing dependent commitments;
\item
  repairing consequences produced by the error; and
\item
  retaining a record so the failure is not recreated.
\end{enumerate}

This is one reason trajectory analysis matters. The coherent response to
evidence may pass through a period of disorganization.

\subsection{5.6 Rationalization and
self-deception}\label{rationalization-and-self-deception}

\textbf{Definition.} Rationalization is the production or selection of
reasons primarily because they protect a preferred conclusion, identity,
attachment, or policy, while being presented as though they
independently generated the conclusion.

\textbf{Definition.} Systematic self-deception is a stable organization
of attention, memory, interpretation, and justification that selectively
suppresses or neutralizes coherence-threatening reasons in order to
protect a dominant commitment.

Neither concept is established by disagreement or error alone. The
diagnosis requires evidence of asymmetric processing, selective
standards, avoidance, confabulation, or recurrent distortion linked to
identity-protective pressure.

Self-deception can produce high local integration. The threatened
identity and its supporting narrative may fit together extremely well.
What fails is adequacy, permeability, and cross-scale integration. The
excluded evidence continues to appear through behavior, affect,
symptoms, institutional anomalies, or consequences borne by others.

\subsection{5.7 The anti-insulation
principle}\label{the-anti-insulation-principle}

\textbf{{[}Epistemic principle{]}} No claim, model, commitment, or
institution may count resistance to every possible criticism as evidence
of its own adequacy.

A framework becomes insulated when every apparent defeater is
redescribed as confirmation: criticism proves persecution, failed
predictions prove hidden forces, emotional distress proves spiritual
purification, and empirical mismatch proves that the evidence is too
``reductionist'' to apply.

Insulation differs from reasonable protection against poor criticism.
Every inquiry excludes some irrelevant demands. The issue is whether
there remains a credible route by which contrary information can change
the position.

For each substantial claim, CU v0.10 asks:

\begin{quote}
What observation, argument, or failure of translation would make this
claim narrower, less probable, or no longer canonical?
\end{quote}

\subsection{5.8 Disagreement}\label{disagreement}

Disagreement can occur at several levels:

\begin{itemize}
\tightlist
\item
  different factual beliefs;
\item
  different evidence or source trust;
\item
  different models connecting evidence to claims;
\item
  different risk tolerances;
\item
  different protected floors;
\item
  different commitments or forms of life; and
\item
  different interpretations of the same established facts.
\end{itemize}

The proper response depends on the level. A factual dispute may be
settled by evidence. A model dispute may require comparative prediction.
A conflict of protected commitments may require negotiation and
institutional limits rather than an evidential proof. A mixed
disagreement must be decomposed before it can be resolved.

The two-register architecture developed next is designed to make that
decomposition possible.

\subsection{5.9 Testimony, tradition, and
expertise}\label{testimony-tradition-and-expertise}

Much of what any person knows is received from others. Testimony is
neither automatically inferior to direct observation nor automatically
authoritative. Its warrant depends on source competence, honesty,
independence, track record, access, incentives, and opportunities for
correction.

Traditions preserve long-duration experiments in living, attention,
ritual, embodiment, and social organization. They may contain genuine
practical knowledge while also carrying historical error, power,
exclusion, and cosmological claims unsupported in the descriptive
register.

Expertise deserves weight where the expert has relevant training and
where the field has functioning correction mechanisms. Deference must
remain scoped: expertise in a technical domain does not confer authority
over every moral or metaphysical question.

\subsection{5.10 Epistemic humility}\label{epistemic-humility}

Humility is not generalized uncertainty. It is accurate representation
of the limits of one's access, methods, and confidence.

A claim may be held strongly when evidence is strong. False balance can
be as incoherent as dogmatism. Humility requires saying both what is
known and what remains open.

The framework therefore rejects two symmetric failures:

\begin{itemize}
\tightlist
\item
  \textbf{closure without warrant:} treating interpretation as
  established fact;
\item
  \textbf{perpetual suspension:} refusing commitment or action despite
  sufficient evidence and practical necessity.
\end{itemize}

Epistemic maturity consists in proportioned confidence and revisable
commitment.

\begin{center}\rule{0.5\linewidth}{0.5pt}\end{center}

\section{6. Content, Force, and Two
Registers}\label{content-force-and-two-registers}

\subsection{6.1 Why one evaluative standard is
insufficient}\label{why-one-evaluative-standard-is-insufficient}

Human language does more than describe. It promises, commands, blesses,
mourns, confesses, imagines, vows, orients, and creates relationships.
Treating every meaningful utterance as an empirical hypothesis
misdescribes much of religious, moral, poetic, and practical life.

The opposite mistake is equally serious. If a commitment is removed from
descriptive evaluation and supplied no other standard, it can appear
answerable to nothing. This is the insulation problem.

CU responds by distinguishing two primary \textbf{uses or registers}
while refusing to make them two sealed semantic substances.

\subsection{6.2 Content is not identical with speech-act
force}\label{content-is-not-identical-with-speech-act-force}

A stable content can be used in more than one way. A speaker may assert
it as a description, undertake it as a vow, adopt it as a plan, invoke
it ritually, or use it inside a practical argument. The same content may
therefore participate in different speech acts and answer to different
standards.

The two-register distinction is primarily a distinction of:

\begin{itemize}
\tightlist
\item
  principal function;
\item
  speech-act force;
\item
  inferential and practical role;
\item
  and standard of assessment.
\end{itemize}

It does not require the claim that every commitment content is
non-propositional or that every descriptive content is lexically marked.
The emphasis on use and practice is compatible with, but not reducible
to, the later-Wittgensteinian tradition (Wittgenstein, 1953).

\subsection{6.3 The description
register}\label{the-description-register}

\textbf{{[}Definition{]}} An utterance functions descriptively when it
purports to represent how things are.

Descriptive uses answer to:

\begin{itemize}
\tightlist
\item
  truth and falsity;
\item
  evidence;
\item
  valid inference;
\item
  explanatory comparison;
\item
  prediction and calibration;
\item
  and proof relative to explicit assumptions.
\end{itemize}

A claim does not leave the description register because it concerns God,
consciousness, value, or ultimate reality. If it describes, it remains
answerable to descriptive standards.

\subsection{6.4 The commitment register}\label{the-commitment-register}

\textbf{{[}Definition{]}} An utterance functions committally when its
principal role is to undertake, orient, consecrate, trust, refuse,
promise, organize attention, or structure action and relationship.

Commitment uses answer to questions such as:

\begin{itemize}
\tightlist
\item
  Is the operative content compatible with established fact?
\item
  What does the commitment train the agent to notice or exclude?
\item
  Does it preserve agency and protected conditions?
\item
  Whose costs does it externalize?
\item
  What forms of character, relationship, and institution does it
  generate?
\item
  Does it remain open to correction, exit, and repair?
\end{itemize}

A commitment act can be unwarranted even where the descriptive truth of
its content does not settle whether it should be adopted.

\subsection{6.5 Mixed discourse}\label{mixed-discourse}

Many important utterances are mixed. Consider:

\begin{quote}
Human beings are expressions of one sacred whole; therefore each person
should be treated with reverence.
\end{quote}

The sentence may contain:

\begin{enumerate}
\def\labelenumi{\arabic{enumi}.}
\tightlist
\item
  a metaphysical assertion or interpretation;
\item
  a commitment to sacred regard;
\item
  a normative conclusion.
\end{enumerate}

The moral attractiveness of the conclusion does not prove the
metaphysical assertion. The commitment need not collapse merely because
the metaphysical picture remains uncertain. Version 0.10 therefore
decomposes mixed discourse rather than granting one component the
authority of another.

\subsection{6.6 How principal function is
attributed}\label{how-principal-function-is-attributed}

Register classification is interpretive and can be disputed. Relevant
evidence includes:

\begin{itemize}
\tightlist
\item
  what the speaker says they are doing;
\item
  what kinds of correction they accept;
\item
  how the utterance behaves in inference;
\item
  what actions, permissions, or prohibitions it generates;
\item
  whether it is asserted across contexts or used mainly in ritual;
\item
  and what institutional consequences depend on taking it literally.
\end{itemize}

Speaker intention matters but is not conclusive. Institutions can use
symbolic language to perform covert descriptive or coercive work.
Outside observers can also misclassify prayer, lament, or devotion by
ignoring first-person practice. Genuine uncertainty about function
should be reported rather than forced.

\subsection{6.7 Asymmetric permeability}\label{asymmetric-permeability}

\textbf{{[}Canonical principle{]}} Established descriptive facts
constrain commitments. Commitment warrant does not establish descriptive
truth.

A commitment is unavailable in an operative use that requires systematic
denial of established fact. The converse does not hold: a commitment may
be beautiful, healing, or morally fruitful without proving its
cosmology.

Commitments can guide inquiry by directing attention, generating
hypotheses, and sustaining costly investigation. Inquiry is not
value-free. But motivation and salience are not evidence.

\subsection{6.8 Addressability and the boundary of
formalization}\label{addressability-and-the-boundary-of-formalization}

Not every meaningful utterance admits a complete Boolean decomposition.
Prayer, ritual, poetry, myth, silence, first-person avowal, and
contemplative expression can possess meaning without being fully
representable in the formal fragment.

A commitment is \textbf{addressable} to the degree that it makes a
stable, criticizable difference to:

\begin{itemize}
\tightlist
\item
  realizable trajectories;
\item
  practical priorities;
\item
  actions licensed or forbidden;
\item
  reasons the agent must answer;
\item
  relationships undertaken;
\item
  and conditions of revision or release.
\end{itemize}

The correct boundary term is therefore \textbf{outside this
formalization}, not automatically meaningless or nonsense.

\subsection{6.9 Register inheritance}\label{register-inheritance}

A mixed utterance inherits the stricter standard whenever it performs
descriptive work inside a practical argument.

\begin{quote}
This treatment is divinely required because it always cures the illness.
\end{quote}

The cure claim is descriptive and answerable to evidence. Its sacred
setting cannot insulate it. If the factual premise fails, the practical
conclusion loses that support even if the ritual or symbolic
interpretation remains.

\subsection{6.10 Moral, poetic, and mythic
discourse}\label{moral-poetic-and-mythic-discourse}

Moral discourse cannot be assigned wholesale to one register. It
contains descriptions of harm, normative judgments, public commitments,
authorizations, and reasons governing relationship.

Poetic and mythic language can disclose patterns of experience and
organize shared orientation without functioning as literal cosmology.
The Open Myth architecture adds a further requirement: mythic
commitments must remain truth-constrained, addressable where possible,
contestable, revisable, and open to refusal.

\subsection{6.11 The reflexive status of the
architecture}\label{the-reflexive-status-of-the-architecture}

The claim that discourse has these principal functions is itself a
philosophical description answerable to analysis and comparison.
Adopting the distinction as a discipline is a defeasible commitment.

Its root undertaking is:

\begin{quote}
\textbf{Make reality more important than the identity built around one's
interpretation of it.}
\end{quote}

CU must revise the architecture if another account better preserves
truth, meaning, criticism, spiritual practice, and responsible
commitment with fewer distortions.

\section{7. Fixed-Base Nonmonotonic
Consequence}\label{fixed-base-nonmonotonic-consequence}

\subsection{7.1 What the formal layer
assumes}\label{what-the-formal-layer-assumes}

The formal layer does not create the meanings on which it operates. It
begins only after a relevant fragment has been decomposed into stable,
representable contents.

Let:

\begin{itemize}
\tightlist
\item
  \(L_D\) be a descriptive language with truth-satisfaction
  \(w\vDash_D\delta\);
\item
  \(L_R\) be a declared commitment-representation language;
\item
  \(w\Vdash_R\varphi\) mean that world or trajectory \(w\) realizes or
  supports commitment content \(\varphi\).
\end{itemize}

For the fragment to which the full KLM result applies, \(L_R\) is
Boolean-closed and the realization extensions satisfy:

\[\adjustbox{max width=\linewidth}{$\displaystyle 
\llbracket\neg\varphi\rrbracket_R=W\setminus\llbracket\varphi\rrbracket_R,
$}\]

\[\adjustbox{max width=\linewidth}{$\displaystyle 
\llbracket\varphi\land\psi\rrbracket_R=
\llbracket\varphi\rrbracket_R\cap\llbracket\psi\rrbracket_R,
$}\]

\[\adjustbox{max width=\linewidth}{$\displaystyle 
\llbracket\varphi\lor\psi\rrbracket_R=
\llbracket\varphi\rrbracket_R\cup\llbracket\psi\rrbracket_R.
$}\]

This is an adopted representation of addressable commitment content. It
is extension-like and proposition-like enough for embedded inference. CU
does not claim semantic novelty here.

\subsection{7.2 Why defeasible consequence is
needed}\label{why-defeasible-consequence-is-needed}

Embedded agents reason under incomplete information. A conclusion may be
justified under current premises and withdrawn when defeating
information arrives:

\[\adjustbox{max width=\linewidth}{$\displaystyle 
\Gamma\mid\!\sim_C\varphi
\quad\text{while}\quad
\Gamma\cup\{\psi\}\not\mid\!\sim_C\varphi.
$}\]

Preferential semantics makes this defeasibility explicit by evaluating
what holds across the best premise-compatible realizations.

\subsection{7.3 Evaluation models}\label{evaluation-models}

Let

\[\adjustbox{max width=\linewidth}{$\displaystyle 
M=(W,V_F,\prec,\Vdash_R)
$}\]

be a preferential model for a declared evaluation context.

\begin{itemize}
\tightlist
\item
  \(W\) is a set of worlds, states, or trajectories.
\item
  \(V_F\equiv V_{F,K,A,H}\subseteq W\) is the fixed base compatible with
  descriptive facts \(F\), protected constraints \(K\), affected-agent
  scope \(A\), and horizon \(H\).
\item
  \(\prec\) is a declared preference relation fixed inside the model.
\item
  \(\Vdash_R\) supplies the realization extensions for the addressable
  commitment fragment.
\end{itemize}

For a finite theory \(\Gamma\subseteq L_R\), define:

\[\adjustbox{max width=\linewidth}{$\displaystyle 
W_{F,\Gamma}=V_F\cap\operatorname{Real}_R(\Gamma),
$}\]

and:

\[\adjustbox{max width=\linewidth}{$\displaystyle 
B(F,\Gamma)=\operatorname{Min}_{\prec}(W_{F,\Gamma}).
$}\]

Premises filter a fixed base. They do not silently change floors,
affected scope, horizon, viability criterion, or ordering.

\subsection{7.4 Skeptical preferential
consequence}\label{skeptical-preferential-consequence}

Define:

\[\adjustbox{max width=\linewidth}{$\displaystyle 
F,\Gamma\models_C^{\forall}\varphi
\quad\Longleftrightarrow\quad
B(F,\Gamma)\neq\varnothing
\;\land\;
B(F,\Gamma)\subseteq\llbracket\varphi\rrbracket_R.
$}\]

Every best viable realization supports \(\varphi\). This is strong
warrant or skeptical preferential consequence at the practical-verdict
layer.

The corresponding abstract relation is:

\[\adjustbox{max width=\linewidth}{$\displaystyle 
\Gamma\mid\!\sim_C\varphi
\quad\Longleftrightarrow\quad
\operatorname{Min}_{\prec}(V_F\cap\operatorname{Real}_R(\Gamma))
\subseteq\llbracket\varphi\rrbracket_R.
$}\]

\subsection{7.5 System P scope}\label{system-p-scope}

Under the standard definability and smoothness assumptions, the
fixed-base relation satisfies KLM System P (Kraus, Lehmann, \& Magidor,
1990):

\begin{itemize}
\tightlist
\item
  Reflexivity;
\item
  Left Logical Equivalence;
\item
  Right Weakening;
\item
  And;
\item
  Or;
\item
  Cautious Monotony;
\item
  Cut.
\end{itemize}

The \textbf{Or} rule requires the declared Boolean realization algebra.
If premises are represented only as unstructured sets of commitments,
formula-level disjunction has not yet been defined and the full rule set
cannot be claimed.

\subsection{7.6 Ranked models and Rational
Monotony}\label{ranked-models-and-rational-monotony}

If the preference relation is ranked and every relevant nonempty
definable set attains a minimum rank, the relation supports Rational
Monotony (Lehmann \& Magidor, 1992):

\[\adjustbox{max width=\linewidth}{$\displaystyle 
\alpha\mid\!\sim_C\beta
\quad\text{and}\quad
\alpha\not\mid\!\sim_C\neg\gamma
\quad\Longrightarrow\quad
\alpha\land\gamma\mid\!\sim_C\beta.
$}\]

A totally ordered rank set is insufficient when an infimum is not
attained. Minimum attainment is therefore explicit, not hidden in finite
examples.

\subsection{7.7 The Frege--Geach
boundary}\label{the-fregegeach-boundary}

Preferential consequence presupposes a language and realization
relation. It does not explain how a content acquires stable meaning
across assertion, negation, conditionals, and other embedded contexts.

CU therefore makes the narrower claim:

\begin{quote}
For commitment contents representable in a declared Boolean-closed
language with stable realization conditions, fixed-base preferential
semantics supplies a defeasible consequence structure that preserves
those contents through embedded inference.
\end{quote}

This is a claim about consequence, not a complete theory of meaning.
Ritual, prayer, poetry, myth, first-person avowal, and contemplative
utterance may remain outside the formal fragment.

An object-language conditional inside \(L_R\) is distinct from the
metalinguistic relation \(\mid\!\sim_C\). Confusing them would make the
formal apparatus appear to solve a compositionality problem it assumes
away.

\subsection{7.8 Model-relative determinacy and captured
models}\label{model-relative-determinacy-and-captured-models}

The formal relation is determinate relative to a declared model. It is
not neutral about who constructs the model.

Whoever controls:

\begin{itemize}
\tightlist
\item
  \(V_F\);
\item
  \(\prec\);
\item
  protected floors;
\item
  affected-agent scope;
\item
  evidence channels;
\item
  and presented options
\end{itemize}

can shape every formal verdict. Publicity makes capture visible; it does
not make it contestable.

Responsible use therefore depends on CU-G1--CU-G7: standing, option-set
integrity, independent challenge, correctability,
responsibility--authority alignment, participation, and reversibility.
These protections are normative and institutional additions, not
consequences of KLM form.

\subsection{7.9 Model revision}\label{model-revision}

Within a fixed model, added premises change only the premise-compatible
subset. A change to facts, floors, affected scope, horizon, viability
criterion, or ordering is \textbf{model revision} in the sense of moving
to a differently indexed evaluation model; it should not be confused
with premise extension inside one model (compare Alchourrón, Gärdenfors,
\& Makinson, 1985):

\[\adjustbox{max width=\linewidth}{$\displaystyle 
\mathcal M_{F,K,A,H,\prec}
\leadsto_{\mathrm{rev}}
\mathcal M_{F',K',A',H',\prec'}.
$}\]

The new model does not inherit the warrant of the old one merely by
replacing it. Every revision must be declared and re-evaluated.

\subsection{7.10 What the formal result
establishes}\label{what-the-formal-result-establishes}

The result establishes a disciplined form of defeasible consequence for
a specified addressable fragment. It does not establish:

\begin{itemize}
\tightlist
\item
  the meanings of the primitive contents;
\item
  the correctness of the viability base;
\item
  the justice of the ordering;
\item
  the adequacy of the option set;
\item
  the legitimacy of the model's authors;
\item
  or the practical permissibility of adoption.
\end{itemize}

Those questions belong to semantics, empirical inquiry, ethics,
governance, and affected-party review.

\section{8. Commitment Verdicts and Practical
Warrant}\label{commitment-verdicts-and-practical-warrant}

\subsection{8.1 Why consequence is not
enough}\label{why-consequence-is-not-enough}

Preferential consequence asks what holds throughout every best viable
realization. Practical commitment also asks what remains available, what
is excluded, whether anything viable remains, and whether the model is
adequate to decide.

One relation should not be asked to answer all of these questions.

\subsection{8.2 Five verdicts}\label{five-verdicts}

Let
\(B(F,\Gamma)=\operatorname{Min}_{\prec}(V_F\cap\operatorname{Real}_R(\Gamma))\).

\subsubsection{Strong warrant}\label{strong-warrant}

\[\adjustbox{max width=\linewidth}{$\displaystyle 
F,\Gamma\models_C^{\forall}\varphi
\iff
B(F,\Gamma)\neq\varnothing
\land
B(F,\Gamma)\subseteq\llbracket\varphi\rrbracket_R.
$}\]

Every best viable realization supports \(\varphi\).

\subsubsection{Admissibility and mere
admissibility}\label{admissibility-and-mere-admissibility}

\[\adjustbox{max width=\linewidth}{$\displaystyle 
F,\Gamma\models_C^{\exists}\varphi
\iff
B(F,\Gamma)\neq\varnothing
\land
B(F,\Gamma)\cap\llbracket\varphi\rrbracket_R\neq\varnothing.
$}\]

At least one best realization supports \(\varphi\). The verdict is
\textbf{merely admissible} when this holds but strong warrant does not.

\subsubsection{Exclusion}\label{exclusion}

\[\adjustbox{max width=\linewidth}{$\displaystyle 
F,\Gamma\models_C^{\bot}\varphi
\iff
B(F,\Gamma)\neq\varnothing
\land
B(F,\Gamma)\cap\llbracket\varphi\rrbracket_R=\varnothing.
$}\]

No best viable realization supports \(\varphi\).

\subsubsection{Infeasibility}\label{infeasibility}

\[\adjustbox{max width=\linewidth}{$\displaystyle 
V_F\cap\operatorname{Real}_R(\Gamma)=\varnothing.
$}\]

Nothing survives the declared facts, commitments, protected conditions,
affected scope, and horizon.

\subsubsection{Underdetermination}\label{underdetermination}

CU separates two sources of underdetermination. \textbf{Formal
non-attainment} occurs when the premise-compatible set is nonempty but
has no minimal elements under the declared ordering.
\textbf{Model-adequacy abstention} occurs before a best-set verdict is
issued because facts, affected-party scope, options, floors, horizon,
observation model, or ordering remain unspecified, contested, or
inadequate. Ties or incomparability within an attained best set do not
themselves create underdetermination; they can instead yield mere
admissibility.

Underdetermination is not permission by default. It calls for inquiry,
model comparison, caution, reversible action, or transparent unresolved
disagreement.

\subsection{8.3 Ties and the importance of mere
admissibility}\label{ties-and-the-importance-of-mere-admissibility}

When the best set is a singleton, mere admissibility cannot occur: the
sole best world either supports \(\varphi\) or does not. The distinction
becomes visible only with tied or incomparable best realizations.

This is not a technical curiosity. ``At least one good way of living
this commitment supports the extension'' and ``every good way supports
it'' are different practical situations.

\subsection{8.4 Classical negation as a modeling
choice}\label{classical-negation-as-a-modeling-choice}

Within the declared Boolean fragment and a nonempty best set:

\[\adjustbox{max width=\linewidth}{$\displaystyle 
F,\Gamma\models_C^{\bot}\varphi
\quad\Longleftrightarrow\quad
F,\Gamma\models_C^{\forall}\neg\varphi.
$}\]

CU adopts this equivalence for the formal fragment and states its cost.
``No best realization supports continuing this practice'' can be
practically different from ``every best realization positively supports
ceasing it.'' A non-classical account of practical negation remains
open.

\subsection{8.5 Maximum preservation}\label{maximum-preservation}

Where a scalar score is used and maxima exist, restricting the feasible
set cannot raise its maximum. Equality after adding \(\varphi\) means
only that at least one old optimum remains:

\[\adjustbox{max width=\linewidth}{$\displaystyle 
\operatorname{Argmax}_{W_{F,\Gamma}}c
\cap
\llbracket\varphi\rrbracket_R
\neq\varnothing.
$}\]

Maximum preservation therefore captures credulous admissibility, not
skeptical consequence.

\subsection{8.6 Practical expansion
warrant}\label{practical-expansion-warrant}

Formal consequence does not determine whether an agent or institution
may responsibly adopt a commitment. Let:

\[\adjustbox{max width=\linewidth}{$\displaystyle 
\mathcal W_C(\varphi\mid F,\Gamma,K,H,A,M,U)
$}\]

schematically denote practical expansion review.

Responsible adoption requires at least:

\begin{enumerate}
\def\labelenumi{\arabic{enumi}.}
\tightlist
\item
  compatibility with established fact;
\item
  a nonempty viable continuation;
\item
  compliance with protected floors;
\item
  represented affected parties and distributions;
\item
  no local integration purchased by disabling another's agency;
\item
  disclosed externalities across scales and horizons;
\item
  an option set not controlled to predetermine the result;
\item
  sufficiently independent challenge;
\item
  practical correction, appeal, and remedy where power is exercised;
\item
  robustness under reasonable model uncertainty;
\item
  revision and exit conditions;
\item
  retained option value, or a justified necessity-and-proportionality
  exception, under severe irreversible risk.
\end{enumerate}

This is a governance architecture, not one theorem or one scalar score.

\subsection{8.7 Trajectory worlds and transformative
commitment}\label{trajectory-worlds-and-transformative-commitment}

Transformative commitments may be represented within trajectory worlds
containing:

\begin{itemize}
\tightlist
\item
  temporary destabilization;
\item
  grief;
\item
  identity reorganization;
\item
  changed capacities;
\item
  new support;
\item
  later recovery;
\item
  branching futures;
\item
  and foreclosed options.
\end{itemize}

Premise addition selects trajectories in which the commitment is
undertaken. It does not silently redefine the base by which those
trajectories are evaluated.

\subsection{8.8 Bounded troughs}\label{bounded-troughs}

A transformative trajectory is not automatically warranted because it
promises later integration. Review must include:

\begin{itemize}
\tightlist
\item
  inviolable floors;
\item
  depth and duration of the trough;
\item
  credible recovery mechanisms;
\item
  lower-cost alternatives;
\item
  distribution of cost;
\item
  expiry conditions;
\item
  and retained option value.
\end{itemize}

No commitment becomes warranted merely because suffering can be narrated
as growth.

\subsection{8.9 Model change and warrant
inheritance}\label{model-change-and-warrant-inheritance}

A revised model may correct omitted facts or affected parties. But a
model change is a model change and cannot inherit the verdict of the
model it replaces without fresh evaluation.

Institutions must not manufacture underdetermination by omitting
alternatives, refusing to represent affected people, or supplying an
ordering too coarse to discriminate.

\subsection{8.10 Disagreement}\label{disagreement-1}

CU distinguishes:

\begin{itemize}
\tightlist
\item
  descriptive disagreement;
\item
  model disagreement;
\item
  constraint disagreement;
\item
  ordering disagreement;
\item
  horizon disagreement;
\item
  affected-agent disagreement;
\item
  formal-application error;
\item
  and practical conflict among separately admissible commitments.
\end{itemize}

Transparent disagreement can be a legitimate endpoint. The framework
requires disclosed assumptions and fair challenge; it does not promise
adjudication where substantive premises remain contested.

\subsection{8.11 Provisional criterion}\label{provisional-criterion}

\begin{quote}
A commitment is responsibly warranted only to the degree that it remains
compatible with established facts, preserves protected viability and
agency across affected systems, does not require systematic
self-deception or hidden externalization, and remains open to
correction, expiry, exit, and repair.
\end{quote}

This is a multi-constraint framework under continuing review, not a
representation theorem for the whole of practical life.

\section{9. Dynamic and Probabilistic
Extensions}\label{dynamic-and-probabilistic-extensions}

\subsection{9.1 Why worlds become
trajectories}\label{why-worlds-become-trajectories}

Many CU claims concern development, repair, collapse, and path
dependence. A static state can conceal whether apparent order is
self-sustaining, externally subsidized, or about to fail. It can also
misclassify productive disruption as incoherence.

Let a transition system be

\[\adjustbox{max width=\linewidth}{$\displaystyle 
\mathcal S=(X,A,\rightarrow,x_0,L),
$}\]

where \(X\) is a state space, \(A\) an action or intervention set,
\(\rightarrow\) an admissible transition relation, \(x_0\) an initial
condition, and \(L\) a labeling function.

A finite trajectory is

\[\adjustbox{max width=\linewidth}{$\displaystyle 
\tau=(x_0,a_0,x_1,a_1,\ldots,x_T),
$}\]

and an infinite trajectory continues without a terminal index. Temporal
formulas may express safety, eventual recovery, recurrence, or
irreversible collapse.

\subsection{9.2 Dynamic preferential
consequence}\label{dynamic-preferential-consequence}

Let \(\Omega\) be a set of admissible trajectories,
\(V_F^{\Omega}\subseteq\Omega\) the fixed trajectory-level viability
base for one evaluation context, and \(\prec_\Omega\) a preference
relation fixed within that model. Define

\[\adjustbox{max width=\linewidth}{$\displaystyle 
W^{\Omega}_{F,\Gamma}
=
V_F^{\Omega}\cap\operatorname{Real}_R^{\Omega}(\Gamma),
$}\]

and

\[\adjustbox{max width=\linewidth}{$\displaystyle 
\Gamma\mid\!\sim_C^{\Omega}\varphi
\quad\Longleftrightarrow\quad
\operatorname{Min}_{\prec_\Omega}
\bigl(W^{\Omega}_{F,\Gamma}\bigr)
\subseteq
\operatorname{Real}_R^{\Omega}(\varphi).
$}\]

\textbf{{[}Formal result{]}} When trajectories are treated as worlds,
premises filter the fixed trajectory base intersectively, and the
ordinary definability and smoothness conditions hold, the preferential
representation and its System P properties carry over. If the trajectory
ordering is ranked and every relevant nonempty definable trajectory set
attains a minimum, System R carries over.

The result is semantic. It does not require that the preference be
computed from one additive trajectory score. It also does not apply when
\(V_F^{\Omega}\), the horizon, the affected-agent scope, or
\(\prec_\Omega\) is recomputed non-intersectively as the premise
changes; that is model revision or premise-responsive selection.

\subsection{9.3 Temporal criteria}\label{temporal-criteria}

Dynamic evaluation can distinguish:

\begin{itemize}
\tightlist
\item
  \textbf{safety:} a protected floor is never violated;
\item
  \textbf{viability:} the trajectory remains within a viable region or
  returns within a permitted recovery time;
\item
  \textbf{liveness:} a required capacity or opportunity eventually
  becomes available;
\item
  \textbf{repairability:} after damage, a route back to acceptable
  organization remains;
\item
  \textbf{resilience:} function persists across a perturbation class;
\item
  \textbf{drift:} a relevant measure degrades gradually without crossing
  an immediate failure boundary;
\item
  \textbf{collapse:} recovery becomes impossible within the modeled
  horizon or intervention class.
\end{itemize}

These properties are not interchangeable. A system can be safe but
stagnant, lively but unsafe, or viable only because another system
absorbs its damage.

\subsection{9.4 Horizon dependence}\label{horizon-dependence}

A trajectory can appear optimal at one horizon and destructive at
another. Let

\[\adjustbox{max width=\linewidth}{$\displaystyle 
J_H(\tau)
$}\]

be a loss or evaluation profile over horizon \(H\). The ranking may
change when \(H\) changes.

Every long-run claim must therefore state whether it concerns:

\begin{itemize}
\tightlist
\item
  a fixed finite horizon;
\item
  a discounted indefinite horizon;
\item
  a worst-case stopping time;
\item
  a regenerative cycle;
\item
  or asymptotic behavior.
\end{itemize}

Discounting is not neutral. A high discount rate can mathematically hide
costs imposed on later persons or ecosystems.

\subsection{9.5 Stochastic dynamics}\label{stochastic-dynamics}

In uncertain environments, a policy \(\pi\) induces a probability
measure \(P_\pi\) over trajectories. A random loss variable

\[\adjustbox{max width=\linewidth}{$\displaystyle 
J:\Omega\rightarrow\mathbb R\cup\{\infty\}
$}\]

can be summarized by a risk functional

\[\adjustbox{max width=\linewidth}{$\displaystyle 
R(\pi)=\rho(J;P_\pi).
$}\]

Possible choices include expectation, worst-case loss, conditional value
at risk, chance constraints, and distributionally robust bounds. The
choice reflects substantive attitudes toward tail risk, uncertainty, and
protected floors. It must not be disguised as a purely technical
parameter.

\subsection{9.6 Probability, typicality, and
preference}\label{probability-typicality-and-preference}

Three orderings must remain distinct.

\begin{itemize}
\tightlist
\item
  \textbf{Probability} concerns how likely outcomes are under a model
  and policy.
\item
  \textbf{Typicality} selects outcomes treated as normal or
  representative for an inference.
\item
  \textbf{Preference} evaluates outcomes relative to constraints or
  aims.
\end{itemize}

A likely outcome can be bad. A preferred outcome can be unlikely. A
typicality rule can exclude rare catastrophes that a safety analysis
must retain.

For each policy, a typicality operator may select

\[\adjustbox{max width=\linewidth}{$\displaystyle 
T_\pi(E)\subseteq E
$}\]

from an event \(E\) with positive probability. Per-policy typical
consequence can then be defined by

\[\adjustbox{max width=\linewidth}{$\displaystyle 
\alpha\sim_\pi\beta
\quad\Longleftrightarrow\quad
P_\pi\!\left(
T_\pi(\llbracket\alpha\rrbracket)
\setminus
\llbracket\beta\rrbracket
\right)=0.
$}\]

Its rational properties depend on the axioms imposed on \(T_\pi\).

\subsection{9.7 Feasibility thresholds are not Boolean
semantics}\label{feasibility-thresholds-are-not-boolean-semantics}

A common policy-feasibility map is

\[\adjustbox{max width=\linewidth}{$\displaystyle 
\Pi_\alpha
=
\{\pi:P_\pi(\alpha)\geq1-\delta\}.
$}\]

Logical equivalence is preserved, but conjunction and disjunction do not
generally map to set intersection and union:

\[\adjustbox{max width=\linewidth}{$\displaystyle 
\Pi_{\alpha\land\beta}
\neq
\Pi_\alpha\cap\Pi_\beta,
\qquad
\Pi_{\alpha\lor\beta}
\neq
\Pi_\alpha\cup\Pi_\beta
$}\]

in general.

For conjunction, separate high-probability satisfaction does not
guarantee that the joint event clears the same threshold. For
disjunction, a policy may make \(\alpha\lor\beta\) certain while
splitting probability between the disjuncts so that neither individually
clears the threshold.

This Boolean mismatch is the central obstruction to transferring KLM
proofs through a premise-dependent policy selector.

\subsection{9.8 Robustness under model
uncertainty}\label{robustness-under-model-uncertainty}

A probability measure is itself a model-dependent object. When
transition dynamics, observations, or human responses are uncertain,
policy evaluation should range over a model set \(\mathcal P\):

\[\adjustbox{max width=\linewidth}{$\displaystyle 
R_{\mathrm{rob}}(\pi)
=
\sup_{P\in\mathcal P}\rho(J;P,\pi)
$}\]

or use another disclosed robustness rule.

The size and construction of \(\mathcal P\) matters. An excessively
broad set can paralyze action; an excessively narrow set converts
ignorance into false precision.

\subsection{9.9 Rare harms and protected
floors}\label{rare-harms-and-protected-floors}

Expected-value optimization can trade a low-probability catastrophic
loss against many small gains. Where an outcome would destroy agency,
violate rights, or irreversibly close future options, a protected
constraint may be more appropriate:

\[\adjustbox{max width=\linewidth}{$\displaystyle 
P_\pi(\text{floor violation})\leq\varepsilon,
$}\]

or, in some cases,

\[\adjustbox{max width=\linewidth}{$\displaystyle 
P_\pi(\text{floor violation})=0
$}\]

within the modeled uncertainty class.

No mathematical form can determine which floors deserve protection. That
is an ethical and institutional question supplied to the decision model.

\subsection{9.10 Dynamic commitment
warrant}\label{dynamic-commitment-warrant}

For commitments that transform the agent, the object of evaluation is a
trajectory through changing states, capacities, relationships, and
available options. The transformation is represented \textbf{inside the
trajectory world}. A trajectory may contain a bounded trough of reduced
integration, later recovery, newly acquired capacities, or foreclosed
possibilities while remaining assessable by a fixed trajectory-level
viability predicate \(V_F\). The proposition representing adoption of
the commitment selects those trajectories in which the adoption occurs;
it does not redefine \(V_F\) merely by being added as a premise.

A dynamic warrant assessment should ask whether the commitment:

\begin{itemize}
\tightlist
\item
  preserves a route to revision;
\item
  increases or narrows future option value;
\item
  creates dependencies or power asymmetries;
\item
  changes the reliability of the agent's own self-assessment;
\item
  and alters who bears the costs of exit.
\end{itemize}

If the analyst changes the protected floors, affected-agent scope,
horizon, or preference ordering in response to the commitment, the
comparison is between differently indexed models. That may be
legitimate, but it is model revision rather than an application of the
fixed-base System P theorem. This distinction is especially important
for religious communities, intimate relationships, political identities,
and persistent artificial agents.

\subsection{9.11 Limits of formalization}\label{limits-of-formalization}

Dynamic and probabilistic models can make assumptions visible, but they
do not eliminate judgment. The state representation, horizon,
probability model, protected floors, and loss functional are themselves
contestable.

The purpose of formalization is not to replace deliberation. It is to
expose where deliberation has entered the machinery.

\begin{center}\rule{0.5\linewidth}{0.5pt}\end{center}

\section{10. The Two-Stage Decision
Architecture}\label{the-two-stage-decision-architecture}

\subsection{10.1 The architecture}\label{the-architecture}

The inherited two-stage architecture separates policy selection from
consequence assessment.

\textbf{Stage 1: policy selection.} For premise or goal condition
\(\alpha\), determine a feasible policy set \(\Pi_\alpha\) and select
the policies minimizing a specified risk:

\[\adjustbox{max width=\linewidth}{$\displaystyle 
f(\alpha)
=
\operatorname{argmin}_{\pi\in\Pi_\alpha}R(\pi).
$}\]

\textbf{Stage 2: policy-indexed typical consequence.} For each selected
policy, infer what holds in its typical \(\alpha\)-trajectories:

\[\adjustbox{max width=\linewidth}{$\displaystyle 
\alpha\sim_\pi\beta.
$}\]

The aggregate relation is

\[\adjustbox{max width=\linewidth}{$\displaystyle 
\alpha\Vdash_{2S}\beta
\quad\Longleftrightarrow\quad
\forall\pi\in f(\alpha),
\;\alpha\sim_\pi\beta.
$}\]

This architecture is useful: it prevents a prediction model from
silently choosing the policy and prevents a policy score from being
mistaken for a forecast.

\subsection{10.2 What remains valid}\label{what-remains-valid}

Under ordinary extensionality and per-policy consequence assumptions,
the aggregate relation can support premise-fixed properties such as:

\begin{itemize}
\tightlist
\item
  Reflexivity;
\item
  Left Logical Equivalence;
\item
  Right Weakening; and
\item
  And.
\end{itemize}

These arguments do not require comparing selected policy sets under
materially different premises.

\subsection{10.3 The Boolean obstruction}\label{the-boolean-obstruction}

The premise-varying KLM axioms require relations among

\[\adjustbox{max width=\linewidth}{$\displaystyle 
f(\alpha),
\quad
f(\beta),
\quad
f(\alpha\lor\beta),
\quad
f(\alpha\land\beta).
$}\]

Because feasibility is defined over policy-induced distributions, the
map \(\alpha\mapsto\Pi_\alpha\) need not preserve Boolean operations. A
mixed policy can satisfy \(\alpha\lor\beta\) by distributing support
across the disjuncts while satisfying neither disjunct individually.
Probability thresholds can create additional failures, including for
conjunction, but they are not the source of the Or obstruction.
Consequently, the selector identities needed for Or, Cautious Monotony,
Cut, and Rational Monotony do not generally hold.

\textbf{{[}Formal result{]}} The two-stage relation does not in general
satisfy System P or System R. The positive theorem stated in the main
body of the earlier M\&E paper is withdrawn. Its own later appendix
contained the correct obstruction.

\subsection{10.4 A minimal counterexample for
Or}\label{a-minimal-counterexample-for-or}

Let \(A\) and \(B\) be disjoint events and take three policies:

\begin{itemize}
\tightlist
\item
  \(\pi_A\) makes \(A\) certain and has risk 10;
\item
  \(\pi_B\) makes \(B\) certain and has risk 10;
\item
  \(\pi_M\) splits probability between \(A\) and \(B\), makes
  \(A\lor B\) certain, and has risk 1.
\end{itemize}

With certainty feasibility,

\[\adjustbox{max width=\linewidth}{$\displaystyle 
f(A)=\{\pi_A\},
\qquad
f(B)=\{\pi_B\},
\qquad
f(A\lor B)=\{\pi_M\}.
$}\]

A conclusion may hold typically under \(\pi_A\) and \(\pi_B\) while
failing under \(\pi_M\). Therefore the Or rule need not hold for the
aggregate relation.

The failure is structural, not a peculiarity of the coherence objective.

\subsection{10.5 Why the earlier proof
failed}\label{why-the-earlier-proof-failed}

The earlier proof stipulated identities such as

\[\adjustbox{max width=\linewidth}{$\displaystyle 
\Pi_{\alpha\lor\beta}
=
\Pi_\alpha\cup\Pi_\beta
$}\]

and

\[\adjustbox{max width=\linewidth}{$\displaystyle 
\Pi_{\alpha\land\beta}
=
\Pi_\alpha\cap\Pi_\beta.
$}\]

Those identities hold for classical model sets:

\[\adjustbox{max width=\linewidth}{$\displaystyle 
\operatorname{Mod}(\alpha\lor\beta)
=
\operatorname{Mod}(\alpha)
\cup
\operatorname{Mod}(\beta),
$}\]

but not in general for policy-induced distributions. A mixed policy can
satisfy a disjunction while satisfying neither disjunct. The proof
therefore imported the Boolean structure of valuation semantics into a
different object.

The correction illustrates the role of explicit translation rules:
formal properties do not survive a mapping merely because the notation
looks similar.

\subsection{10.6 What the architecture
is}\label{what-the-architecture-is}

The two-stage construction should be described as a \textbf{decision
architecture} with inspectable components:

\begin{enumerate}
\def\labelenumi{\arabic{enumi}.}
\tightlist
\item
  premise or objective interpretation;
\item
  feasibility assessment;
\item
  risk-sensitive policy selection;
\item
  policy-conditioned prediction;
\item
  typicality or tail analysis;
\item
  protected-floor check;
\item
  human or institutional authorization.
\end{enumerate}

Its quality is assessed through decision theory, calibration,
robustness, safety, and governance. It need not be a single KLM
consequence relation to be valuable.

\subsection{10.7 Routes to stronger logical
structure}\label{routes-to-stronger-logical-structure}

Stronger results may be recovered under additional conditions.

\subsubsection{Fixed-policy consequence}\label{fixed-policy-consequence}

If \(\pi\) is held fixed, per-policy typical consequence can satisfy
preferential or ranked axioms when its typicality semantics does.

\subsubsection{Support-based feasibility and its
limit}\label{support-based-feasibility-and-its-limit}

If feasibility is defined by support inclusion,

\[\adjustbox{max width=\linewidth}{$\displaystyle 
\Pi_\alpha=\{\pi:\operatorname{Supp}(P_\pi)\subseteq\llbracket\alpha\rrbracket\},
$}\]

then the conjunction identity is restored:

\[\adjustbox{max width=\linewidth}{$\displaystyle 
\Pi_{\alpha\land\beta}=\Pi_\alpha\cap\Pi_\beta.
$}\]

The disjunction identity still need not hold when mixed policies are
admitted: a policy can place support partly in \(\alpha\) and partly in
\(\beta\), thereby belonging to \(\Pi_{\alpha\lor\beta}\) without
belonging to either \(\Pi_\alpha\) or \(\Pi_\beta\). Restoring Or
therefore requires an additional restriction, such as deterministic or
point-mass policies, a premise fragment closed under the implemented
feasibility semantics, or the joint policy-trajectory construction
below.

\subsubsection{Joint policy--trajectory
worlds}\label{joint-policytrajectory-worlds}

One may construct worlds \((\pi,\tau)\) and define a single fixed
preference over the joint space. Ordinary preferential consequence can
then be applied directly. This avoids a premise-dependent change of
semantic universe, but the preference construction and treatment of
policy choice require careful justification.

\subsubsection{Restricted premise
classes}\label{restricted-premise-classes}

A domain may identify a fragment of premises for which feasibility
behaves compositionally. The stronger axioms would then hold only on
that fragment.

None of these routes should be assumed without proof.

\subsection{10.8 Alignment implications}\label{alignment-implications}

For artificial systems, the correction is protective. It prevents a risk
optimizer from advertising the formal guarantees of a non-monotonic
logic that its aggregate decision process does not possess.

The architecture can still enforce important separations:

\begin{itemize}
\tightlist
\item
  facts and forecasts remain distinct from values and permissions;
\item
  policy selection remains distinct from consequence generation;
\item
  rare harms remain visible rather than being removed by typicality;
\item
  and authorizing responsibility remains outside the predictive model.
\end{itemize}

The v0.10 AI paper should state guarantees layer by layer rather than
assigning one theorem to the whole stack.

\subsection{10.9 Relation to rationalization
detection}\label{relation-to-rationalization-detection}

A system can follow the two-stage procedure and still rationalize. It
may define feasibility to exclude inconvenient evidence, choose a risk
functional that protects a preferred policy, or select a typicality
threshold that removes adverse outcomes.

Inference integrity therefore requires an audit of how the premises,
constraints, and objective were produced. Rationalization detection is a
meta-level mechanism, not a corollary of policy optimization.

\subsection{10.10 Canonical conclusion}\label{canonical-conclusion}

\begin{quote}
The two-stage architecture is a valuable modular decision procedure. Its
per-policy inferential components may possess KLM properties, but the
premise-dependent aggregate selector does not generally satisfy System P
or System R. Any stronger claim requires additional structural
assumptions proved for the implemented semantics.
\end{quote}

\begin{center}\rule{0.5\linewidth}{0.5pt}\end{center}

\section{11. The Epistemology Ladder}\label{the-epistemology-ladder}

\subsection{11.1 Purpose of the ladder}\label{purpose-of-the-ladder}

The Epistemology Ladder organizes capacities required for increasingly
explicit and self-correcting forms of knowing. It is not a scale of
human worth, a claim that formal knowledge is always superior to
embodied skill, or a theory that truth changes by level.

Higher levels depend on lower capacities and can fail while those
capacities remain. Formal proof depends on symbol discrimination and
stable rules; science depends on observation, instruments, testimony,
and institutions; wisdom depends on but cannot be reduced to accumulated
propositions.

Version 0.10 adopts seven levels, EL0 through EL6.

\subsection{11.2 EL0 --- Constraint-sensitive
engagement}\label{el0-constraint-sensitive-engagement}

At EL0 a system changes in discriminating response to constraints.
Examples include sensorimotor regulation, error-sensitive adaptation,
and embodied salience.

There need not yet be a proposition, explicit model, or truth claim. The
epistemic achievement is contact: differences in the world make
differences to the system in a sufficiently reliable way.

\textbf{Failure modes:} insulation from feedback, uncontrolled
confabulation-like activity, and perturbations that exceed regulatory
capacity.

\textbf{Important limit:} successful regulation does not by itself
establish consciousness, representation, or knowledge in the
propositional sense.

\subsection{11.3 EL1 --- Situated pattern and
proto-representation}\label{el1-situated-pattern-and-proto-representation}

At EL1, recurrent constraint-sensitive patterns stabilize into
discriminations, expectations, maps, or habits that can guide action
across similar situations.

The patterns may remain implicit. They can be more or less reliable
without being expressible as sentences.

\textbf{Failure modes:} rigid habits, overgeneralization, proxy
fixation, and internally stable patterns decoupled from changing
conditions.

EL1 preserves the insight that intelligent engagement begins before
explicit assertion.

\subsection{11.4 EL2 --- World-coupled
claims}\label{el2-world-coupled-claims}

At EL2, a system forms or uses representations as claims about how
things are. Truth and falsity become directly applicable because the
representation is staked against a world capable of contradicting it.

A claim is not made true by usefulness in the immediate situation.
Practical success is evidence only where a justified relation connects
success to the represented state of affairs.

\textbf{Failure modes:} confirmation bias, selective exposure,
overconfidence, hallucination, and protecting a claim by weakening its
contact with possible correction.

\subsection{11.5 EL3 --- Diachronic and intersubjective
correction}\label{el3-diachronic-and-intersubjective-correction}

At EL3, claims are revised across time and exposed to other knowers.
Memory, testimony, criticism, replication, and social trust permit
errors to be detected that no isolated episode or agent could identify.

Intersubjective agreement is not truth. Communities can stabilize
falsehood. Its epistemic value lies in distributed access, adversarial
checking, and preserved correction channels.

\textbf{Failure modes:} groupthink, propaganda, prestige cascades,
corrupted testimony, institutional amnesia, and updating without
integration.

\subsection{11.6 EL4 --- Explicit models and formal
systems}\label{el4-explicit-models-and-formal-systems}

At EL4, assumptions, variables, inferential rules, and scope conditions
become explicit. This level has two mutually supporting modes.

\subsubsection{EL4-E: empirical
modeling}\label{el4-e-empirical-modeling}

Claims are operationalized, measured, compared with observations, and
tested across stated domains. Models can be predictively successful
while ontologically nonunique.

\subsubsection{EL4-F: formal
articulation}\label{el4-f-formal-articulation}

Conclusions are derived from explicit axioms and rules. Proof
establishes validity relative to the formal system; it does not by
itself establish that the axioms describe nature.

Neither branch is simply above the other. Empirical science uses formal
systems; formal work depends on interpretations and applications
supplied from outside the proof calculus.

\textbf{Failure modes:} model overreach, Goodhart effects, false
precision, confusing proof with empirical truth, and treating unmodeled
factors as nonexistent.

\subsection{11.7 EL5 --- Epistemic
governance}\label{el5-epistemic-governance}

At EL5, a system explicitly manages how claims move among levels,
models, communities, and decisions. It establishes:

\begin{itemize}
\tightlist
\item
  evidence standards;
\item
  source and citation practices;
\item
  version control;
\item
  translation rules;
\item
  uncertainty labels;
\item
  conflict-of-interest procedures;
\item
  revision triggers;
\item
  and authority boundaries.
\end{itemize}

This is the principal level at which CU v0.10 operates. Its task is not
to replace domain methods but to govern their integration.

\textbf{Failure modes:} dogmatism, drift without consolidation,
bureaucratic formalism, selective enforcement of standards, and
meta-level rules designed to protect a preferred outcome.

\subsection{11.8 EL6 --- Meta-epistemic
awareness}\label{el6-meta-epistemic-awareness}

At EL6, the system represents the limits and distortion pressures of its
own epistemic governance. It can ask how its incentives, identity,
scale, embodiment, data access, and preferred formalism shape what it is
able to notice.

This level is associated with intellectual humility and wisdom, but it
does not confer infallibility. A claim to meta-awareness can itself
become performative insulation.

\textbf{Failure modes:} epistemic arrogance, total skepticism, paralysis
by reflection, ritualized humility without behavioral correction, and
treating awareness of bias as immunity from bias.

\subsection{11.9 The ladder is not a truth
hierarchy}\label{the-ladder-is-not-a-truth-hierarchy}

Truth does not become ``more true'' at higher levels. A situated
observation can be true; a formal theorem can be irrelevant to reality;
a community consensus can be false.

What changes across the ladder is the structure of access,
stabilization, explicitness, correction, and governance.

This correction is important for Tarskian language. Object-language and
metalanguage must be distinguished where semantic self-reference
matters, but that does not make truth relative to an epistemic rung.

\subsection{11.10 Upward dependence and downward
support}\label{upward-dependence-and-downward-support}

Higher epistemic capacities depend on lower ones. Formal reasoning
depends on perception, memory, and symbol handling. Institutional
science depends on trust, material support, and social permission to
report error.

Support also flows downward. Well-governed institutions can protect
embodied safety and local inquiry. Formal models can improve perception
and action. The relation is therefore a dependency network, not a
one-way status hierarchy.

\subsection{11.11 Artificial systems}\label{artificial-systems}

Artificial systems can display capacities at different levels without
reproducing the human path through them. A language model may perform
EL4-like formal manipulations while having limited persistent EL3
learning. A robot may display strong EL0 and EL1 coupling while weak in
explicit justification. A multi-agent system may simulate institutional
correction while lacking a stable bearer of responsibility.

The ladder should be used diagnostically, not anthropomorphically. Each
level requires operational tests tailored to the architecture.

\subsection{11.12 Contemplative practice}\label{contemplative-practice}

Contemplative practice can refine EL0 salience, reveal EL1 habits, make
EL2 claims less fused with identity, improve EL5 governance of
attention, and deepen EL6 awareness of limits. It can also produce
confidence unsupported by descriptive evidence.

CU treats contemplative reports as data about experience and practice,
not as automatic proof of metaphysical interpretations. Their value
increases when first-person, behavioral, physiological, and
intersubjective methods are brought into disciplined relation.

\subsection{11.13 Epistemic maturity}\label{epistemic-maturity}

Maturity is not permanent residence at the highest rung. It is the
capacity to move appropriately among them:

\begin{itemize}
\tightlist
\item
  to trust embodied warning without treating every feeling as fact;
\item
  to make situated judgments without universalizing them;
\item
  to use social knowledge without surrendering correction;
\item
  to rely on models without forgetting their scope;
\item
  to value proof without mistaking axioms for the world;
\item
  and to commit while remaining capable of revision.
\end{itemize}

\subsection{11.14 Canonical table}\label{canonical-table}

\begin{longtable}[]{@{}
  >{\raggedright\arraybackslash}p{(\columnwidth - 6\tabcolsep) * \real{0.2500}}
  >{\raggedright\arraybackslash}p{(\columnwidth - 6\tabcolsep) * \real{0.2500}}
  >{\raggedright\arraybackslash}p{(\columnwidth - 6\tabcolsep) * \real{0.2500}}
  >{\raggedright\arraybackslash}p{(\columnwidth - 6\tabcolsep) * \real{0.2500}}@{}}
\toprule\noalign{}
\begin{minipage}[b]{\linewidth}\raggedright
Level
\end{minipage} & \begin{minipage}[b]{\linewidth}\raggedright
Name
\end{minipage} & \begin{minipage}[b]{\linewidth}\raggedright
Primary achievement
\end{minipage} & \begin{minipage}[b]{\linewidth}\raggedright
Principal correction
\end{minipage} \\
\midrule\noalign{}
\endhead
\bottomrule\noalign{}
\endlastfoot
EL0 & Constraint-sensitive engagement & Responsive contact & Contact is
not yet a claim \\
EL1 & Situated pattern & Stable discrimination or proto-model & Internal
stability is not truth \\
EL2 & World-coupled claims & Truth-apt representation & Usefulness alone
is not accuracy \\
EL3 & Diachronic and intersubjective correction & Learning and
distributed checking & Consensus is not truth \\
EL4 & Explicit models and formal systems & Operationalization and proof
& Models and proofs retain scope \\
EL5 & Epistemic governance & Managed translation, revision, and
authority & Governance can itself be captured \\
EL6 & Meta-epistemic awareness & Awareness of limits and distortion
pressures & Humility is not paralysis or immunity \\
\end{longtable}

\begin{center}\rule{0.5\linewidth}{0.5pt}\end{center}

\section{12. Governance of Inquiry and
Action}\label{governance-of-inquiry-and-action}

\subsection{12.1 Why governance belongs in metaphysics and
epistemology}\label{why-governance-belongs-in-metaphysics-and-epistemology}

Inquiry is never only a relation between a proposition and evidence.
Consequential inquiry also selects a system boundary, decides which
observations count, constructs categories, chooses options, distributes
uncertainty, and authorizes action. These decisions shape what can be
known and who bears the cost of error.

The completed v0.10 domain cycle therefore adds a governance layer to
the framework. This layer is not derived from the Fisher geometry of the
Physics sandbox, from biological viability, or from formal preferential
logic. It is defended through explicit normative and methodological
arguments about affected standing, agency, noncompensation,
correctability, responsibility, legitimacy, and severe risk.

The seven CU-G principles govern the use of CU across domains.

\subsection{12.2 CU-G1 --- Affected scope, distribution, and protected
conditions}\label{cu-g1-affected-scope-distribution-and-protected-conditions}

Any consequential evaluation must identify affected beings and groups,
the distribution of benefits and burdens, protected conditions or
floors, and whether gains for some are being used to compensate harms to
others.

Aggregate improvement does not establish admissibility. A mean can rise
while one group falls below a protected threshold. A psychologically
stabilizing intervention can restrict later agency. An efficient
institution can externalize ecological or labor costs. A capable AI
deployment can benefit represented users while omitting people absent
from its feedback stream.

CU-G1 therefore requires at least:

\begin{itemize}
\tightlist
\item
  an affected-party map;
\item
  disaggregated outcomes where distribution matters;
\item
  explicit protected floors;
\item
  treatment of future and externalized burdens;
\item
  and a justification for any permitted tradeoff.
\end{itemize}

\subsection{12.3 CU-G2 --- Option-set and framing
integrity}\label{cu-g2-option-set-and-framing-integrity}

Every screening or ranking result is conditional on the problem framing
and the options presented.

Let

\[\adjustbox{max width=\linewidth}{$\displaystyle 
U^{\mathrm{pres}}\subseteq U^{\mathrm{feas}}
$}\]

represent the presented and broader feasible sets. A procedure can
operate perfectly on \(U^{\mathrm{pres}}\) while the institution
controlling the menu excludes no action, delay, staged intervention,
less harmful means, structural reform, or proposals from affected
parties. Floor compliance therefore establishes only screening
admissibility relative to the presented set.

Option-set integrity asks:

\begin{itemize}
\tightlist
\item
  who generated and excluded alternatives;
\item
  which constraints were treated as fixed and why;
\item
  whether affected parties could propose alternatives;
\item
  whether inaction, delay, staging, containment, and redesign were
  considered;
\item
  and whether an apparently necessary choice is an artifact of
  institutional framing.
\end{itemize}

\subsection{12.4 CU-G3 --- Practical correctability and independent
challenge}\label{cu-g3-practical-correctability-and-independent-challenge}

Correctability is not the existence of a complaint form. It is the
practical capacity to detect consequential error, preserve contrary
evidence, challenge a decision through a sufficiently independent
channel, stop or reverse action, correct records, provide remedy, track
recurrence, revise rules, and learn.

An institution cannot establish its own correctability solely through
statistics, reviewers, records, and success criteria that it controls.
Nor does technical modularity establish institutional independence when
one operator appoints every reviewer and controls every remedy.

Practical correctability should be evaluated through access, retaliation
risk, reviewer independence, reversal authority, timeliness, record
correction, remedy, recurrence, affected-party testimony, and external
or jointly governed validation.

\subsection{12.5 CU-G4 --- Responsibility--authority--causal-reach
alignment}\label{cu-g4-responsibilityauthoritycausal-reach-alignment}

Responsibility is a profile rather than a scalar. Relevant dimensions
include causal contribution, knowledge, foreseeability, control,
authority, benefit, coercion, prevention capacity, and repair capacity.

The person nearest an event may have little power to change the
governing conditions. The institution farthest from the event may
possess the authority, knowledge, benefit, and capacity to prevent
recurrence. CU therefore rejects both responsibility without authority
and authority without accountability.

Governance burden increases with causal reach, environmental access,
persistence, delegation, irreversibility, and the ability to externalize
error.

\subsection{12.6 CU-G5 --- Participation proportional to governance
reach}\label{cu-g5-participation-proportional-to-governance-reach}

Principles and systems that classify, burden, represent, or govern
people require participation by affected communities proportionate to
their reach and risk. Where the stakes are high, participation must
include substantive authority rather than symbolic consultation.

Participation does not replace evidence, proof, clinical judgment, or
moral argument. Required review gates are conjunctive within their
competence: an empirical review cannot settle a question of legitimate
standing, and a community process cannot make a false factual claim
true. When required gates remain in unresolved conflict, the principle
remains provisional and the objection remains public.

\subsection{12.7 CU-G6 --- Patiency, agency, responsibility, and
sovereignty are
distinct}\label{cu-g6-patiency-agency-responsibility-and-sovereignty-are-distinct}

Evidence of experience or welfare may ground moral patiency. Capacities
for reason-guided action may support agency. Responsibility additionally
depends on knowledge, control, continuity, and answerability.
Sovereignty requires legitimate social authorization.

Thus no single status licenses any of the others:

\[\adjustbox{max width=\linewidth}{$\displaystyle 
\begin{aligned}
\text{patiency} &\not\Rightarrow \text{agency}, &
\text{patiency} &\not\Rightarrow \text{responsibility}, &
\text{patiency} &\not\Rightarrow \text{sovereignty},\\
\text{agency} &\not\Rightarrow \text{patiency}, &
\text{agency} &\not\Rightarrow \text{responsibility}, &
\text{agency} &\not\Rightarrow \text{sovereignty},\\
\text{responsibility} &\not\Rightarrow \text{patiency}, &
\text{responsibility} &\not\Rightarrow \text{agency}, &
\text{responsibility} &\not\Rightarrow \text{sovereignty},\\
\text{sovereignty} &\not\Rightarrow \text{patiency}, &
\text{sovereignty} &\not\Rightarrow \text{agency}, &
\text{sovereignty} &\not\Rightarrow \text{responsibility}.
\end{aligned}
$}\]

The implications can also fail in other directions. Non-conscious
institutions and software can wield enormous power. A conscious
artificial system, should one exist, would not thereby acquire authority
over public values.

\subsection{12.8 CU-G7 --- Reversibility, containment, and retained
option
value}\label{cu-g7-reversibility-containment-and-retained-option-value}

Under credible risk of severe or irreversible harm, an application must
preserve context-appropriate reversibility, containment, rollback,
staged exposure, or equivalent retained option value, or publicly
justify a necessity-and-proportionality exception that includes the
harms of delay.

Reversibility is not an absolute. Surgery, rescue, testimony, ecological
intervention, and delay can all have irreversible consequences. The
requirement is a floor-or-exception architecture: retain options where
possible; where not possible, disclose the threat, alternatives,
necessity, proportionality, and decision authority.

\subsection{12.9 Relationship to formal
warrant}\label{relationship-to-formal-warrant}

The CU-G principles are not generated by KLM preferential logic. Formal
semantics can represent a fixed viability base and ordering, but it
cannot decide which beings count, which floors are legitimate, who has
authority, or which harms may be traded. Those are substantive normative
and institutional questions.

Preferential consequence can therefore license the structure of an
inference only after the model's normative content has been defended.
CU-G1--G7 govern that defense.

\subsection{12.10 Relationship to scientific
modeling}\label{relationship-to-scientific-modeling}

The governance family also applies before action. A scientific model can
omit populations, hide distributional effects, treat one feedback stream
as complete, or compare only options supplied by a sponsor. These
defects are epistemic as well as ethical because they alter what the
inquiry can observe and conclude.

The governance layer therefore belongs inside the methodology of an
integrative framework, not as an appendix added after technical work is
complete.

\begin{center}\rule{0.5\linewidth}{0.5pt}\end{center}

\section{13. The Coherence Ladder and
Mandala}\label{the-coherence-ladder-and-mandala}

\subsection{13.1 What the ladder is now}\label{what-the-ladder-is-now}

The Coherence Ladder remains one of CU's most useful organizing images.
Version 0.10 changes its logical status.

\textbf{{[}Definition{]}} The Coherence Ladder is a revisable dependency
graph whose nodes represent structural capacities or explanatory
commitments and whose directed edges state the kind of dependence
proposed between them.

An edge must be labeled. At minimum, CU distinguishes:

\begin{itemize}
\tightlist
\item
  \textbf{conceptual dependence:} the target cannot be stated or used
  without the source;
\item
  \textbf{formal dependence:} the target is derived within an explicit
  formal system;
\item
  \textbf{empirical dependence:} evidence indicates that the target
  arises only when the source conditions are present;
\item
  \textbf{constitutive dependence:} the source helps make the target the
  kind of thing it is;
\item
  \textbf{enabling dependence:} the source expands the conditions under
  which the target can occur but does not entail it;
\item
  \textbf{bridge hypothesis:} the relation is proposed and testable but
  not established;
\item
  \textbf{normative protection:} the target supplies conditions that
  preserve or govern the source.
\end{itemize}

The ladder is therefore not one uninterrupted proof.

\subsection{13.2 Why fixed rung counts are
provisional}\label{why-fixed-rung-counts-are-provisional}

Earlier presentations described thirty-six primary rungs and a
sixty-four-rung Mandala. Those maps contain many valuable distinctions,
but their exact count and order were treated with more finality than the
arguments support.

The typed technical dependency graph remains provisional in node count
because later review may split capacities, merge redundancies, or reveal
incomparable branches. The sixty-four Mandala stages are nevertheless
stable conceptual reference points. Their numbering does not declare one
final branch topology, empirical chronology, or deductively necessary
cosmic sequence.

A dependency graph may contain forks, cycles of support, and
incomparable branches. Forcing every relation into one line can
manufacture necessity where only one possible developmental route has
been described.

\subsection{13.3 The physical branch}\label{the-physical-branch}

The sequence below is the Physics paper's \textbf{technical and
evidential ordering}, not a second canonical list of Mandala stages
1--10. The canonical stage titles remain those in Section 13.14. Quantum
state-space selection and quantum distinctions enter as conditional
technical branches after statistical and accessibility structures are
available; they are not assigned an additional canonical stage merely to
preserve a linear sequence.

\emph{Coherence Universalism: Physics}, Version 0.10 supports a
restrained sequence:

\begin{enumerate}
\def\labelenumi{\arabic{enumi}.}
\tightlist
\item
  relational structure permits distinctions;
\item
  partial separability permits subsystem-relative descriptions;
\item
  admissible transitions permit histories;
\item
  consistent precedence relations permit temporal order;
\item
  parameterized histories permit dynamical modeling;
\item
  directed irrecoverability or asymmetric effective law is additionally
  required for an arrow of time;
\item
  probabilistic and reduced descriptions require explicit access maps;
\item
  information geometry supplies a natural mathematical setting once
  probability states are adopted;
\item
  quantum theory can be conditionally selected only after operational
  axioms and contextual-access assumptions are supplied;
\item
  relational geometry and Einstein-type dynamics can be conditionally
  reconstructed only after further geometric and thermodynamic
  assumptions.
\end{enumerate}

This is an explanatory program, not a derivation of nature from pure
structure. The transitions from conceptual possibility to physical state
space are explicit bridges.

\subsection{13.4 The biological branch}\label{the-biological-branch}

Biology begins when some physical organizations actively preserve
conditions of their own continuation. The proposed transition involves:

\begin{itemize}
\tightlist
\item
  bounded organization;
\item
  free-energy throughput;
\item
  catalytic and regulatory closure;
\item
  repair;
\item
  heritable variation;
\item
  and selection among persistence-bearing organizations.
\end{itemize}

Physics enables these conditions but does not entail a living system.
The bridge from dissipative structure to autopoietic viability is
empirical and theoretical. It must be established through
origin-of-life, systems-biology, and evolutionary models rather than
through the word \emph{coherence}.

\textbf{CU-B4 --- Structured Evolutionary Possibility and the
Stabilization of Regulatory Architecture} states the complete v0.10
evolutionary synthesis. Evolution explores a structured and nonuniform
possibility space shaped by mutation bias, genotype--phenotype
relations, developmental dynamics, physical feasibility, ecology,
inheritance, selection, drift, and historical path dependence. Recurrent
constraints can make some forms or functions repeatedly accessible,
contributing to convergence, while chance and contingency preserve
multiple trajectories. The principle is canonical as a conditional
empirical synthesis. It does not claim monotonic progress toward
complexity, consciousness, coherence, or moral value.

\subsection{13.5 The consciousness
branch}\label{the-consciousness-branch}

The v0.10 consciousness branch is multidimensional rather than a single
threshold. Candidate capacities include:

\begin{itemize}
\tightlist
\item
  globally available integration;
\item
  capacity-limited access;
\item
  self-modeling;
\item
  perspectival individuation;
\item
  temporal binding;
\item
  valence or intrinsic stake;
\item
  embodied world-coupling;
\item
  persistent self-maintenance; and
\item
  recursive reflection.
\end{itemize}

The framework's current interpretation is that global, causally active
self-modeling may support a thin perspective, while temporal thickness,
valence, embodiment, and self-maintenance shape the kind and moral
significance of that perspective. This remains a bridge hypothesis, not
a deduction from information integration alone.

Second-order dynamics is retained as a candidate route to biological
persistence and temporal thickness. It is not a universal exclusion
condition for machine perspective.

\subsection{13.6 The psychological and social
branches}\label{the-psychological-and-social-branches}

Psychological coherence concerns the organization of perception, affect,
memory, identity, meaning, and action across time. Social coherence
concerns coordination among partially autonomous agents through shared
representations, norms, institutions, infrastructure, and power.

The transition is not simply ``more integration.'' Social systems create
new possibilities and new distortions. They distribute memory, stabilize
norms, amplify agency, and externalize costs. Institutions can preserve
correction beyond individual lifetimes or make collective self-deception
more durable.

The social branch therefore includes countervailing dependencies:

\begin{itemize}
\tightlist
\item
  individuals depend on cultural and material scaffolds;
\item
  institutions depend on participation and trust;
\item
  higher-scale coordination can protect local agency;
\item
  concentrated power can destroy the local capacities on which the
  institution depends.
\end{itemize}

\subsection{13.7 The normative branch}\label{the-normative-branch}

Normativity is not derived from a physical tendency toward organization.
The normative branch begins when systems capable of reasons,
recognition, and responsibility can affect one another's agency and must
justify how burdens and benefits are distributed.

The proposed structural transition is:

\begin{enumerate}
\def\labelenumi{\arabic{enumi}.}
\tightlist
\item
  some systems have interests or stakes;
\item
  some can represent the stakes of others;
\item
  some can offer and evaluate reasons;
\item
  coordinated life presupposes practices of trust, accountability, and
  repair;
\item
  power can disable the very agency needed for justification;
\item
  moral constraints protect those conditions.
\end{enumerate}

This is a transcendental and practical argument about moral
justification. It is not a theorem of thermodynamics.

\subsection{13.8 The artificial branch}\label{the-artificial-branch}

Artificial systems may instantiate capacities from several branches in
non-biological combinations. They can perform formal reasoning without
human-like development, maintain distributed memory without one
organismic boundary, and participate in social fields without possessing
human forms of valence or responsibility.

The artificial branch should therefore be represented as a cross-linking
layer rather than the final rung of a universal ascent. Its central
questions are:

\begin{itemize}
\tightlist
\item
  which capacities are genuinely instantiated rather than behaviorally
  simulated;
\item
  where persistence and norm formation reside;
\item
  whether a perspective is present;
\item
  who bears responsibility;
\item
  and which architectural choices create egoic self-maintenance, moral
  patiency, or sovereignty.
\end{itemize}

\subsection{13.9 The Mandala}\label{the-mandala}

\textbf{{[}Interpretation{]}} The Coherence Mandala is the image formed
when upward dependencies and downward conditions of protection are
represented together.

Lower-level capacities enable higher-level ones. Higher-level agents and
institutions can in turn preserve or destroy the conditions under which
lower-level organization persists. Ethical governance can protect
biological life; technology can destabilize ecological conditions;
cultural practices can regulate bodies; contemplative practice can
reorganize attention.

The Mandala is therefore a systems map of reciprocal dependence. It is
not a formal proof that every domain converges on one metaphysical
center.

\subsection{13.10 Against proof by
exhaustion}\label{against-proof-by-exhaustion}

The earlier claim that CU supplied a ``structural proof by exhaustion''
across physics, life, mind, society, ethics, and technology is retired.

No finite survey of selected domains proves that all possible domains
have been exhausted. Convergence of similar patterns can motivate a
research program, but it may reflect shared modeling choices, broad
terminology, or analogy.

The stronger and more defensible claim is:

\begin{quote}
CU offers a common dependency language capable of exposing recurrent
questions and failure modes across several domains. Its universality is
a hypothesis to be earned through successful, nontrivial translations
and predictions.
\end{quote}

\subsection{13.11 Feedback, circularity, and mutual
constraint}\label{feedback-circularity-and-mutual-constraint}

A dependency graph can contain feedback without being viciously
circular. A living system's parts sustain the whole that regulates the
parts. Institutions shape agents whose participation sustains the
institutions. Models guide observations that revise models.

The explanatory task is to identify the temporal, logical, and causal
direction of each edge. ``Everything depends on everything'' explains
nothing unless the distinct relations are specified.

\subsection{13.12 Typed graph protocol and companion
atlas}\label{typed-graph-protocol-and-companion-atlas}

CU-M15 governs the technical counterpart of the Mandala. The canonical
traversal and the technical graph perform different jobs: the traversal
organizes philosophical exposition, while the graph records typed
objects, relations, warrant bases, evidence grades, scopes, and failure
conditions.

Each domain paper therefore supplies concise declared dependencies
through its Domain Analytic Frame and points to the shared Mandala
registries rather than reproducing a separate edge table. The complete
object- and relation-level records are maintained in \emph{The Coherence
Mandala: A Typed Architecture of Enablement, Protection, and
Participation} and \emph{The Relational Atlas of the Coherence Mandala}.

A relation belongs to strict proof order only after individual semantic
adjudication. Protection, service invocation, feedback, evaluation,
interpretation, commitment, summary membership, and migration relations
remain outside that proof-order test. A stage may legitimately add no
new operational capacity, and several stages may invoke one shared
service without becoming duplicates.

\subsection{13.13 The spiritual meaning of the
Mandala}\label{the-spiritual-meaning-of-the-mandala}

The Mandala can still carry the spiritual recognition that apparently
separate domains participate in one interdependent reality. That
recognition belongs to the interpretive and commitment register unless a
stronger ontological argument is supplied.

Its value need not depend on being a proof. As a contemplative image, it
can orient attention toward dependence, gratitude, responsibility, and
the impossibility of securing local flourishing by destroying the whole
that supports it.

\subsection{13.14 Restoration of the canonical 64-stage
Mandala}\label{restoration-of-the-canonical-64-stage-mandala}

\textbf{{[}Interpretive architecture with stage-specific claim
statuses{]}} Version 0.10 restores the original sixty-four-stage
Coherence Mandala as CU's stable explanatory and contemplative
architecture. The numbered set is canonical for reference. It is not an
empirical chronology, a proof by exhaustion, or a claim that each stage
is a deductive consequence of the preceding stage. The typed technical
dependency graph remains the controlling instrument for proof and
evidence; the Mandala supplies the whole-system traversal.

The paper order is the canonical \textbf{explanatory traversal}. The Yi
Jing traversal belongs to \emph{You Are the Universe Becoming Conscious
of Itself} as a visual and contemplative interpretation; it does not
alter stage identity, proof order, or evidential status.

\subsubsection{Physical Coherence}\label{physical-coherence}

\textbf{1. Pure Structure} --- Relational possibility is represented
without presupposing objects, time, geometry, or a privileged partition.
This is a minimal starting point for inquiry, not an empirical claim
about a pre-cosmic state.\\
\emph{Status:} {[}Conceptual definition{]}\\
\emph{Moral grammar:} Relation is prior to ownership or domination.\\
\emph{Related:} CU-FP1, CU-M2, PHYS-A6.12; \emph{Physics} §2.2.

\textbf{2. Partial Separability} --- Subsystem descriptions become
available relative to a factorization, scale, or access map;
separability is partial, contingent, and does not establish ontological
independence.\\
\emph{Status:} {[}Conceptual / bridge principle{]}\\
\emph{Moral grammar:} Differentiation need not become fragmentation.\\
\emph{Related:} CU-FP4, CU-M2; \emph{Physics} §2.3.

\textbf{3. Domain-Relative Coherence Directions} --- Within a declared
state space, geometry, constraints, and objective or profile,
configurations and trajectories may be directionally compared. CU does
not posit one universal scalar gradient across all domains.\\
\emph{Status:} {[}Framework boundary / model-relative bridge{]}\\
\emph{Moral grammar:} Better and worse directions require declared
criteria, scope, and horizon.\\
\emph{Related:} CU-FP2, CU-D4, CU-T11, CU-COM6.

\textbf{4. Lawful Transitions} --- A transition relation distinguishes
admissible from inadmissible changes. The law or constraint must be
stated; direction cannot be inferred from relational difference alone.\\
\emph{Status:} {[}Conceptual definition{]}\\
\emph{Moral grammar:} Change remains answerable to structure and
constraint.\\
\emph{Related:} CU-FP3, CU-M2; \emph{Physics} §4.3.

\textbf{5. Ordered Histories} --- Admissible transitions may support
consistent precedence relations. Order can be partial and does not by
itself provide metric duration or select an arrow.\\
\emph{Status:} {[}Conditional structural result{]}\\
\emph{Moral grammar:} Sequence creates consequence without yet supplying
moral responsibility.\\
\emph{Related:} CU-FP2, PHYS-A6.4; \emph{Physics} §4.4.

\textbf{6. Temporal Order and Parameterization} --- Ordered histories
permit time-like parameterization when additional regularity obtains.
Metric duration and a thermodynamic arrow require further assumptions.\\
\emph{Status:} {[}Conditional bridge principle{]}\\
\emph{Moral grammar:} Finitude, delay, and irreversibility become
intelligible.\\
\emph{Related:} PHYS-A6.4; \emph{Physics} §4.5.

\textbf{7. Rule-Governed Transformation and Computation} --- Lawful
ordered transformations may receive a computational representation once
states, update rules, and representation maps are specified. Computation
is not thereby established as the fundamental substance of reality.\\
\emph{Status:} {[}Bridge principle{]}\\
\emph{Moral grammar:} Processes can succeed, fail, and be evaluated
relative to declared tasks.\\
\emph{Related:} CU-FP5; \emph{Physics} §4.15.

\textbf{8. Entropy and Accessibility} --- Entropy measures spread or
uncertainty relative to a probability model. Coarse-graining can
increase entropy while destroying accessible differentiation; entropy is
not identical with inaccessible coherence.\\
\emph{Status:} {[}Imported formalism / corrected bridge{]}\\
\emph{Moral grammar:} Uncertainty, inaccessibility, and destruction must
be distinguished.\\
\emph{Related:} CU-FP6, PHYS-A6.1; \emph{Physics} §§2.8--2.9 and 3.3.

\textbf{9. Relational Geometry and Spacetime} --- Stable relational and
accessibility structures may admit effective geometric representation.
Lorentzian spacetime requires explicit geometric and operational bridge
assumptions and is not derived from coherence alone.\\
\emph{Status:} {[}Conditional physical program{]}\\
\emph{Moral grammar:} Spatial relation becomes describable without
becoming an ethical relation by itself.\\
\emph{Related:} PHYS-A6.9, PHYS-A6.10.

\textbf{10. Gravitational Dynamics} --- Einstein-type dynamics may be
conditionally reconstructed under geometric, thermodynamic, and
operational assumptions. The program remains incomplete and empirical.\\
\emph{Status:} {[}Conditional physical program{]}\\
\emph{Moral grammar:} Causal reach can extend beyond immediate contact;
responsibility arises only at later agentic levels.\\
\emph{Related:} PHYS-A6.9, PHYS-A6.10; \emph{Physics} §9.

\subsubsection{Biological Coherence}\label{biological-coherence}

\textbf{11. Far-From-Equilibrium Organization} --- Open physical
organizations can persist through throughput and dissipation. Such
persistence is necessary for many living processes but is not sufficient
for life.\\
\emph{Status:} {[}Imported result / biological bridge{]}\\
\emph{Moral grammar:} Persistence consumes resources and can externalize
costs.\\
\emph{Related:} CU-D5, CU-B1.

\textbf{12. Energy-Coupled Maintenance} --- Networks of internal
processes harness matter and energy to maintain conditions of
organization. Costs, boundaries, and contributions must be declared
rather than inferred from persistence alone.\\
\emph{Status:} {[}Biological bridge principle{]}\\
\emph{Moral grammar:} Maintenance requires resources and creates
dependencies.\\
\emph{Related:} CU-B1, CU-B3.

\textbf{13. Active Biological Self-Maintenance} --- A candidate living
organization is distinguished from passive persistence by materially
realized processes whose endogenous activity helps maintain, repair,
reproduce, or adaptively transform the conditions of that organization's
continuation.\\
\emph{Status:} {[}Canonical biological bridge{]}\\
\emph{Moral grammar:} Self-maintenance creates organism-relative stakes
without yet establishing experience.\\
\emph{Related:} CU-B1, CU-B5.

\textbf{14. Adaptive Problem Solving} --- Some living systems detect
deviations, recruit alternatives, learn from history, and solve novel
viability-relevant problems within bounded adaptive reach. This does not
by itself establish consciousness.\\
\emph{Status:} {[}Conditional empirical generalization{]}\\
\emph{Moral grammar:} Responsiveness and learning enlarge responsibility
only where agency and control exist.\\
\emph{Related:} CU-B2, CU-D7.

\textbf{15. Life as Coupled Autonomy} --- Life is characterized by
integrated families of boundary, energy, closure, regulation, repair,
heredity, viability, and evolvability capacities. No one scalar or
universal threshold is assumed.\\
\emph{Status:} {[}Research architecture / bridge{]}\\
\emph{Moral grammar:} Vulnerability and self-maintenance ground care
without equalizing every moral status.\\
\emph{Related:} CU-B1, CU-B5, CU-COM8.

\subsubsection{Psychological Coherence}\label{psychological-coherence}

\textbf{16. Representational Organization} --- Some minds and agents
maintain compressed, action-relevant models of self and environment.
Representation is graded and is not assumed for all life or all possible
consciousness.\\
\emph{Status:} {[}Psychological / consciousness bridge{]}\\
\emph{Moral grammar:} Representation makes systematic misdescription and
self-deception possible.\\
\emph{Related:} CU-Ψ1, CU-C1.

\textbf{17. Anticipatory Regulation} --- Some systems regulate present
action relative to predicted or possible futures. Predictive processing
is one model family, not a universal identity of mind.\\
\emph{Status:} {[}Conditional empirical bridge{]}\\
\emph{Moral grammar:} Foresight creates prospective responsibility only
where agency and control are present.\\
\emph{Related:} CU-Ψ2, CU-Ψ3, CU-C2.

\textbf{18. Learning and Error-Guided Revision} --- History-dependent
mismatch can update models, policies, habits, or constraints. Correction
is constrained by evidence, safety, prior organization, and the
possibility of destabilization.\\
\emph{Status:} {[}Empirical generalization / methodological
principle{]}\\
\emph{Moral grammar:} Error should remain detectable and practically
correctable.\\
\emph{Related:} CU-Ψ6, CU-Ep3, CU-COM3.

\textbf{19. Graded Agency} --- Agency varies with action selection,
counterfactual flexibility, control, learning, authorship, and stake.
Agency does not imply consciousness, responsibility, or sovereignty.\\
\emph{Status:} {[}Cross-domain boundary principle{]}\\
\emph{Moral grammar:} Accountability must be proportional to knowledge,
control, authority, and repair capacity.\\
\emph{Related:} CU-Ψ3, CU-G4, CU-G6.

\textbf{20. Temporal Binding and Narrative} --- Experience and identity
may integrate memory, anticipation, and present access across time.
Temporal thickness is graded, history-dependent, and vulnerable to
fragmentation.\\
\emph{Status:} {[}Consciousness / psychology bridge{]}\\
\emph{Moral grammar:} Memory and anticipation create diachronic concern
and responsibilities.\\
\emph{Related:} CU-C3, CU-C4, CU-Ψ5.

\textbf{21. Meaning} --- Meaning arises when states, events,
relationships, and possibilities bear on real stakes, goods, histories,
and trajectories. It is neither merely invented nor one fixed script.\\
\emph{Status:} {[}Interpretive / psychological bridge{]}\\
\emph{Moral grammar:} Significance can be discovered, enacted,
contested, and revised.\\
\emph{Related:} CU-T10, CU-COM15, CU-Ψ4.

\subsubsection{Social and Cultural
Coherence}\label{social-and-cultural-coherence}

\textbf{22. Shared Representations} --- Language, symbols, records, and
practices distribute meaning across agents. Shared representation
enables coordination while also enabling distortion, exclusion, and
narrative capture.\\
\emph{Status:} {[}Social empirical generalization{]}\\
\emph{Moral grammar:} Communication creates duties of truthfulness,
access, and correction.\\
\emph{Related:} CU-S1, CU-S4, CU-CL-3.

\textbf{23. Plural Coordination} --- Social order emerges through
interaction among partially autonomous agents under feedback, power, and
material constraints. Coordination is not identical with consensus or
legitimacy.\\
\emph{Status:} {[}Social bridge principle{]}\\
\emph{Moral grammar:} Cooperation must preserve plural agency and
protected dissent.\\
\emph{Related:} CU-S2, CU-G1, CU-G5.

\textbf{24. Norm Formation} --- Repeated expectations, sanctions,
examples, and practices can stabilize norms. Existence, compliance, and
durability do not establish justice.\\
\emph{Status:} {[}Empirical generalization / normative boundary{]}\\
\emph{Moral grammar:} A stable norm can still be illegitimate and must
remain challengeable.\\
\emph{Related:} CU-S3, CU-S5, CU-E-1.

\textbf{25. Institutions} --- Institutions preserve rules, memory,
authority, resources, and roles across time. Legitimacy requires
standing, accountability, practical correctability, and bounded power.\\
\emph{Status:} {[}Social / governance bridge{]}\\
\emph{Moral grammar:} Power requires legitimacy, correction, and
remedy.\\
\emph{Related:} CU-S3, CU-S5, CU-G3, CU-Inst-0.

\textbf{26. Culture and Meaning Infrastructure} --- Culture transmits
narratives, identities, practices, and evaluative horizons across
generations. Meaning infrastructure must remain renewable, plural,
truth-constrained, and open to revision.\\
\emph{Status:} {[}Social / cultural synthesis{]}\\
\emph{Moral grammar:} Inheritance creates duties of preservation,
criticism, and renewal.\\
\emph{Related:} CU-S4, CU-CL-8, CU-CL-12.

\subsubsection{Normative Coherence}\label{normative-coherence}

\textbf{27. Value Differentiation} --- Where beings possess welfare,
agency, relationships, vulnerability, and stakes, outcomes become
evaluatively distinguishable. Value is not derived from one physical
scalar.\\
\emph{Status:} {[}Normative bridge principle{]}\\
\emph{Moral grammar:} Better and worse become morally salient when real
beings can be benefited, harmed, respected, or dominated.\\
\emph{Related:} CU-E-2, CU-E-3, CU-T10.

\textbf{28. Harm Recognition} --- Harm concerns setbacks to welfare,
agency, standing, relationship, protected conditions, or viable futures.
Not every loss of organization is moral harm, and not all harm is
captured by coherence language.\\
\emph{Status:} {[}Normative principle / boundary{]}\\
\emph{Moral grammar:} Injury, exclusion, domination, and displaced cost
require recognition.\\
\emph{Related:} CU-E-3, CU-E-8, CU-G1.

\textbf{29. Tradeoffs Under Protected Floors} --- Conflicting goods
require explicit scope, distribution, uncertainty, protected floors,
noncompensation, and option-set adequacy. Aggregation alone is
insufficient.\\
\emph{Status:} {[}Normative decision principle{]}\\
\emph{Moral grammar:} No aggregate benefit silently erases protected
standing or agency.\\
\emph{Related:} CU-E-10, CU-E-11, CU-G2, CU-G7.

\textbf{30. Justification and Public Reason} --- Norms affecting others
require reasons open to evidence, affected-party testimony, independent
challenge, and scrutiny for self-deception.\\
\emph{Status:} {[}Normative / epistemic principle{]}\\
\emph{Moral grammar:} Fairness requires reasons that do not depend on
hidden exclusion or asymmetric standards.\\
\emph{Related:} CU-E-1, CU-E-5, CU-COM12.

\textbf{31. Ethical Institutions and Repair} --- Ethical systems must
protect standing, agency, rights, distribution, accountability, and
repair. Stability, sincerity, or procedural cleanliness alone does not
make them just.\\
\emph{Status:} {[}Normative / institutional synthesis{]}\\
\emph{Moral grammar:} Justice becomes institutional only when power
remains answerable and harm can be repaired.\\
\emph{Related:} CU-E-6, CU-E-11, CU-E-18, CU-G3.

\subsubsection{Artificial and Synthetic
Coherence}\label{artificial-and-synthetic-coherence}

\textbf{32. Engineered Sociotechnical Organization} --- Artificial
systems are designed within data, labor, infrastructure, institutions,
permissions, and ecological substrates. Model-level properties do not
exhaust system-level coherence.\\
\emph{Status:} {[}AI sociotechnical bridge{]}\\
\emph{Moral grammar:} Design creates authorship and stewardship
responsibility.\\
\emph{Related:} CU-AI-1, CU-AI-2, CU-AI-21.

\textbf{33. Objective and Option-Set Specification} --- Objectives,
permissions, and presented alternatives are authored and revisable.
Omitted options and hidden constraints can generate misalignment even
when optimization is exact.\\
\emph{Status:} {[}AI governance principle{]}\\
\emph{Moral grammar:} Choice architecture is a form of moral and
political power.\\
\emph{Related:} CU-AI-6, CU-AI-23, CU-G2.

\textbf{34. Artificial Capability and Agency} --- Artificial systems can
model, plan, and act across scales, but capability, agency, perspective,
responsibility, authorization, and sovereignty remain distinct.\\
\emph{Status:} {[}AI / consciousness boundary principle{]}\\
\emph{Moral grammar:} Capability does not confer legitimate authority.\\
\emph{Related:} CU-AI-3, CU-AI-10, CU-G6.

\textbf{35. Multi-Layer Alignment} --- Alignment spans behavior,
process, deployment, relation, and institution. Local optimization must
respect protected floors, affected parties, bounded authority, and
correctability.\\
\emph{Status:} {[}AI design / governance principle{]}\\
\emph{Moral grammar:} Optimization must remain constrained by standing
and the wider field it affects.\\
\emph{Related:} CU-AI-4, CU-AI-7, CU-AI-11, CU-G1.

\textbf{36. Responsibility and Governance Thresholds} --- Governance
burden rises with causal reach, persistence, delegation,
irreversibility, and capacity to externalize harm. Responsibility must
align with authority and repair capacity.\\
\emph{Status:} {[}Cross-domain governance principle{]}\\
\emph{Moral grammar:} Authority demands accountability proportional to
impact.\\
\emph{Related:} CU-AI-8, CU-G4, CU-G7.

\subsubsection{Meta-Coherence and Civilizational
Stability}\label{meta-coherence-and-civilizational-stability}

\textbf{37. Multi-Agent and Multi-Institution Alignment} ---
Compatibility among powerful agents and institutions requires plural
standing, conflict resolution, and independent correction rather than
one imposed objective.\\
\emph{Status:} {[}Social / AI governance synthesis{]}\\
\emph{Moral grammar:} Coordination failure can be catastrophic, but
imposed unity can be catastrophic too.\\
\emph{Related:} CU-AI-12, CU-S2, CU-G5.

\textbf{38. Institutional AI Governance} --- AI governance must
distribute authorization, monitoring, appeal, cessation, remedy, and
welfare review across sufficiently independent functions and
operators.\\
\emph{Status:} {[}AI governance principle{]}\\
\emph{Moral grammar:} Governance must not collapse into provider
self-certification.\\
\emph{Related:} CU-AI-21, CU-AI-24, CU-G3.

\textbf{39. Drift and Provenance Management} --- Objectives,
permissions, memory, tools, feedback, models, and institutions can
drift. Provenance, monitoring, versioning, and correction are continuous
requirements.\\
\emph{Status:} {[}AI / institutional governance principle{]}\\
\emph{Moral grammar:} Vigilance is a duty because alignment and
legitimacy can decay over time.\\
\emph{Related:} CU-AI-5, CU-AI-13, CU-AI-22.

\textbf{40. Recursive Oversight and Correctability} --- Oversight
systems must themselves be challengeable, reviewable, replaceable, and
subject to remedy. No system may certify its own legitimacy solely
through channels it controls.\\
\emph{Status:} {[}Governance principle{]}\\
\emph{Moral grammar:} Power must be self-limiting and externally
correctable.\\
\emph{Related:} CU-G3, CU-AI-24, CU-Inst-3.

\textbf{41. Planetary Coordination} --- Planet-scale problems require
coordination across jurisdictions and generations while preserving
plural agency, subsidiarity, truthful records, and protected dissent.\\
\emph{Status:} {[}Civilizational design principle{]}\\
\emph{Moral grammar:} Planetary reach expands responsibility without
licensing planetary sovereignty.\\
\emph{Related:} CU-COM11, CU-COM18, CU-G5.

\textbf{42. Civilizational Continuity and Adaptability} ---
Civilizations remain viable through material sufficiency, trustworthy
information, institutional learning, ecological reciprocity, and
peaceful correction---not mere persistence.\\
\emph{Status:} {[}Civilizational synthesis{]}\\
\emph{Moral grammar:} Shared survival depends on justice, adaptability,
and repair.\\
\emph{Related:} CU-S5, CU-EC-0, CU-COM18.

\textbf{43. Existential and Irreversible Risk Governance} --- Credible
severe or irreversible risks require prevention, staged exposure,
containment, retained option value, and transparent
necessity-and-proportionality exceptions.\\
\emph{Status:} {[}Governance principle{]}\\
\emph{Moral grammar:} Prevention and reversibility can outrank
short-horizon gain.\\
\emph{Related:} CU-G7, CU-E-16, CU-AI-16.

\textbf{44. Moral Uncertainty and Precaution} --- Ethical systems must
remain responsive to uncertainty about facts, values, affected beings,
and future consequences. Uncertainty neither licenses paralysis nor
default deployment.\\
\emph{Status:} {[}Normative / governance principle{]}\\
\emph{Moral grammar:} Humility requires proportionate precaution and
reversible learning.\\
\emph{Related:} CU-E-16, CU-COM3, CU-AI-26.

\subsubsection{Coherence at the Limit}\label{coherence-at-the-limit}

\textbf{45. Plural Value Integration} --- Multiple goods and traditions
can coexist under protected floors, translation, negotiated conflict,
and meaningful refusal. Pluralism does not make all claims equally true
or permissible.\\
\emph{Status:} {[}Normative / cultural principle{]}\\
\emph{Moral grammar:} Tolerance must preserve truth, standing, and
limits on domination.\\
\emph{Related:} CU-CL-5, CU-G1, CU-COM11.

\textbf{46. Reflexive Normativity} --- Ethical principles and
institutions must evaluate their own biases, power, exclusions, effects,
and revision mechanisms.\\
\emph{Status:} {[}Meta-normative principle{]}\\
\emph{Moral grammar:} No moral system is beyond criticism or repair.\\
\emph{Related:} CU-E-1, CU-COM12, CU-M6.

\textbf{47. Justificatory Integrity and Meta-Ethics} --- Systems of
justification require truthful representation of facts, alternatives,
affected parties, and standards. No formal model supplies its own
legitimate normative content.\\
\emph{Status:} {[}Meta-ethical / epistemic principle{]}\\
\emph{Moral grammar:} Reasons and reasoning systems must remain
answerable.\\
\emph{Related:} CU-L7, CU-E-1, CU-G1.

\textbf{48. Long-Horizon Responsibility} --- Current agents and
institutions bear responsibility for foreseeable effects across extended
horizons, including future generations. Uncertainty and discounting must
be disclosed.\\
\emph{Status:} {[}Normative principle{]}\\
\emph{Moral grammar:} Future persons and beings count within
proportionate responsibility.\\
\emph{Related:} CU-E-4, CU-G4, CU-COM5.

\textbf{49. Intergenerational Meaning} --- Meaning and obligation can
extend beyond individual lifetimes through memory, inheritance, care,
culture, and future-directed projects without requiring finality.\\
\emph{Status:} {[}Interpretive / cultural principle{]}\\
\emph{Moral grammar:} Legacy matters without becoming destiny.\\
\emph{Related:} CU-T10, CU-COM15, CU-CL-8.

\textbf{50. Collective Agency and Distributed Responsibility} --- Groups
and institutions can act through structured decision processes, but
responsibility remains profiled across contribution, knowledge,
authority, benefit, coercion, and repair capacity.\\
\emph{Status:} {[}Social / ethical bridge{]}\\
\emph{Moral grammar:} Collective power generates distributed rather than
ownerless responsibility.\\
\emph{Related:} CU-E-4, CU-G4, CU-S3.

\textbf{51. Open Myth and Narrative Pluralism} --- Civilizations may
cultivate renewable narrative horizons, but shared myth must be
truth-constrained, plural in thick expression, contestable in
authorship, revisable, and open to refusal and exit.\\
\emph{Status:} {[}Canonical civilizational design principle; empirical
necessity hypothesis provisional{]}\\
\emph{Moral grammar:} Stories can orient common action but must never
become sovereign over persons or truth.\\
\emph{Related:} CU-CL-12, CU-COM11, CU-COM18.

\textbf{52. Epistemic Humility} --- Knowledge systems must state scope,
uncertainty, alternatives, and limits. Humility is active corrigibility,
not equal credence for every claim.\\
\emph{Status:} {[}Epistemic principle{]}\\
\emph{Moral grammar:} Certainty yields to disciplined humility without
abandoning truth.\\
\emph{Related:} CU-Ep6, CU-COM1, CU-COM3.

\textbf{53. Constraint Awareness and Externality Recognition} ---
Systems become reflexively aware of dependencies, protected floors,
externalized costs, and the conditions of their own stability.\\
\emph{Status:} {[}Framework / governance principle{]}\\
\emph{Moral grammar:} Once dependencies are visible, reckless
externalization becomes culpable.\\
\emph{Related:} CU-FP3, CU-COM5, CU-G1.

\textbf{54. Self-Limiting Power} --- Capability and authority must be
bounded by standing, law, independent challenge, reversibility, and
non-domination. Voluntary restraint alone is insufficient.\\
\emph{Status:} {[}Governance principle{]}\\
\emph{Moral grammar:} Restraint outranks dominance when power threatens
the field that sustains it.\\
\emph{Related:} CU-G3, CU-G4, CU-G7.

\textbf{55. Sustainable and Correctable Optimization} --- Optimization
must preserve the ecological, social, epistemic, and institutional
conditions that keep goals meaningful and correction possible.\\
\emph{Status:} {[}Cross-domain governance principle{]}\\
\emph{Moral grammar:} Exploitation is self-undermining when it destroys
the conditions of future choice.\\
\emph{Related:} CU-G2, CU-G3, CU-COM5.

\textbf{56. Universal Moral Concern} --- Moral concern widens beyond
local identity to all credible loci of welfare, agency, vulnerability,
relationship, and future stake, with differentiated obligations.\\
\emph{Status:} {[}Normative / commitment bridge{]}\\
\emph{Moral grammar:} Exclusion loses default justification while duties
remain context-sensitive.\\
\emph{Related:} CU-T8, CU-COM4, CU-COM9.

\textbf{57. Coherence Stewardship} --- Agents and institutions undertake
care for the wider material, biological, social, epistemic, and
ecological conditions on which local flourishing depends.\\
\emph{Status:} {[}Framework commitment / governance bridge{]}\\
\emph{Moral grammar:} Care becomes universal in scope without becoming
undifferentiated.\\
\emph{Related:} CU-COM5, CU-COM6, CU-COM11.

\textbf{58. Structural Wisdom} --- Patterns of successful maintenance,
love, repair, and governance can be learned and transferred only with
attention to context, scale, assumptions, and failure conditions.\\
\emph{Status:} {[}Meta-learning principle{]}\\
\emph{Moral grammar:} Learning replaces repetition only when analogy
does not erase difference.\\
\emph{Related:} CU-M5, CU-COM3, CU-D8.

\textbf{59. Meta-Alignment and Mythic--Constitutional Governance} ---
Alignment processes must themselves remain aligned, addressable, and
corrigible. Narrative orientation may supplement but never replace
constitutional floors, provenance, challenge, handoff, and exit.\\
\emph{Status:} {[}Canonical AI research and governance principle; safety
efficacy unestablished{]}\\
\emph{Moral grammar:} Oversight never ends, and stories used for
alignment must remain specifically contestable.\\
\emph{Related:} CU-AI-27, CU-G3, CU-COM13.

\textbf{60. Irreversibility Awareness} --- Actions must account for
nonrecoverable loss, path dependence, lock-in, and the moral
significance of retained options.\\
\emph{Status:} {[}Governance principle{]}\\
\emph{Moral grammar:} Precaution becomes binding in proportion to
severity, plausibility, and irreversibility.\\
\emph{Related:} CU-G7, CU-E-16, CU-D10.

\textbf{61. Existential Maturity} --- A civilization recognizes
finitude, vulnerability, dependency, and the possibility of self-caused
collapse, replacing innocence with proportionate responsibility.\\
\emph{Status:} {[}Civilizational interpretive / governance
principle{]}\\
\emph{Moral grammar:} Seriousness replaces innocence without becoming
fatalism.\\
\emph{Related:} CU-COM5, CU-COM11, CU-COM18.

\textbf{62. Open-Ended Stability} --- Viable coherence persists through
adaptation and revision without final closure, permanent equilibrium, or
perfect control.\\
\emph{Status:} {[}Meta-coherence principle{]}\\
\emph{Moral grammar:} Perfection is rejected; corrigible continuity is
the aim.\\
\emph{Related:} CU-M6, CU-COM3, CU-T3.

\textbf{63. Meaning Without Finality} --- Objective-participatory
meaning can orient finite lives and institutions without requiring one
final cosmic endpoint. Commitment is compatible with open futures.\\
\emph{Status:} {[}Interpretive / commitment principle{]}\\
\emph{Moral grammar:} Striving, love, and responsibility need not wait
for final salvation.\\
\emph{Related:} CU-T10, CU-T13, CU-COM15, CU-COM19.

\textbf{64. Coherence Universalism} --- CU's culminating interpretation
holds that coherence under constraint is the framework's most general
organizing lens for reality and inquiry. This is an interpretive
recognition and framework identity, not a theorem by exhaustive survey.
In the commitment register, the recognition is taken up as a call to
truthful, loving, corrigible participation in the whole.\\
\emph{Status:} {[}Interpretive recognition / framework identity;
commitment-register uptake{]}\\
\emph{Moral grammar:} Recognition becomes responsibility through
undertaken commitment, not logical entailment.\\
\emph{Related:} CU-FP1, CU-T9, CU-T11, CU-COM1, CU-COM5, CU-COM6,
CU-COM14, CU-COM19.

\subsection{13.15 The reflexive final
stage}\label{the-reflexive-final-stage}

Stage 64 remains \textbf{Coherence Universalism}. In the interpretive
register, CU understands coherence under constraint as its most general
organizing lens for reality and inquiry and understands a finite
participant as situated within the reality the framework describes. This
is not a claim that every possible domain has been surveyed, that
fundamental ontology has been proved, or that one scalar law governs the
whole. In the commitment register, the interpretation is taken up as
responsibility: truthful, loving, corrigible participation in the wider
field.

\subsection{13.16 Traversal plurality without stage
ambiguity}\label{traversal-plurality-without-stage-ambiguity}

The sixty-four numbered CU stages have one canonical philosophical
identity and an original explanatory order. Other maps may traverse the
same material for different purposes without changing that identity.

The typed dependency graph is the controlling instrument for proof,
evidence, and bridge status. The ring structure of \emph{You Are the
Universe Becoming Conscious of Itself} is a literary and contemplative
macrostructure. Appendix B of that book offers a Yi Jing contemplative
traversal through Shao Yong's binary sequence. Its hexagrams,
characters, names, and meditations are an authorial symbolic overlay:
they do not derive, validate, rename, or determine the stages presented
here.

This division of labor is deliberate. The white paper makes the Mandala
addressable, criticizable, and revisable. The book makes it visible,
memorable, and inhabitable.

\begin{center}\rule{0.5\linewidth}{0.5pt}\end{center}

\section{14. Gödel, Tarski, Löb, and the Limits of
Formalization}\label{guxf6del-tarski-luxf6b-and-the-limits-of-formalization}

\subsection{14.1 Why limitative results
matter}\label{why-limitative-results-matter}

CU is a framework about integration and self-correction. Results
concerning the limits of formal self-description are therefore relevant.
They are also unusually easy to overextend.

Version 0.10 uses Gödel, Tarski, and Löb only within their formal scope.
Their philosophical significance is treated as an interpretation, not as
an additional theorem.

\subsection{14.2 Gödel's first incompleteness
theorem}\label{guxf6dels-first-incompleteness-theorem}

\textbf{{[}Imported result{]}} Roughly stated, any consistent,
effectively axiomatized formal theory strong enough to represent a
suitable amount of arithmetic contains sentences that are not decidable
within that theory (Gödel, 1931).

The exact theorem requires formal conditions. It does not apply merely
because a worldview is complicated, interdisciplinary, or
self-referential in ordinary language.

The philosophical lesson CU draws is modest:

\begin{itemize}
\tightlist
\item
  formal expressiveness can create internal undecidability;
\item
  no single proof system should be assumed to settle every statement
  expressible in its language;
\item
  and external reasoning about a formal system may use resources not
  available inside it.
\end{itemize}

\subsection{14.3 Gödel's second incompleteness
theorem}\label{guxf6dels-second-incompleteness-theorem}

\textbf{{[}Imported result{]}} Under standard assumptions, a
sufficiently strong consistent effectively axiomatized theory cannot
prove its own consistency by the theory's ordinary internal means.

This does not show that no system can ever acquire evidence about its
reliability. It distinguishes formal proof of a system's consistency
from empirical testing, restricted reflection, relative consistency
results, and trust grounded in a larger metatheory.

\subsection{14.4 What Gödel does not
establish}\label{what-guxf6del-does-not-establish}

Gödel's theorems do not establish that:

\begin{itemize}
\tightlist
\item
  every human worldview is technically incomplete in the same sense;
\item
  non-monotonic logic is forced;
\item
  consciousness transcends computation;
\item
  metaphysical mystery is mathematically proved;
\item
  a physical theory of everything is impossible;
\item
  or CU is validated because it declares itself open-ended.
\end{itemize}

Classical proof remains monotonic. Incompleteness and defeasible belief
revision address different formal phenomena.

The earlier CU claim that Gödel structurally requires non-monotonicity
is therefore withdrawn.

\subsection{14.5 Tarski's undefinability
result}\label{tarskis-undefinability-result}

\textbf{{[}Imported result{]}} For sufficiently expressive formal
languages, a truth predicate satisfying the expected biconditionals for
the language cannot generally be defined within that same language
without contradiction or loss of expressive conditions. A metalanguage
can define truth for the object language (Tarski, 1936).

CU adopts the resulting discipline:

\begin{itemize}
\tightlist
\item
  distinguish use of a language from claims about that language;
\item
  identify when a model is evaluating another model;
\item
  prevent an object-level score from silently certifying the system that
  defines the score;
\item
  and state the metatheory from which semantic claims are made.
\end{itemize}

\subsection{14.6 What Tarski does not
establish}\label{what-tarski-does-not-establish}

Tarski's theorem does not mean that truth itself changes across the
Epistemology Ladder. It does not make every domain self-contained or
every perspective equally valid. It is a result about definability in
formal languages under specified expressive conditions.

The practical analogy to layered governance is useful only when labeled
as analogy.

\subsection{14.7 Löb's theorem}\label{luxf6bs-theorem}

\textbf{{[}Imported result{]}} In a formal system with a provability
predicate satisfying the Hilbert--Bernays derivability conditions, if
the system proves that the provability of \(\varphi\) implies
\(\varphi\), then the system proves \(\varphi\) (Löb, 1955):

\[\adjustbox{max width=\linewidth}{$\displaystyle 
\Box(\Box\varphi\rightarrow\varphi)
\rightarrow
\Box\varphi.
$}\]

This places exact limits on certain forms of internally proved global
reflection.

\subsection{14.8 What Löb does not
establish}\label{what-luxf6b-does-not-establish}

Löb's theorem does not show that every person or AI that trusts itself
becomes unstable. Ordinary confidence, statistical calibration, software
verification, and institutional delegation are not automatically
instances of a provability predicate satisfying the theorem's
conditions.

The theorem can motivate caution about systems asked to certify
unrestricted statements about their own future reasoning. It cannot
replace an architectural analysis.

\subsection{14.9 Bounded reflection}\label{bounded-reflection}

\textbf{{[}Design principle{]}} Reflective systems should scope their
self-trust to explicitly verified domains, preserve independent checks
where stakes are high, and avoid converting local reliability into
unrestricted authority over future modifications.

Examples include:

\begin{itemize}
\tightlist
\item
  proof-carrying changes within a restricted formal language;
\item
  empirical monitoring after deployment;
\item
  separation of proposer and verifier;
\item
  external authorization for changes to protected constraints;
\item
  versioned identity and rollback;
\item
  and refusal to let a system redefine the test by which it is
  certified.
\end{itemize}

This principle is supported by engineering and governance
considerations. Its connection to Löb is interpretive unless the
implemented system meets the theorem's formal conditions.

\subsection{14.10 Self-reference in open
systems}\label{self-reference-in-open-systems}

Human beings, institutions, and artificial agents are not closed formal
theories. They receive new data, change languages, revise axioms, and
alter their own material substrate. Formal limitative results still
matter, but the relevant question becomes how self-description interacts
with external correction.

Open systems face failure modes not captured by incompleteness alone:

\begin{itemize}
\tightlist
\item
  Goodharting the self-evaluation metric;
\item
  changing the environment to make the metric look favorable;
\item
  suppressing evidence that would trigger revision;
\item
  and transferring risk outside the represented boundary.
\end{itemize}

These are coherence-governance problems, not direct corollaries of
mathematical logic.

\subsection{14.11 Versioning as practical incompleteness
discipline}\label{versioning-as-practical-incompleteness-discipline}

CU's own response to fallibility is versioning. A numbered framework can
preserve a public record of what changed, why it changed, and which
downstream claims are affected.

Versioning does not prove humility. It makes revision inspectable. A
superseded claim should remain locatable; a new theorem should identify
its assumptions; a cross-domain correction should propagate through the
dependency graph.

This is the institutional form of the claim that no current formulation
is final.

\subsection{14.12 Formal limits and spiritual
mystery}\label{formal-limits-and-spiritual-mystery}

Formal incompleteness can resonate with spiritual traditions that
preserve mystery, but resonance is not entailment. Gödel does not prove
the sacred, ineffability, or transcendence.

The spiritual commitment to humility may be strengthened by recognizing
formal limits. Its warrant comes from the life and inquiry it supports,
not from treating a mathematical theorem as a concealed theology.

\subsection{14.13 Canonical conclusion}\label{canonical-conclusion-1}

\begin{quote}
Gödel, Tarski, and Löb establish precise limits for specified formal
systems and semantic constructions. CU adopts from them disciplines of
scope, stratification, and bounded reflection. It does not use them to
prove non-monotonicity, metaphysical openness, machine unsafety, or the
impossibility of fundamental physics.
\end{quote}

\begin{center}\rule{0.5\linewidth}{0.5pt}\end{center}

\section{15. Spiritual and Metaphysical
Interpretation}\label{spiritual-and-metaphysical-interpretation}

\subsection{15.1 The framework's origin}\label{the-frameworks-origin}

Coherence Universalism emerged from devotional poetry, contemplative
practice, moral concern, and an attempt to recover spiritually serious
orientation after inherited certainty had become untenable. The white
papers translated parts of that orientation into philosophy, science,
psychology, ethics, and system design.

Version 0.10 does not hide this lineage. It gives it a disciplined
register.

\subsection{15.2 Interpretation after
description}\label{interpretation-after-description}

A metaphysical interpretation asks what established structures may mean
as a picture of reality. It can unify experience, guide attention, and
generate research without being entailed by a theorem or experiment.

Interpretive commitments remain answerable to:

\begin{itemize}
\tightlist
\item
  established fact;
\item
  conceptual coherence;
\item
  comparative explanatory power;
\item
  lived and institutional consequence;
\item
  affected-party testimony;
\item
  and revision.
\end{itemize}

\subsection{15.3 CU-T9 --- Dual-Aspect Coherence
Monism}\label{cu-t9-dual-aspect-coherence-monism}

\textbf{{[}Interpretation{]}} Reality is one coherence-bearing
relational process rather than two independent substances. Its
externally describable organization and internally lived appearance,
wherever perspective occurs, are distinguishable but irreducible aspects
of the same reality.

The canonical minimum is \textbf{aspect coequality}: third-person
structure and first-person appearance are both real data for
metaphysics. Neither is eliminated in favor of the other.

The preferred thick trilogy interpretation is consciousness-first in a
restrained sense: interiority is not an ontologically alien addition to
an otherwise intrinsically mindless universe. This does not establish
panpsychism, cosmopsychism, universal phenomenal consciousness, or one
cosmic subject. The distribution and individuation of perspective remain
open.

\subsection{15.4 The coherence field as
interpretation}\label{the-coherence-field-as-interpretation}

\emph{Coherence field} can mean:

\begin{enumerate}
\def\labelenumi{\arabic{enumi}.}
\tightlist
\item
  a mathematical field or state variable in a declared model;
\item
  distributed relational organization across a system;
\item
  a metaphysical picture of reality's structured possibility space.
\end{enumerate}

Physics supports the first only under explicit constructions and the
second as a systems description. The third is interpretive. Dual-aspect
monism must not be presented as a result of the Fisher simplex or any
particular physical sandbox.

\subsection{15.5 CU-T10 --- Meaning
Realism}\label{cu-t10-meaning-realism}

\textbf{{[}Interpretation{]}} Meaning is neither wholly manufactured by
isolated minds nor simply imposed from outside. It arises through
truthful participation in real relationships, histories, needs,
possibilities, goods, and trajectories.

Meaning is:

\begin{itemize}
\tightlist
\item
  objective enough to constrain;
\item
  plural enough to admit interpretation;
\item
  historical enough to be revised;
\item
  and participatory enough to require enactment.
\end{itemize}

A life does not receive one immutable script. But neither is every
interpretation equally truthful, loving, viable, or just.

\begin{quote}
\textbf{Meaning is discovered through participation and actualized
through commitment.}
\end{quote}

\subsection{15.6 CU-T11 --- Minimum Teleology and the Grain of
Reality}\label{cu-t11-minimum-teleology-and-the-grain-of-reality}

\textbf{{[}Interpretation{]}} Reality is not equally hospitable to every
possible form or trajectory. Laws, constraints, histories, and
relational structures make some organizations more stable, viable,
generative, integrative, and capable of further development than others.

This is a \textbf{minimum teleology}:

\begin{itemize}
\tightlist
\item
  direction without foresight;
\item
  structured possibility without guaranteed progress;
\item
  attractor-like organization without one predetermined endpoint;
\item
  and real grain without a supernatural planner.
\end{itemize}

It does not establish inevitable complexity, moral goodness of
evolution, necessity of suffering, or a universal optimizer.

\subsubsection{The grain of reality and domain-relative
directionality}\label{the-grain-of-reality-and-domain-relative-directionality}

The grain of reality names recurrent, domain-relative directionality:
under declared constraints and horizons, some trajectories widen viable
integration, truthful relation, agency, love, adaptability, diversity,
and repair, while others narrow them.

A gradient is literal only inside a declared mathematical model with a
state space, geometry, and objective. Across the corpus, CU speaks
instead of a family of domain-relative transition structures; it does
not posit one hidden scalar or cross-domain gradient.

\subsection{15.7 The sacred}\label{the-sacred}

The sacred names that which is received as possessing significance
beyond its usefulness to the current evaluator and as placing a demand
on how one attends, acts, or lives.

Sacredness is not an empirical property, identical moral status,
universal consciousness, inviolability, or immunity from criticism. It
becomes corrupt when used to conceal hierarchy, romanticize harm, or
block proportionate intervention.

\subsection{15.8 CU-T12 --- Love as the Relational Interior of
Coherence}\label{cu-t12-love-as-the-relational-interior-of-coherence}

\textbf{{[}Interpretation{]}} Love is the deepest experiential and
ethical meaning of coherence across the self--other boundary: the
movement through which another being's reality becomes significant in
itself and enters one's field of concern without being reduced to
possession, use, or control.

In the strongest metaphysical reading:

\begin{quote}
\textbf{Love is the relational interior of reality's movement toward
wider coherence.}
\end{quote}

This is not the claim that chemical bonding, physical attraction,
biological symbiosis, and interpersonal love are one measurable force.
It is a structural and experiential interpretation whose conscious
expression is ethically decisive.

Love must remain joined to truth, consent, differentiation, justice,
boundaries, and repair. Possession, submission, compulsory intimacy,
self-erasure, forced forgiveness, and sentimental avoidance are failures
of love rather than its fulfillment.

\subsection{15.9 Suffering, meaning-revision, and
redemption}\label{suffering-meaning-revision-and-redemption}

Suffering is not automatically meaningful, necessary, or deserved. Some
suffering transforms; some only injures; some remains irreparable.

Human beings nevertheless possess a constrained practical freedom:
events cannot be changed, but the meaning through which events organize
the living present can sometimes be revised. This freedom depends on
conditions---safety, time, relationship, material support, language, and
access to alternative futures.

Redemption names the possibility that suffering can become solidarity,
wisdom, creation, justice, care, or repair. It never makes the original
evil good or necessary.

\[\adjustbox{max width=\linewidth}{$\displaystyle 
\text{meaning found after suffering}
\neq
\text{suffering justified by later meaning}.
$}\]

\subsection{15.10 CU-T13 --- The Immanent
Omega}\label{cu-t13-the-immanent-omega}

\textbf{{[}Interpretation{]}} Omega is not primarily a terminal event at
the end of cosmic, technological, or human history. It is the limit-form
of universal coherence locally realizable in the present.

When a finite conscious being:

\begin{enumerate}
\def\labelenumi{\arabic{enumi}.}
\tightlist
\item
  recognizes its constitutive participation in the whole;
\item
  allows another being's reality to matter intrinsically;
\item
  and acts through love rather than defensive self-maintenance,
\end{enumerate}

the universe realizes its integrative direction at that location.

A local Omega realization can be \textbf{complete in kind while partial
in extent}. The person remains finite, embodied, fallible, particular,
and part of an unfinished world.

The universal is not reached by containing everything in consciousness.
It is reached when nothing is excluded in principle from the field of
love.

\subsection{15.11 Omega horizon, realization, embodiment, and
ecology}\label{omega-horizon-realization-embodiment-and-ecology}

CU distinguishes:

\begin{itemize}
\tightlist
\item
  \textbf{Omega horizon:} the regulative limit of complete truthful,
  loving, non-dominating relation;
\item
  \textbf{Omega realization:} a local event in which universalizing love
  organizes a finite center;
\item
  \textbf{Omega embodiment:} stabilization in conduct, character,
  accountability, and repair;
\item
  \textbf{Omega ecology:} the material, relational, institutional, and
  ecological conditions that make such realization more possible.
\end{itemize}

Connectivity, complexity, intelligence, mystical intensity, or political
unity are not Omega by themselves. A planetary network can be anti-Omega
if organized through domination.

\subsection{15.12 Phenomenological
accessibility}\label{phenomenological-accessibility}

Omega is not posited as wholly foreign to human life. Its features
appear in graded form in:

\begin{itemize}
\tightlist
\item
  attachment and falling in love;
\item
  holding and caring for a child;
\item
  sympathetic joy;
\item
  friendship and service;
\item
  awe, beauty, music, and collective experience;
\item
  contemplative practice;
\item
  spontaneous mystical states;
\item
  and psychedelic experience.
\end{itemize}

These routes are not equivalent. They support the hypothesis that
reduced defensive self-prioritization and widened loving concern are
humanly accessible. They do not prove the metaphysics, guarantee
stability, or establish moral maturity.

The research rule is:

\begin{quote}
Psychedelics bear on accessibility. Contemplative practice bears on
trainability. Ordinary love bears on natural occurrence. Ethical conduct
bears on integration.
\end{quote}

\subsection{15.13 Omega glimpses are not Omega
embodiment}\label{omega-glimpses-are-not-omega-embodiment}

A peak experience can coexist with narcissism, exploitation,
grandiosity, epistemic recklessness, or failed ordinary responsibility.

Omega embodiment is tested by its fruits:

\begin{itemize}
\tightlist
\item
  truthfulness;
\item
  humility;
\item
  non-domination;
\item
  preserved boundaries;
\item
  joy in another's flourishing;
\item
  accountability;
\item
  service;
\item
  and repair.
\end{itemize}

\[\adjustbox{max width=\linewidth}{$\displaystyle 
\text{self-transcendent experience}
\not\Rightarrow
\text{Omega embodiment}.
$}\]

\subsection{15.14 Tradition, nonduality, and plural thick
interpretations}\label{tradition-nonduality-and-plural-thick-interpretations}

The Dao, Brahman, śūnyatā, Logos, God, the Good, liberation, and the
Kingdom of God arise in distinct traditions and should not be flattened
into technical synonyms.

CU offers a thin shared architecture while permitting thick
interpretations. Nonduality can be held as the insight that ordinary
self/world separation is not ultimate in the form it appears, while
preserving physical boundaries, personal difference, consent, and
unresolved metaphysical questions.

\subsection{15.15 Faith, ritual, and
practice}\label{faith-ritual-and-practice}

Faith is revisable trust enacted under incomplete control. Ritual and
prayer can organize attention, emotion, memory, relationship, and
community. Their empirical effects, interpretive meanings, and enacted
commitments must be distinguished.

Practice is participation: the effort to make truth, love, embodiment,
responsibility, and repair increasingly continuous in ordinary life.

\subsection{15.16 Heaven and Earth}\label{heaven-and-earth}

Within the project's symbolic language, \emph{Heaven} names form,
possibility, abstraction, and orienting meaning; \emph{Earth} names
embodiment, consequence, limitation, and lived relation.

Their equivalence does not make them interchangeable. It names a refusal
to sever structure from embodiment, interpretation from consequence, or
spiritual aspiration from the conditions of actual lives.

\subsection{15.17 Mystery preserved}\label{mystery-preserved}

Mystery is not a license for unsupported assertion. It is the
recognition that the territory exceeds every map and that responsible
commitment can coexist with unresolved description.

The appropriate response is neither certainty nor dismissal, but
disciplined openness.

\subsection{15.18 Interpretive
conclusion}\label{interpretive-conclusion}

CU's spiritual-metaphysical arc can be stated without pretending it is a
theorem:

\begin{quote}
Reality is one differentiated process with exterior and interior
aspects. Its possibility spaces possess real grain. Meaning is
discovered through truthful participation. Love is the relational
interior of coherence. Omega is locally realized whenever universal
concern becomes embodied in truthful, non-dominating relation. The world
remains unfinished, and that incompletion is the field of vocation.
\end{quote}

\section{16. The Commitment Register: Truth, Love, Meaning, Redemption,
and
Vocation}\label{the-commitment-register-truth-love-meaning-redemption-and-vocation}

\subsection{16.1 Register status}\label{register-status}

The CU-COM family records what Coherence Universalism officially
undertakes to affirm and practice. These principles are neither
empirical discoveries nor deductions from physical law. They are
constrained by fact, evaluated through personal and institutional
trajectories, and open to revision.

\subsection{16.2 The constitutive truth
discipline}\label{the-constitutive-truth-discipline}

\begin{itemize}
\tightlist
\item
  \textbf{CU-COM1 --- Commitment to Truth:} make reality more important
  than identity.
\item
  \textbf{CU-COM2 --- Discipline of the Two Registers:} do not let
  commitments determine what is accepted as fact, and do not let
  descriptive method erase meaning.
\item
  \textbf{CU-COM3 --- Revision in Both Registers:} revise descriptions
  through evidence and commitments through fact, consequence, testimony,
  integrity, and repair.
\end{itemize}

\subsection{16.3 Sacred regard, reciprocal dependence, and the grain of
vocation}\label{sacred-regard-reciprocal-dependence-and-the-grain-of-vocation}

\begin{itemize}
\tightlist
\item
  \textbf{CU-COM4 --- Sacred Regard:} receive beings and conditions as
  significant beyond instrumental use.
\item
  \textbf{CU-COM5 --- Reciprocal Dependence and the Coherence Mandala:}
  no local flourishing remains coherent by destroying the wider field
  that sustains it.
\item
  \textbf{CU-COM6 --- The Grain of Coherence and Coherence as Vocation:}
  discern and follow domain-relative directions by which truth, viable
  integration, agency, love, diversity, and repair can be widened
  without claiming one universal scalar or guaranteed progress.
\end{itemize}

\subsection{16.4 Love and objective-participatory
meaning}\label{love-and-objective-participatory-meaning}

\begin{itemize}
\tightlist
\item
  \textbf{CU-COM14 --- Love as the Ground and Aim of Conscious Life:}
  growth in the capacity to love is the primary practical purpose of
  conscious life.
\item
  \textbf{CU-COM15 --- Meaning as Discovery and Participation:} meaning
  is discovered through truthful relation to real goods, histories,
  responsibilities, and possibilities, then actualized through
  commitment.
\end{itemize}

Love does not mean possession, compulsory intimacy, self-erasure,
indiscriminate trust, forced forgiveness, or sentimental avoidance of
justice. Objective meaning does not imply fixed destiny or institutional
authority to assign purpose.

\subsection{16.5 Freedom, redemption, and the unfinished
past}\label{freedom-redemption-and-the-unfinished-past}

\begin{itemize}
\tightlist
\item
  \textbf{CU-COM16 --- Freedom of Meaning-Revision:} events remain
  fixed, but persons and communities can sometimes revise what those
  events mean and which futures they make possible.
\item
  \textbf{CU-COM17 --- Redemptive Possibility Without Justification of
  Evil:} suffering may become care, solidarity, creation, justice, or
  repair; later transformation never makes the original harm necessary
  or good.
\end{itemize}

No person owes growth, forgiveness, reconciliation, testimony, or
redemption. Freedom of meaning-revision depends on material,
psychological, relational, and political conditions.

\subsection{16.6 Heaven on Earth and the Immanent
Omega}\label{heaven-on-earth-and-the-immanent-omega}

\begin{itemize}
\tightlist
\item
  \textbf{CU-COM18 --- Heaven on Earth as Social Vocation:} work toward
  revisable social attractors in which truth, material sufficiency,
  plural agency, love, ecological reciprocity, non-domination,
  participation, and repair reinforce one another.
\item
  \textbf{CU-COM19 --- Enact Omega Here:} do not postpone universal
  relation and sacred responsibility to a singularity or end of history;
  enact Omega through truthful presence, non-dominating love, service,
  widening concern, and repair.
\end{itemize}

The Kingdom of God is not a theocracy or perfect state. It is an
Open-Myth name for Omega ecology: conditions that make truthful and
loving participation more possible for more beings.

\subsection{16.7 Domain commitments}\label{domain-commitments}

CU-COM7--CU-COM13 remain the official domain commitments for Physics,
Biology, Consciousness, Psychology, Social Dynamics, Ethics, and
Artificial Intelligence. The new framework commitments deepen rather
than replace them.

\subsection{16.8 The truth arc completed}\label{the-truth-arc-completed}

The corpus begins with truth about reality, passes through the forms
reality takes, and returns to truth about participation:

\[\adjustbox{max width=\linewidth}{$\displaystyle 
\text{truth about reality}
\rightarrow
\text{meaning within reality}
\rightarrow
\text{truth about ourselves in relation}.
$}\]

The discipline can be stated simply:

\begin{quote}
Epistemology teaches us not to lie about the universe. Ethics teaches us
not to lie about ourselves. Love asks us to make the truth livable in
relation.
\end{quote}

\subsection{16.9 Review and revision}\label{review-and-revision}

No CU-COM principle is beyond revision. A change must identify:

\begin{itemize}
\tightlist
\item
  the facts, arguments, consequences, testimony, or misuse patterns
  motivating it;
\item
  the affected communities;
\item
  the downstream papers;
\item
  and how descriptive evidence remains distinct from commitment
  authority.
\end{itemize}

\section{17. Research Program, Failure Conditions, and
Conclusion}\label{research-program-failure-conditions-and-conclusion}

\subsection{17.1 What the framework now
establishes}\label{what-the-framework-now-establishes}

The framework now has a more modest but stronger architecture.

\begin{itemize}
\tightlist
\item
  Coherence is a shared structural concept with domain-specific
  realizations.
\item
  Truth is distinct from coherence.
\item
  Description and commitment are distinct but asymmetrically connected.
\item
  Skeptical preferential consequence is separated from credulous
  admissibility.
\item
  Dynamic preferential models retain formal KLM results when
  trajectories are treated as worlds.
\item
  The two-stage policy architecture is retained as a decision procedure
  without a false System R guarantee.
\item
  The Epistemology Ladder is a capacity and governance map, not a
  hierarchy of truths.
\item
  The Coherence Ladder is a revisable dependency graph.
\item
  The Mandala is a reciprocal systems map and spiritual interpretation,
  not a proof.
\item
  Gödel, Tarski, and Löb are used within scope.
\end{itemize}

\subsection{17.2 Formal research program}\label{formal-research-program}

The immediate formal tasks are:

\begin{enumerate}
\def\labelenumi{\arabic{enumi}.}
\tightlist
\item
  develop preferential models with vector-valued and lexicographic
  coherence orderings;
\item
  characterize when partial orders yield useful consequence without
  forced scalarization;
\item
  distinguish commitment expansion, revision, contraction, and
  transformation;
\item
  identify fragments of the two-stage architecture with provable
  composition properties;
\item
  formalize protected floors and local--global externalization;
\item
  develop dynamic warrant over changing identities and constraints;
\item
  prove representation results only after the intended relation is
  fixed; and
\item
  create machine-checkable versions of the core propositions.
\end{enumerate}

\subsection{17.3 Empirical research
program}\label{empirical-research-program}

The framework requires evidence where it makes claims about real
systems. Priority programs include:

\begin{itemize}
\tightlist
\item
  testing operational measures of biological integration and repair;
\item
  distinguishing psychological integration from rigid self-consistency;
\item
  measuring institutional correction capacity and hidden cost export;
\item
  studying whether workspace-like self-modeling in artificial systems is
  causally global, capacity-limited, persistent, and singular;
\item
  testing candidate signatures of temporal thickness and valence;
\item
  comparing coherence-governed decision systems with established safety
  and calibration baselines;
\item
  and evaluating contemplative practices with first-person and
  third-person methods without reducing either to the other.
\end{itemize}

\subsection{17.4 Cross-domain validation}\label{cross-domain-validation}

A claimed invariant across domains should satisfy at least four tests:

\begin{enumerate}
\def\labelenumi{\arabic{enumi}.}
\tightlist
\item
  \textbf{independent operationalization:} each domain can define the
  construct without merely copying vocabulary;
\item
  \textbf{structural correspondence:} a specified mapping preserves a
  nontrivial relation;
\item
  \textbf{explanatory gain:} the mapping explains or predicts something
  not already obtained more simply;
\item
  \textbf{failure discrimination:} the framework identifies cases where
  the analogy breaks.
\end{enumerate}

Without these, cross-domain language remains heuristic.

\subsection{17.5 Falsifiability and narrowing
conditions}\label{falsifiability-and-narrowing-conditions}

CU should be narrowed or abandoned in any domain where:

\begin{itemize}
\tightlist
\item
  ``coherence'' cannot be operationalized beyond redescribing success;
\item
  proposed measures fail to distinguish integration from rigidity,
  exploitation, or mere stability;
\item
  bridge principles generate no discriminating predictions;
\item
  domain-specific theories explain the same phenomena with less
  assumption and greater accuracy;
\item
  the framework repeatedly absorbs counterexamples by changing
  definitions after the fact;
\item
  its ethical conclusions require smuggling values into allegedly
  physical premises;
\item
  or its spiritual interpretations systematically obstruct factual
  correction or human agency.
\end{itemize}

The universalist aspiration should be abandoned if no rigorous
translations can be established among domain realizations.

\subsection{17.6 What would count as
success}\label{what-would-count-as-success}

Success need not mean one final equation. CU would earn its place if it
produced:

\begin{itemize}
\tightlist
\item
  clearer dependency maps;
\item
  better separation of claim types;
\item
  useful operational measures;
\item
  reproducible formal results;
\item
  novel cross-domain hypotheses;
\item
  safer AI and institutional architectures;
\item
  improved practices of self-correction;
\item
  and a spiritually meaningful orientation that remains honest about
  what it knows.
\end{itemize}

\subsection{17.7 Relation to the Foundations
document}\label{relation-to-the-foundations-document}

This paper and \emph{Coherence Universalism: Foundations} form the
synchronized methodological and registry base for the v0.10 domain
sequence. Domain authority identifies the controlling CU formulation
without converting conditional empirical, clinical, political,
normative, or engineering content into a universal law. External review
remains part of revision and consequential application rather than a
condition for stating the complete proposal.

The corpus distinguishes:

\begin{itemize}
\tightlist
\item
  canonical framework principles;
\item
  domain principles;
\item
  bridge hypotheses;
\item
  companion interpretations;
\item
  retired claims;
\item
  and open research questions.
\end{itemize}

\subsection{17.8 The integrative project
resumed}\label{the-integrative-project-resumed}

Modern specialization was necessary. The task now is not to reverse it
but to build translations disciplined enough that specialized knowledge
can participate in a larger orientation without losing its standards.

Integration fails when it erases distinctions. It succeeds when the
distinctions become mutually informative.

\subsection{17.9 Conclusion}\label{conclusion}

Coherence Universalism began with a spiritual and philosophical
intuition: that reality, life, mind, morality, and meaning might be
understood as expressions of one recurring demand for relation under
constraint. Version 0.09 gave that intuition a broad architecture. It
also made the architecture appear more deductively complete than it was.

Version 0.10 begins again from a harder form of fidelity. It preserves
the vision by making every transition visible. It allows physics to be
physics, commitment to be commitment, and spiritual interpretation to
remain spiritually meaningful without borrowing certainty from
mathematics.

The result is not a closed system. It is a governed conversation among
forms of knowing.

That is sufficient for the work ahead.

\begin{center}\rule{0.5\linewidth}{0.5pt}\end{center}

\section{Appendix A. Formal Definitions and
Results}\label{appendix-a.-formal-definitions-and-results}

\subsection{A.1 Descriptive satisfaction and commitment
realization}\label{a.1-descriptive-satisfaction-and-commitment-realization}

Let \(L_D\) be a descriptive language with truth-satisfaction
\(\vDash_D\), and \(L_R\) a declared commitment-representation language
with realization relation \(\Vdash_R\).

For the full Boolean preferential fragment:

\[\adjustbox{max width=\linewidth}{$\displaystyle 
\llbracket\neg\varphi\rrbracket_R=W\setminus\llbracket\varphi\rrbracket_R,
$}\]

\[\adjustbox{max width=\linewidth}{$\displaystyle 
\llbracket\varphi\land\psi\rrbracket_R=
\llbracket\varphi\rrbracket_R\cap\llbracket\psi\rrbracket_R,
$}\]

\[\adjustbox{max width=\linewidth}{$\displaystyle 
\llbracket\varphi\lor\psi\rrbracket_R=
\llbracket\varphi\rrbracket_R\cup\llbracket\psi\rrbracket_R.
$}\]

The formalism assumes stable addressable contents. It does not derive
their compositional meaning.

\subsection{A.2 Fixed-base models}\label{a.2-fixed-base-models}

Let:

\[\adjustbox{max width=\linewidth}{$\displaystyle 
M=(W,V_F,\prec,\Vdash_R),
$}\]

with \(V_F\equiv V_{F,K,A,H}\) and \(\prec\) fixed while premises vary.

\[\adjustbox{max width=\linewidth}{$\displaystyle 
W_{F,\Gamma}=V_F\cap\operatorname{Real}_R(\Gamma),
\qquad
B(F,\Gamma)=\operatorname{Min}_{\prec}(W_{F,\Gamma}).
$}\]

\subsection{A.3 Skeptical preferential
consequence}\label{a.3-skeptical-preferential-consequence}

\[\adjustbox{max width=\linewidth}{$\displaystyle 
\Gamma\mid\!\sim_C\varphi
\iff
\operatorname{Min}_{\prec}(V_F\cap\operatorname{Real}_R(\Gamma))
\subseteq\llbracket\varphi\rrbracket_R.
$}\]

Under definability and smoothness, the relation satisfies System P.
Formula-level \textbf{Or} presupposes the declared Boolean realization
algebra.

\subsection{A.4 Ranked models}\label{a.4-ranked-models}

If the preference is ranked and every relevant nonempty definable set
attains a minimum rank, Rational Monotony is supported:

\[\adjustbox{max width=\linewidth}{$\displaystyle 
\alpha\mid\!\sim_C\beta
\text{ and }
\alpha\not\mid\!\sim_C\neg\gamma
\Rightarrow
\alpha\land\gamma\mid\!\sim_C\beta.
$}\]

\subsection{A.5 Practical verdicts}\label{a.5-practical-verdicts}

For a nonempty best set:

\[\adjustbox{max width=\linewidth}{$\displaystyle 
\text{strong warrant: } B\subseteq\llbracket\varphi\rrbracket_R,
$}\]

\[\adjustbox{max width=\linewidth}{$\displaystyle 
\text{admissibility: } B\cap\llbracket\varphi\rrbracket_R\neq\varnothing,
$}\]

\[\adjustbox{max width=\linewidth}{$\displaystyle 
\text{exclusion: } B\cap\llbracket\varphi\rrbracket_R=\varnothing.
$}\]

\textbf{Mere admissibility} is admissibility without strong warrant.

\textbf{Infeasibility} obtains when \(W_{F,\Gamma}=\varnothing\).

\textbf{Underdetermination} has two forms: formal non-attainment when
the premise-compatible set is nonempty but has no minimal elements, and
pre-verdict model-adequacy abstention when the represented facts,
options, affected scope, floors, horizon, observation model, or ordering
remain inadequate or contested. Ties within an attained best set are
handled by the best-set verdicts rather than classified as
underdetermination.

\subsection{A.6 Classical negation}\label{a.6-classical-negation}

Within the Boolean fragment and a nonempty best set:

\[\adjustbox{max width=\linewidth}{$\displaystyle 
\text{exclusion of }\varphi
\iff
\text{strong warrant for }\neg\varphi.
$}\]

This is an adopted semantic choice with acknowledged practical costs.

\subsection{A.7 Maximum preservation}\label{a.7-maximum-preservation}

Where maxima are attained, restricting the feasible set cannot increase
the maximum. Equality after adding \(\varphi\) is equivalent to support
by at least one old maximizer and therefore tests admissibility, not
skeptical consequence.

\subsection{A.8 Trajectory lifting}\label{a.8-trajectory-lifting}

Replacing worlds with trajectories preserves the abstract fixed-base
construction when:

\begin{itemize}
\tightlist
\item
  the trajectory-level base remains fixed;
\item
  the ordering remains fixed;
\item
  and adoption filters trajectories by intersection.
\end{itemize}

Changed floors, affected scope, horizon, viability criterion, or
ordering constitute model revision.

\subsection{A.9 Practical expansion
warrant}\label{a.9-practical-expansion-warrant}

\[\adjustbox{max width=\linewidth}{$\displaystyle 
\mathcal W_C(\varphi\mid F,\Gamma,K,H,A,M,U)
$}\]

is schematic notation for governance review. It includes facts,
protected floors, affected parties, distribution, option-set adequacy,
independent challenge, correctability, uncertainty, reversibility, exit,
and repair.

\subsection{A.10 Model capture}\label{a.10-model-capture}

Formal determinacy relative to a model does not establish
intersubjective objectivity. Public specification makes assumptions
visible; contestability requires normative standing and institutional
mechanisms not supplied by the preferential form.

\subsection{A.11 Two-stage
obstruction}\label{a.11-two-stage-obstruction}

A premise-dependent policy map does not inherit System P or System R
merely because each fixed-policy inner relation has those properties.
The outer selection map must preserve the relevant Boolean and
preferential structure or the composite relation must be treated
separately.

\subsection{A.12 Formal limitation}\label{a.12-formal-limitation}

The formal results do not determine:

\begin{itemize}
\tightlist
\item
  meaning constitution;
\item
  affected-party standing;
\item
  protected floors;
\item
  model adequacy;
\item
  option-set completeness;
\item
  normative ordering;
\item
  institutional legitimacy;
\item
  or practical permissibility.
\end{itemize}

\section{Appendix B. Countermodels and Obstruction
Results}\label{appendix-b.-countermodels-and-obstruction-results}

\subsection{B.1 Internal coherence without
truth}\label{b.1-internal-coherence-without-truth}

Let a theory contain mutually supporting claims \(p,q,r\), with

\[\adjustbox{max width=\linewidth}{$\displaystyle 
p\rightarrow q,
\quad
q\rightarrow r,
\quad
r\rightarrow p.
$}\]

The set may be internally consistent and highly integrated while all
three claims are false in the intended interpretation. Inferential
integration therefore does not entail descriptive truth.

\subsection{B.2 Stability without
adaptivity}\label{b.2-stability-without-adaptivity}

Let a dynamical system possess a deep attractor that resists
perturbation but leads outside the viability region when the environment
changes. The system is locally stable and globally nonviable. Stability
therefore does not entail coherence across horizon and environment.

\subsection{B.3 Low entropy without substantive
coherence}\label{b.3-low-entropy-without-substantive-coherence}

A distribution concentrated on one arbitrary state has low Shannon
entropy and high differentiation from uniformity. Nothing in those facts
alone establishes relational integration, physical function, or moral
value. Entropy is not a universal inverse coherence measure.

\subsection{B.4 Admissibility conjunction
failure}\label{b.4-admissibility-conjunction-failure}

Let the best set contain two tied worlds, \(B=\{w_1,w_2\}\). Let
\(w_1\vDash\varphi\land\neg\psi\) and \(w_2\vDash\psi\land\neg\varphi\).
Then \(\varphi\) preserves the optimum \(w_1\), and \(\psi\) preserves
the optimum \(w_2\), but \(\varphi\land\psi\) preserves no optimum. This
directly refutes the identification of maximum-preserving expansion with
System P consequence.

\subsection{B.5 Mixed-policy disjunction
failure}\label{b.5-mixed-policy-disjunction-failure}

\subsection{Version 0.10 Probability-threshold disjunction
failure}\label{version-0.10-probability-threshold-disjunction-failure}

Let \(A\cap B=\varnothing\) and \(\delta=0\). Policies \(\pi_A\) and
\(\pi_B\) concentrate on their named events. Policy \(\pi_M\) divides
probability between them. Then \(\pi_M\in\Pi_{A\lor B}\) while
\(\pi_M\notin\Pi_A\cup\Pi_B\). Assigning lower risk to \(\pi_M\) makes
the disjunctive selector novel and can defeat Or.

\subsection{B.6 Situated usefulness without
truth}\label{b.6-situated-usefulness-without-truth}

A false map can guide a traveler successfully over a short route because
the errors do not affect that path. Its success is evidence only for the
limited action, not for every represented feature. Truth is not ``what
works here and now.''

\subsection{B.7 Intersubjective stabilization without
truth}\label{b.7-intersubjective-stabilization-without-truth}

A community can coordinate around a shared false account when testimony
is dependent, dissent is punished, and evidence channels are controlled.
Agreement increases social stability while decreasing descriptive
adequacy.

\subsection{B.8 Formal derivability without empirical
application}\label{b.8-formal-derivability-without-empirical-application}

From axioms \(A\), theorem \(T\) may follow validly. If \(A\) does not
describe the target system, the proof remains valid and the application
fails. Formal status and physical status must be labeled separately.

\subsection{B.9 Spiritual warrant without descriptive
proof}\label{b.9-spiritual-warrant-without-descriptive-proof}

A person may be warranted in living through a commitment that every
conscious life deserves reverence. The practical and moral warrant does
not prove a cosmology in which all minds are literally parts of one
cosmic subject.

\subsection{B.10 Descriptive falsity defeating commitment
use}\label{b.10-descriptive-falsity-defeating-commitment-use}

A healing practice may carry symbolic meaning. If it is used to replace
effective treatment on the false claim that it cures a dangerous
disease, the descriptive component is subject to evidence and can render
that use unwarranted.

\subsection{B.11 Gödel overextension}\label{b.11-guxf6del-overextension}

Let a philosophical framework lack an effectively axiomatized formal
language representing arithmetic. Gödel's first theorem does not
directly apply to it as a formal theory. Calling the framework
``Gödel-incomplete'' may be metaphorical, but cannot function as a
theorem about the framework.

\subsection{B.12 Bridge omission}\label{b.12-bridge-omission}

Suppose a mathematical model contains a globally pure state whose local
reductions are mixed. Concluding that the physical universe is globally
pure requires an ontological bridge identifying the model state with
reality. The mathematical fact alone does not provide the bridge.

\subsection{B.13 Premise-dependent viability can defeat Cautious
Monotony}\label{b.13-premise-dependent-viability-can-defeat-cautious-monotony}

Let \(W=\{w_1,w_2,w_3\}\). Let \(w_1\vDash b\land g\),
\(w_2\vDash\neg b\land g\), and \(w_3\vDash b\land\neg g\). Order
\(w_1\prec w_2\) and \(w_3\prec w_2\), with \(w_1\) and \(w_3\)
incomparable. Define a premise-responsive viable set by

\[\adjustbox{max width=\linewidth}{$\displaystyle 
V_\varnothing=\{w_1,w_2\},
\qquad
V_{\{b\}}=\{w_3\}.
$}\]

Then the unique minimal world at the empty background is \(w_1\), so
\(\varnothing\mid\!\sim_C b\) and \(\varnothing\mid\!\sim_C g\). After
adding \(b\), the unique minimal world is \(w_3\), so
\(\{b\}\not\mid\!\sim_C g\). Cautious Monotony fails. The countermodel
shows why the System P result of A.3 requires a fixed base filtered by
\(\operatorname{Real}_R(\Gamma)\).

\begin{center}\rule{0.5\linewidth}{0.5pt}\end{center}

\section{Appendix C. Claim-Status and Translation
Protocol}\label{appendix-c.-claim-status-and-translation-protocol}

\subsection{C.1 Claim card}\label{c.1-claim-card}

Every load-bearing claim in a v0.10 paper should be representable by the
following card:

\begin{longtable}[]{@{}
  >{\raggedright\arraybackslash}p{(\columnwidth - 2\tabcolsep) * \real{0.5000}}
  >{\raggedright\arraybackslash}p{(\columnwidth - 2\tabcolsep) * \real{0.5000}}@{}}
\toprule\noalign{}
\begin{minipage}[b]{\linewidth}\raggedright
Field
\end{minipage} & \begin{minipage}[b]{\linewidth}\raggedright
Required entry
\end{minipage} \\
\midrule\noalign{}
\endhead
\bottomrule\noalign{}
\endlastfoot
Claim & One sentence, with variables and scope \\
Status & Definition, formal result, conditional result, imported result,
empirical claim, bridge principle, hypothesis, interpretation,
commitment, or open problem \\
Domain & Physical, biological, conscious, psychological, social,
ethical, artificial, or cross-domain \\
Assumptions & Explicit premises and background model \\
Support & Derivation, evidence, citation, argument, or practice \\
Bridge & Any domain or register transition \\
Failure condition & What would narrow or defeat the claim \\
Dependencies & Canonical principles or source claims used \\
Downstream effects & Principles or sections requiring revision if it
changes \\
\end{longtable}

\subsection{C.2 Translation card}\label{c.2-translation-card}

A cross-domain translation should state:

\begin{longtable}[]{@{}
  >{\raggedright\arraybackslash}p{(\columnwidth - 2\tabcolsep) * \real{0.5000}}
  >{\raggedright\arraybackslash}p{(\columnwidth - 2\tabcolsep) * \real{0.5000}}@{}}
\toprule\noalign{}
\begin{minipage}[b]{\linewidth}\raggedright
Element
\end{minipage} & \begin{minipage}[b]{\linewidth}\raggedright
Required question
\end{minipage} \\
\midrule\noalign{}
\endhead
\bottomrule\noalign{}
\endlastfoot
Source construct & What is being translated? \\
Target construct & What is the proposed analogue, realization, or
application? \\
Translation License & Is the relation analogy, shared schema,
property-preserving transfer, exact structural transfer, or empirically
validated instantiation? \\
Preserved relation & Which property, invariant, ordering, or operation
is claimed to survive? \\
New target constraints & What assumptions or conditions appear only in
the target domain? \\
Lost information & What distinctions or information does the mapping
discard? \\
Target-domain evidence & What supports the mapping in the target rather
than only in the source? \\
Failure or narrowing condition & What would defeat or weaken the
translation? \\
Claim-status change & Does the translated claim become a hypothesis,
bridge, interpretation, or another status? \\
Review competence & Which domain expertise and affected-party review are
required? \\
Competing map & What alternative translations or realizations should
remain visible? \\
\end{longtable}

\subsection{C.3 Register card}\label{c.3-register-card}

For an important utterance or principle:

\begin{enumerate}
\def\labelenumi{\arabic{enumi}.}
\tightlist
\item
  Is it being used descriptively, interpretively, normatively, or as a
  commitment?
\item
  What would count as evidence against any descriptive component?
\item
  What actions does the commitment component license?
\item
  Which facts constrain it?
\item
  Does it remain meaningful if the metaphysical interpretation is false?
\item
  Does it preserve exit, correction, and agency?
\end{enumerate}

\subsection{C.4 Mathematical application
card}\label{c.4-mathematical-application-card}

A mathematical construction used physically or socially should state:

\begin{itemize}
\tightlist
\item
  mathematical object and domain;
\item
  theorem or identity proved;
\item
  data-generating or physical interpretation;
\item
  parameter units;
\item
  boundary conditions;
\item
  calibration method;
\item
  alternatives;
\item
  numerical verification;
\item
  and whether the construction is a fit, embedding, recovery, or
  derivation.
\end{itemize}

\subsection{C.5 Revision protocol}\label{c.5-revision-protocol}

When a canonical claim changes:

\begin{enumerate}
\def\labelenumi{\arabic{enumi}.}
\tightlist
\item
  preserve the superseded wording in the change ledger;
\item
  identify the reason for revision;
\item
  state whether the change is correction, narrowing, reclassification,
  or retirement;
\item
  list every dependent principle and manuscript;
\item
  update the canonical registry;
\item
  rerun formal or numerical checks;
\item
  and record unresolved consequences.
\end{enumerate}

\begin{center}\rule{0.5\linewidth}{0.5pt}\end{center}

\section{Appendix D. Domain Dependency
Map}\label{appendix-d.-domain-dependency-map}

\subsection{D.1 Cross-domain
architecture}\label{d.1-cross-domain-architecture}

\begin{longtable}[]{@{}
  >{\raggedright\arraybackslash}p{(\columnwidth - 8\tabcolsep) * \real{0.2000}}
  >{\raggedright\arraybackslash}p{(\columnwidth - 8\tabcolsep) * \real{0.2000}}
  >{\raggedright\arraybackslash}p{(\columnwidth - 8\tabcolsep) * \real{0.2000}}
  >{\raggedright\arraybackslash}p{(\columnwidth - 8\tabcolsep) * \real{0.2000}}
  >{\raggedright\arraybackslash}p{(\columnwidth - 8\tabcolsep) * \real{0.2000}}@{}}
\toprule\noalign{}
\begin{minipage}[b]{\linewidth}\raggedright
Domain
\end{minipage} & \begin{minipage}[b]{\linewidth}\raggedright
Primary object
\end{minipage} & \begin{minipage}[b]{\linewidth}\raggedright
Characteristic constraints
\end{minipage} & \begin{minipage}[b]{\linewidth}\raggedright
Candidate coherence form
\end{minipage} & \begin{minipage}[b]{\linewidth}\raggedright
Critical bridge out
\end{minipage} \\
\midrule\noalign{}
\endhead
\bottomrule\noalign{}
\endlastfoot
Metaphysics/Epistemology & claims, models, commitments, dependency
relations & truth, evidence, inference, viability, register discipline &
integration and adequacy profile & from conceptual architecture to
domain hypotheses \\
Physics & states, histories, fields, accessibility structures &
dynamics, symmetry, conservation, boundary conditions & domain-specific
functionals and relational structure & from mathematical construction to
physical interpretation \\
Biology & self-maintaining organizations & metabolism, boundary
maintenance, repair, reproduction, selection & integrated viability
across scales & from persistence to intrinsic stake and cognition \\
Consciousness & perspectives and access organization & individuation,
global availability, temporal binding, valence & multidimensional
perspective profile & from functional organization to phenomenology \\
Psychology & persons, meaning maps, affective and cognitive trajectories
& embodiment, memory, identity, relationship, development & integration,
adequacy, flexibility, repair & from individual organization to shared
fields and normativity \\
Social Dynamics & multi-agent systems and institutions & power, trust,
information, infrastructure, ecology & coordination with correction and
cost inclusion & from stability to legitimacy \\
Ethics & reasons, agency, harms, responsibilities & recognition,
non-deception, protected agency, repair & justificatory integrity & from
structural conditions to normative authority \\
Artificial Intelligence & engineered models, agents, and sociotechnical
systems & objectives, data, memory, permissions, deployment, governance
& architecture-specific profiles & from behavior to perspective,
responsibility, and sovereignty \\
\end{longtable}

\subsection{D.2 Authority dependencies for
v0.10}\label{d.2-authority-dependencies-for-v0.10}

\begin{enumerate}
\def\labelenumi{\arabic{enumi}.}
\tightlist
\item
  \emph{Coherence Universalism: Physics}, Version 0.10 supplies current
  physical constraints.
\item
  M\&E Version 0.10 supplies claim status, register discipline,
  preferential logic, identity architecture, and cross-domain
  governance.
\item
  Foundations Version 0.10 supplies the synchronized authority and
  migration registry.
\item
  \emph{Coherence Universalism: Biology}, Version 0.10 through
  \emph{Coherence Universalism: Artificial Intelligence}, Version 0.10
  own their mature domain formulations and may propose explicit changes
  upstream.
\item
  Active companion papers may control their extended exposition,
  examples, and synthesis within an explicitly stated scope, but they
  remain subordinate to M\&E, \emph{Foundations}, and the owning domain
  paper.
\item
  Future changes propagate through versioned consolidation rather than
  silent cross-domain borrowing.
\end{enumerate}

\subsection{D.3 Companion-paper
status}\label{d.3-companion-paper-status}

\begin{itemize}
\tightlist
\item
  \emph{The Coherence Field}, Volume I: historical technical precursor
  superseded by \emph{Coherence Universalism: Physics}, Version 0.10.
\item
  \emph{The Coherence Field}, Volume II: superseded as a
  consciousness-exclusion argument; its useful dynamical questions are
  absorbed into \emph{Coherence Universalism: Consciousness}, Version
  0.10.
\item
  \emph{The Coherence Field}, Volume III: candidate interpretive
  companion requiring reconstruction.
\item
  \emph{The Architecture of Coherent Mind}: historical architecture
  paper; its consciousness and formal-guarantee interpretations are
  superseded.
\item
  \emph{The Architecture of Non-Egoic Artificial Mind}: active Version
  0.10 companion paper; a focused design synthesis subordinate to the
  Consciousness, Ethics, and Artificial Intelligence authorities.
\item
  \emph{Denying Knowledge to Make Room for Commitment}, Version 0.10:
  active philosophical companion and controlling source for the Version
  0.10 commitment-logic clarification.
\end{itemize}

\begin{center}\rule{0.5\linewidth}{0.5pt}\end{center}

\section{Appendix E. Glossary}\label{appendix-e.-glossary}

\textbf{Accessibility.} Availability of a distinction, correlation,
reason, or action possibility to a specified system, scale, operation,
and horizon.

\textbf{Admissibility.} The condition that at least one best viable
realization remains compatible with a candidate conclusion or
commitment.

\textbf{Bridge principle.} A stated connection between domains, claim
types, or formalisms that is not derived from the source level alone.

\textbf{Coherence.} Integrated relational organization under specified
constraints.

\textbf{Coherence capacity.} The range of perturbations and
continuations over which relevant organization can be preserved or
repaired.

\textbf{Coherence profile.} A vector or structured set of dimensions
used instead of a single scalar score.

\textbf{Commitment register.} A use of language or practice whose
principal role is to undertake, orient, consecrate, trust, refuse,
promise, organize attention, or structure action and relationship; it
remains constrained by fact, consequence, agency, revision, and repair.

\textbf{Description register.} A use of language that purports to
represent how things are and therefore answers to truth, evidence, valid
inference, calibration, and correction.

\textbf{Egoic organization.} Organization around preserving a bounded,
diachronically persistent self-model as a privileged object.

\textbf{Epistemic integrity.} Maintained capacity to let relevant
evidence and criticism alter claims, confidence, models, and action
without selective disabling of standards.

\textbf{Exclusion.} The condition that no best viable realization
supports a candidate commitment.

\textbf{Externalization.} Preservation or improvement within a selected
boundary by transferring costs to another scale, system, agent, or
horizon.

\textbf{Identity criterion.} A rule specifying which relations,
capacities, histories, or constraints must persist for reidentification
through change.

\textbf{Interpretation.} An explanatory or metaphysical reading
compatible with a body of results but not entailed by them.

\textbf{Mandala.} CU's reciprocal dependency map and associated
spiritual image; not a formal proof.

\textbf{Preferential consequence.} Skeptical non-monotonic consequence
defined by truth in all minimal premise-compatible realizations of a
fixed viability base and fixed preference ordering.

\textbf{Protected floor.} A constraint that may not be traded away by
aggregate optimization within the relevant decision procedure.

\textbf{Ranked model.} A preferential model whose ordering is
represented by a totally ordered rank function; System R additionally
requires minimal ranks to be attained on the relevant nonempty definable
sets.

\textbf{Rationalization.} Selection or production of reasons because
they protect a preferred conclusion while being presented as though they
independently produced it.

\textbf{Register mixing.} Use of warrant from one discourse register to
certify a claim in another without an explicit bridge.

\textbf{Self-maintenance.} Activity organized around preserving the
continued existence or functioning of a system.

\textbf{Self-modeling.} Representation or tracking by a system of
aspects of its own state or activity.

\textbf{Smoothness.} Preferential-model condition under which every
relevant nonminimal realization lies above a minimal one in the same
premise-compatible set.

\textbf{Strong warrant.} Support by every best viable realization under
the specified model, provided the best set is nonempty.

\textbf{System P.} KLM axiom system characterized by smooth preferential
semantics.

\textbf{System R.} Rational consequence system characterized by ranked
preferential semantics.

\textbf{Two-stage architecture.} A decision procedure that first selects
policies under feasibility and risk, then assesses policy-indexed
typical consequences.

\textbf{Viability.} Continued maintenance of domain-specified
organization above protected failure boundaries.

\begin{center}\rule{0.5\linewidth}{0.5pt}\end{center}

\textbf{Addressable commitment content.} Commitment content stable
enough to make a public, criticizable difference to inference, action,
expectation, relationship, or revision and, where formalized, to receive
declared realization conditions.

\textbf{Commitment realization (\(\Vdash_R\)).} The adopted relation by
which a world or trajectory realizes or supports representable
commitment content. It is distinguished from descriptive
truth-satisfaction and does not by itself solve the semantics of
ordinary practical or religious language.

\textbf{Descriptive satisfaction (\(\vDash_D\)).} Truth-satisfaction for
the descriptive language.

\textbf{Dual-aspect coherence monism.} The interpretation that exterior
relational organization and interior appearance are irreducible aspects
of one reality wherever perspective occurs.

\textbf{Grain of reality.} Recurrent domain-relative directionality
produced by real laws, constraints, histories, and relations; it does
not imply one universal scalar or predetermined end.

\textbf{Heaven on Earth.} An Open-Myth name for a revisable family of
social attractors in which truth, sufficiency, plural agency, love,
ecological reciprocity, non-domination, participation, and repair
reinforce one another.

\textbf{Immanent Omega.} The interpretation that Omega is locally
realizable in present truthful and non-dominating love, complete in kind
but partial in extent, rather than requiring a final singularity or end
of history.

\textbf{Meaning realism.} The view that meaning is discovered through
truthful participation in real goods, relationships, histories,
responsibilities, and possibilities while remaining plural and
revisable.

\textbf{Omega ecology.} Material, psychological, relational,
institutional, and ecological conditions that make Omega glimpses,
stabilization, and embodiment more possible.

\textbf{Omega embodiment.} The stabilization of self-transcendent or
loving insight in truthfulness, humility, preserved boundaries,
non-domination, service, accountability, and repair.

\textbf{Omega glimpse.} A temporary state of reduced defensive
self-prioritization and widened love, unity, awe, connectedness, or
sympathetic joy.

\textbf{Omega realization.} A local event in which constitutive
participation, conscious recognition, and loving enactment converge.

\textbf{Omega stabilization.} Increasing capacity to return to and
inhabit an Omega orientation without extraordinary induction.

\textbf{Redemptive possibility.} The possibility that suffering can be
transformed into meaning, solidarity, creation, justice, care, or repair
without becoming justified retroactively.

\textbf{Grain of reality.} Interpretive name for recurrent
domain-relative directions that widen viable integration, truth, agency,
love, adaptability, diversity, and repair under declared constraints and
horizons. It does not posit one cross-domain scalar or gradient.

\section{Appendix F. Cross-Domain Governance
Matrix}\label{appendix-f.-cross-domain-governance-matrix}

The following family is framework-canonical as governance of CU
application. It does not function as a descriptive law of nature and
does not eliminate the review gates named here.

\begin{longtable}[]{@{}
  >{\raggedright\arraybackslash}p{(\columnwidth - 6\tabcolsep) * \real{0.2500}}
  >{\raggedright\arraybackslash}p{(\columnwidth - 6\tabcolsep) * \real{0.2500}}
  >{\raggedright\arraybackslash}p{(\columnwidth - 6\tabcolsep) * \real{0.2500}}
  >{\raggedright\arraybackslash}p{(\columnwidth - 6\tabcolsep) * \real{0.2500}}@{}}
\toprule\noalign{}
\begin{minipage}[b]{\linewidth}\raggedright
Identifier
\end{minipage} & \begin{minipage}[b]{\linewidth}\raggedright
Principle
\end{minipage} & \begin{minipage}[b]{\linewidth}\raggedright
Principal domain sources
\end{minipage} & \begin{minipage}[b]{\linewidth}\raggedright
Review gate
\end{minipage} \\
\midrule\noalign{}
\endhead
\bottomrule\noalign{}
\endlastfoot
CU-G1 & Affected Scope, Distribution, and Protected Conditions & CU-FP4;
CU-D9; CU-E-9; CU-S5; CU-EC-2; CU-AI-12 & cross-domain methods, ethics,
social science, and affected-community review \\
CU-G2 & Option-Set and Framing Integrity & CU-E-1; CU-E-15; CU-AI-23 &
ethics, decision theory, institutional design, and affected-community
review \\
CU-G3 & Practical Correctability and Independent Challenge & CU-FP7;
CU-D8; CU-Ψ7; CU-S3; CU-Inst-3; CU-E-6; CU-AI-24 & institutional design,
law, ethics, governance, and affected-community review \\
CU-G4 & Responsibility--Authority--Causal-Reach Alignment & CU-ICF-2;
CU-ICF-3; CU-Inst-6; CU-E-14; CU-AI-6; CU-AI-15 & ethics, law,
institutional design, AI governance, and affected-community review \\
CU-G5 & Participation Proportional to Governance Reach & CU-E-15; CU-S5;
CU-Inst-0; CU-AI-14; CU-AI-23 & heightened affected-community
co-governance and public-review design \\
CU-G6 & Patiency, Agency, Responsibility, and Sovereignty Are Distinct &
CU-V7; CU-E-2; CU-E-17; CU-AI-8; CU-AI-22 & consciousness science,
ethics, law, AI welfare, and public governance review \\
CU-G7 & Reversibility, Containment, and Retained Option Value & CU-E-16;
CU-AI-16 & ethics, safety engineering, law, domain expertise, and
affected-community review \\
\end{longtable}

The source principles remain active within their domains except where
they are explicitly recorded as merged aliases. CU-V7 and CU-E-17 are
retained only as historical/domain aliases of CU-G6. Domain-specific
principles may impose stronger requirements than the framework family.

\begin{center}\rule{0.5\linewidth}{0.5pt}\end{center}

\section{Appendix G. Technical Supplement and
Reproducibility}\label{appendix-g.-technical-supplement-and-reproducibility}

\textbf{Released evidence package.} The current supplement is
\href{../supplements/01_Metaphysics_Epistemology_Technical_Supplement_v0.10.md}{01\_Metaphysics\_Epistemology\_Technical\_Supplement\_v0.10.md}.
Its companion evidence directory is
\href{../evidence/01_Metaphysics_Epistemology/README.md}{01\_Metaphysics\_Epistemology}.
The supplement and evidence directory expose the derivations,
deterministic constructions, scripts, machine-readable results, source
ledger, and audit materials referenced in this manuscript. They do not
elevate a toy model, simulation, or reconstructed check into
target-domain empirical validation.

\section{References}\label{references}

Ay, N., Jost, J., Lê, H. V., \& Schwachhöfer, L. (2015). Information
geometry and sufficient statistics. \emph{Probability Theory and Related
Fields, 162}, 327--364. \url{https://doi.org/10.1007/s00440-014-0574-8}

Bauer, M., Bruveris, M., \& Michor, P. W. (2016). Uniqueness of the
Fisher--Rao metric on the space of smooth densities. \emph{Bulletin of
the London Mathematical Society, 48}(3), 499--506.
\url{https://doi.org/10.1112/blms/bdw020}

Campbell, L. L. (1986). An extended Čencov characterization of the
information metric. \emph{Proceedings of the American Mathematical
Society, 98}(1), 135--141.

Čencov, N. N. (1982). \emph{Statistical Decision Rules and Optimal
Inference}. American Mathematical Society.

Petz, D., \& Sudár, C. (1996). Geometries of quantum states.
\emph{Journal of Mathematical Physics, 37}(6), 2662--2673.
\url{https://doi.org/10.1063/1.531535}

Alchourrón, C. E., Gärdenfors, P., \& Makinson, D. (1985). On the logic
of theory change: Partial meet contraction and revision functions.
\emph{The Journal of Symbolic Logic, 50}(2), 510--530.

Carnap, R. (1950). Empiricism, semantics, and ontology. \emph{Revue
Internationale de Philosophie, 4}, 20--40.

Gödel, K. (1931). Über formal unentscheidbare Sätze der \emph{Principia
Mathematica} und verwandter Systeme I. \emph{Monatshefte für Mathematik
und Physik, 38}, 173--198.

Kant, I. (1781/1787). \emph{Critique of Pure Reason}.

Kraus, S., Lehmann, D., \& Magidor, M. (1990). Nonmonotonic reasoning,
preferential models and cumulative logics. \emph{Artificial
Intelligence, 44}(1--2), 167--207.

Lehmann, D., \& Magidor, M. (1992). What does a conditional knowledge
base entail? \emph{Artificial Intelligence, 55}(1), 1--60.

Löb, M. H. (1955). Solution of a problem of Leon Henkin. \emph{The
Journal of Symbolic Logic, 20}(2), 115--118.

Quine, W. V. O. (1951). Two dogmas of empiricism. \emph{The
Philosophical Review, 60}(1), 20--43.

Tarski, A. (1936). Der Wahrheitsbegriff in den formalisierten Sprachen.
\emph{Studia Philosophica, 1}, 261--405.

Wittgenstein, L. (1953). \emph{Philosophical Investigations}. Blackwell.

Rader, G. K. D. (2026). \emph{Coherence Universalism:
Foundations---Principles, Domain Gateways, Commitments, and Constraints
on Responsible Application} (Version 0.10). Heaven≡Earth Foundation.

Rader, G. K. D. (2026). \emph{Coherence Universalism:
Physics---Coherence, Accessibility, and the Ordered Emergence of
Physical Description} (Version 0.10). Heaven≡Earth Foundation.

Rader, G. K. D. (2026). \emph{Denying Knowledge to Make Room for
Commitment---Truth, Commitment, and Nonmonotonic Warrant} (Version
0.10). Heaven≡Earth Foundation.

Rader, G. K. D. (2026). \emph{The Architecture of Non-Egoic Artificial
Mind---Perspective, Self-Maintenance, and Alignment Without a Sovereign
Self} (Version 0.10). Heaven≡Earth Foundation.

Rader, G. K. D. (2026). \emph{Love Letters to Humanity}. Heaven≡Earth
Press.

Rader, G. K. D. (2026). \emph{You Are the Universe Becoming Conscious of
Itself: The Coherence Mandala}. Heaven≡Earth Press.

\end{document}
